%============================================================
%  Bd. 35 - Frankls Vermutung
%============================================================

\documentclass[main.tex]{subfiles}

\ifSubfilesClassLoaded{
  \BandSetupStandalone{B35}
}{
}

\title{Bd. 35 - Frankls Vermutung}
\author{Martin Kunze}
\date{}
\setcounter{file}{35}

\hypersetup{
  pdftitle={Bd. 35 - Frankls Vermutung},
  pdfauthor={Martin Kunze}
}

\begin{document}

\title{Bd. 35 - Frankls Vermutung}
\author{Martin Kunze}
\date{}
\hypersetup{pageanchor=false}
\maketitle
\hypersetup{pageanchor=true}
\setcounter{tocdepth}{3}
\tableofcontents
\setcounter{file}{35}
\renewcommand{\FormulaBandID}{35}
\newcommand{\ProofRefStack}[1]{%
  \ensuremath{\begin{gathered}#1\end{gathered}}}

\chapter{Einleitung}

Dieser Band ist ein eigenständiges Sammelbecken für die Mengen-, Quotienten-
und Halbverbandsformen von Frankls Vermutung. Die Vermutung selbst bleibt
klar von den vollständig bewiesenen Teilresultaten und Reduktionen getrennt.
Nur fachübergreifende Werkzeuge bleiben in ihren Grundlagenbänden: die
allgemeine Theorie von Mengenfamilien, Funktionen, Äquivalenzrelationen,
Ordnungen, Endlichkeit und Halbverbänden. Alle vermutungsspezifischen
Anwendungen dieser Werkzeuge werden hier zusammengeführt.

\chapter{Frankls Vermutung für Mengensysteme}

\section{Grundwerkzeuge für Frankl-Familien}

Die elementbezogenen Teilfamilien und ihre injektive Adjunktionsabbildung
wurden bereits in Band~06 eingeführt. Hier werden ihre
Frankl-spezifischen Mengen- und Endlichkeitseigenschaften entwickelt.

\subsection{Elementbezogene Teilfamilien}

Für eine Familie \(M\) und ein Element \(a\) verwenden wir die in Band~06
definierten Teilfamilien \(\FamContaining{M}{a}\) und
\(\FamAvoiding{M}{a}\).

\FormulaThmDelta[Enthaltende Teilfamilie ist Teilfamilie]{%
  \FamContaining{M}{a}\subseteq M
}{%
  \DeltaRow{Mengen}{M\dsep X\dsep a}
}
\begin{tabproof}
  \proofstep{1}{X\in \FamContaining{M}{a}}{\rA}
  \proofstep{1}{X\in M\land a\in X}{%
    \FormulaRefAuto{\FamContaining{M}{a}\coloneqq \{X\in M\mid a\in X\}}{1}}
  \proofstep{1}{X\in M}{\rAEa{2}}
  \proofstep{}{\FamContaining{M}{a}\subseteq M}{%
    \FormulaRefAuto{A \subseteq B \coloneq \forall x\,(x\in A \rightarrow x\in B)}{\rUI{\rRI{1,3}}}}
\end{tabproof}

\FormulaThmDelta[Vermeidende Teilfamilie ist Teilfamilie]{%
  \FamAvoiding{M}{a}\subseteq M
}{%
  \DeltaRow{Mengen}{M\dsep X\dsep a}
}
\begin{tabproof}
  \proofstep{1}{X\in \FamAvoiding{M}{a}}{\rA}
  \proofstep{1}{X\in M\land a\notin X}{%
    \FormulaRefAuto{\FamAvoiding{M}{a}\coloneqq \{X\in M\mid a\notin X\}}{1}}
  \proofstep{1}{X\in M}{\rAEa{2}}
  \proofstep{}{\FamAvoiding{M}{a}\subseteq M}{%
    \FormulaRefAuto{A \subseteq B \coloneq \forall x\,(x\in A \rightarrow x\in B)}{\rUI{\rRI{1,3}}}}
\end{tabproof}

\FormulaThmDeltaK[Enthaltende Teilfamilie als relatives Komplement der vermeidenden Teilfamilie]{%
  M\setminus\FamAvoiding{M}{a}=\FamContaining{M}{a}
}{FamContainingAsAvoidingRelativeComplement}{%
  \DeltaRow{Mengen}{M\dsep X\dsep a}
  \DeltaRow{Teilmengenfamilien}{\FamContaining{M}{a}\dsep \FamAvoiding{M}{a}}
}
\begin{tabproofsplitwide}
  \proofpartwideindR[Teilmengenrichtung
  \(M\setminus\FamAvoiding{M}{a}\subseteq\FamContaining{M}{a}\)]{%
    \vdash M\setminus\FamAvoiding{M}{a}\subseteq\FamContaining{M}{a}
  }{%
    \vdash M\setminus\FamAvoiding{M}{a}\subseteq\FamContaining{M}{a}
  }[FamContainingAsAvoidingRelativeComplementSubsetForward]
    \proofstepwidestar[1]{X\in M\setminus\FamAvoiding{M}{a}}{\rA}
    \proofstepwidestar[1]{X\in M\land X\notin\FamAvoiding{M}{a}}{%
      \FormulaRefAuto{x \in A \setminus B \eqvdash x \in A \land x \notin B}{1}}
    \proofstepwidestar[1]{X\in M}{\rAEa{2}}
    \proofstepwidestar[1]{X\notin\FamAvoiding{M}{a}}{\rAEb{2}}
    \proofstepwidestar[5]{a\notin X}{\rA}
    \proofstepwidestar[1,5]{X\in\FamAvoiding{M}{a}}{%
      \FormulaRefAuto{Vermeidende Teilfamilie}{\rAI{3,5}}}
    \proofstepwidestar[1,5]{\bot}{\rBI{6,4}}
    \proofstepwidestar[1]{\neg(a\notin X)}{\rCI{5,7}}
    \proofstepwidestar[1]{a\in X}{\FormulaRefAuto{rDN}{8}}
    \proofstepwidestar[1]{X\in M\land a\in X}{\rAI{3,9}}
    \proofstepwidestar[1]{X\in\FamContaining{M}{a}}{%
      \FormulaRefAuto{Enthaltende Teilfamilie}{10}}
    \proofstepwidestar{%
      M\setminus\FamAvoiding{M}{a}\subseteq\FamContaining{M}{a}}{%
      \FormulaRefAuto{A \subseteq B \coloneq
      \forall x\,(x\in A \rightarrow x\in B)}{\rUI{\rRI{1,11}}}}
  \closeproofpartwideindR

  \proofpartwideindR[Teilmengenrichtung
  \(\FamContaining{M}{a}\subseteq M\setminus\FamAvoiding{M}{a}\)]{%
    \vdash \FamContaining{M}{a}\subseteq M\setminus\FamAvoiding{M}{a}
  }{%
    \vdash \FamContaining{M}{a}\subseteq M\setminus\FamAvoiding{M}{a}
  }[FamContainingAsAvoidingRelativeComplementSubsetBackward]
    \proofstepwidestar[1]{X\in\FamContaining{M}{a}}{\rA}
    \proofstepwidestar[1]{X\in M\land a\in X}{%
      \FormulaRefAuto{Enthaltende Teilfamilie}{1}}
    \proofstepwidestar[1]{X\in M}{\rAEa{2}}
    \proofstepwidestar[1]{a\in X}{\rAEb{2}}
    \proofstepwidestar[5]{X\in\FamAvoiding{M}{a}}{\rA}
    \proofstepwidestar[5]{X\in M\land a\notin X}{%
      \FormulaRefAuto{Vermeidende Teilfamilie}{5}}
    \proofstepwidestar[5]{a\notin X}{\rAEb{6}}
    \proofstepwidestar[1,5]{\bot}{\rBI{4,7}}
    \proofstepwidestar[1]{X\notin\FamAvoiding{M}{a}}{\rCI{5,8}}
    \proofstepwidestar[1]{X\in M\land X\notin\FamAvoiding{M}{a}}{\rAI{3,9}}
    \proofstepwidestar[1]{X\in M\setminus\FamAvoiding{M}{a}}{%
      \FormulaRefAuto{x \in A \setminus B \eqvdash x \in A \land x \notin B}{10}}
    \proofstepwidestar{%
      \FamContaining{M}{a}\subseteq M\setminus\FamAvoiding{M}{a}}{%
      \FormulaRefAuto{A \subseteq B \coloneq
      \forall x\,(x\in A \rightarrow x\in B)}{\rUI{\rRI{1,11}}}}
  \closeproofpartwideindR

  \proofpartwide{Schluss}
    \proofstepwidestar{%
      M\setminus\FamAvoiding{M}{a}\subseteq\FamContaining{M}{a}}{%
      \FormulaRefAuto{FamContainingAsAvoidingRelativeComplementSubsetForward}}
    \proofstepwidestar{%
      \FamContaining{M}{a}\subseteq M\setminus\FamAvoiding{M}{a}}{%
      \FormulaRefAuto{FamContainingAsAvoidingRelativeComplementSubsetBackward}}
    \proofstepwidestar{M\setminus\FamAvoiding{M}{a}=\FamContaining{M}{a}}{%
      \FormulaRefAuto{A \subseteq B, B \subseteq A \vdash A = B}{1,2}}
  \closeproofpartwide
\end{tabproofsplitwide}

\FormulaThmDelta[Außenstehende Mengen sind nicht enthaltend]{%
  A\notin M\vdash A\notin \FamContaining{M}{a}
}{%
  \DeltaRow{Mengen}{M\dsep A\dsep a}
  \DeltaRow{Teilmengenfamilien}{\FamContaining{M}{a}}
}
\begin{tabproof}
  \proofstep{1}{A\notin M}{\rA}
  \proofstep{2}{A\in \FamContaining{M}{a}}{\rA}
  \proofstep{2}{A\in M\land a\in A}{%
    \FormulaRefAuto{\FamContaining{M}{a}\coloneqq
    \{X\in M\mid a\in X\}}{2}}
  \proofstep{2}{A\in M}{\rAEa{3}}
  \proofstep{1,2}{\bot}{\rBI{4,1}}
  \proofstep{1}{A\notin \FamContaining{M}{a}}{\rCI{2,5}}
\end{tabproof}

\FormulaThmDelta[Außenstehende Mengen sind nicht vermeidend]{%
  A\notin M\vdash A\notin \FamAvoiding{M}{a}
}{%
  \DeltaRow{Mengen}{M\dsep A\dsep a}
  \DeltaRow{Teilmengenfamilien}{\FamAvoiding{M}{a}}
}
\begin{tabproof}
  \proofstep{1}{A\notin M}{\rA}
  \proofstep{2}{A\in \FamAvoiding{M}{a}}{\rA}
  \proofstep{2}{A\in M\land a\notin A}{%
    \FormulaRefAuto{\FamAvoiding{M}{a}\coloneqq
    \{X\in M\mid a\notin X\}}{2}}
  \proofstep{2}{A\in M}{\rAEa{3}}
  \proofstep{1,2}{\bot}{\rBI{4,1}}
  \proofstep{1}{A\notin \FamAvoiding{M}{a}}{\rCI{2,5}}
\end{tabproof}

\FormulaThmDelta[Enthaltende Teilfamilie ist vereinigungsabgeschlossen]{%
  \UnionClosedFam(M)\vdash \UnionClosedFam(\FamContaining{M}{a})
}{%
  \DeltaRow{Mengen}{M\dsep X\dsep Y\dsep a}
}
\begin{tabproof}
  \proofstep{1}{\UnionClosedFam(M)}{\rA}
  \proofstep{2}{X\in \FamContaining{M}{a}}{\rA}
  \proofstep{3}{Y\in \FamContaining{M}{a}}{\rA}
  \proofstep{2}{X\in M\land a\in X}{%
    \FormulaRefAuto{\FamContaining{M}{a}\coloneqq \{X\in M\mid a\in X\}}{2}}
  \proofstep{2}{X\in M}{\rAEa{4}}
  \proofstep{2}{a\in X}{\rAEb{4}}
  \proofstep{3}{Y\in M\land a\in Y}{%
    \FormulaRefAuto{\FamContaining{M}{a}\coloneqq \{X\in M\mid a\in X\}}{3}}
  \proofstep{3}{Y\in M}{\rAEa{7}}
  \proofstep{1,2,3}{X\cup Y\in M}{%
    \FormulaRefAuto{X \in M \dsep Y \in M \vdash X \cup Y \in M}{5,8}}
  \proofstep{2}{a\in X\cup Y}{%
    \FormulaRefAuto{z \in A \vdash z \in A \cup B}{6}}
  \proofstep{1,2,3}{X\cup Y\in M\land a\in X\cup Y}{\rAI{9,10}}
  \proofstep{1,2,3}{X\cup Y\in \FamContaining{M}{a}}{%
    \FormulaRefAuto{\FamContaining{M}{a}\coloneqq \{X\in M\mid a\in X\}}{11}}
  \proofstep{1}{\UnionClosedFam(\FamContaining{M}{a})}{%
    \FormulaRefAuto{\UnionClosedFam(M)}{\rUI{\rRI{2,\rRI{3,12}}}}}
\end{tabproof}

\FormulaThmDeltaK[Zwei vermiedene Mengenbestandteile bleiben in der Vereinigung vermieden]{%
  a\notin X\dsep a\notin Y\vdash a\notin X\cup Y%
}{UnionNotMemberFromBoth}{%
  \DeltaRow{Mengen}{X\dsep Y\dsep a}%
}
\begin{tabproof}
  \proofstep{1}{a\notin X}{\rA}
  \proofstep{2}{a\notin Y}{\rA}
  \proofstep{3}{a\in X\cup Y}{\rA}
  \proofstep{3}{a\in X\lor a\in Y}{%
    \FormulaRefAuto{z\in A\cup B \eqvdash z\in A\lor z\in B}{3}}
  \proofstep{5}{a\in X}{\rA}
  \proofstep{1,5}{\bot}{\rBI{5,1}}
  \proofstep{7}{a\in Y}{\rA}
  \proofstep{2,7}{\bot}{\rBI{7,2}}
  \proofstep{1,2,3}{\bot}{\rOE{4,5,6,7,8}}
  \proofstep{1,2}{a\notin X\cup Y}{\rCI{3,9}}
\end{tabproof}

\FormulaThmDelta[Vermeidende Teilfamilie ist vereinigungsabgeschlossen]{%
  \UnionClosedFam(M)\vdash \UnionClosedFam(\FamAvoiding{M}{a})
}{%
  \DeltaRow{Mengen}{M\dsep X\dsep Y\dsep a}
}
\begin{tabproof}
  \proofstep{1}{\UnionClosedFam(M)}{\rA}
  \proofstep{2}{X\in \FamAvoiding{M}{a}}{\rA}
  \proofstep{3}{Y\in \FamAvoiding{M}{a}}{\rA}
  \proofstep{2}{X\in M\land a\notin X}{%
    \FormulaRefAuto{\FamAvoiding{M}{a}\coloneqq \{X\in M\mid a\notin X\}}{2}}
  \proofstep{2}{X\in M}{\rAEa{4}}
  \proofstep{2}{a\notin X}{\rAEb{4}}
  \proofstep{3}{Y\in M\land a\notin Y}{%
    \FormulaRefAuto{\FamAvoiding{M}{a}\coloneqq \{X\in M\mid a\notin X\}}{3}}
  \proofstep{3}{Y\in M}{\rAEa{7}}
  \proofstep{3}{a\notin Y}{\rAEb{7}}
  \proofstep{1,2,3}{X\cup Y\in M}{%
    \FormulaRefAuto{X \in M \dsep Y \in M \vdash X \cup Y \in M}{5,8}}
  \proofstep{2,3}{a\notin X\cup Y}{%
    \FormulaRefAuto{UnionNotMemberFromBoth}[thm]{6,9}}
  \proofstep{1,2,3}{X\cup Y\in M\land a\notin X\cup Y}{\rAI{10,11}}
  \proofstep{1,2,3}{X\cup Y\in \FamAvoiding{M}{a}}{%
    \FormulaRefAuto{\FamAvoiding{M}{a}\coloneqq \{X\in M\mid a\notin X\}}{12}}
  \proofstep{1}{\UnionClosedFam(\FamAvoiding{M}{a})}{%
    \FormulaRefAuto{\UnionClosedFam(M)}{\rUI{\rRI{2,\rRI{3,13}}}}}
\end{tabproof}

\subsection{Endlichkeit elementbezogener Teilfamilien}

\FormulaThmDeltaK[Enthaltende Teilfamilie einer endlichen Familie ist endlich]{%
  \Finite{M}\vdash \Finite{\FamContaining{M}{a}}
}{FiniteContainingSubfamily}{%
  \DeltaRow{Mengen}{M\dsep a}
}
\begin{tabproof}
  \proofstep{1}{\Finite{M}}{\rA}
  \proofstep{}{\FamContaining{M}{a}\subseteq M}{%
    \FormulaRefAuto{\FamContaining{M}{a}\subseteq M}}
  \proofstep{1}{\Finite{\FamContaining{M}{a}}}{%
    \FormulaRefAuto{FiniteSubset}[thm]{2,1}}
\end{tabproof}

\FormulaThmDeltaK[Vermeidende Teilfamilie einer endlichen Familie ist endlich]{%
  \Finite{M}\vdash \Finite{\FamAvoiding{M}{a}}
}{FiniteAvoidingSubfamily}{%
  \DeltaRow{Mengen}{M\dsep a}
}
\begin{tabproof}
  \proofstep{1}{\Finite{M}}{\rA}
  \proofstep{}{\FamAvoiding{M}{a}\subseteq M}{%
    \FormulaRefAuto{\FamAvoiding{M}{a}\subseteq M}}
  \proofstep{1}{\Finite{\FamAvoiding{M}{a}}}{%
    \FormulaRefAuto{FiniteSubset}[thm]{2,1}}
\end{tabproof}

\section{Frankl-Vermutung}

Die Frankl-Vermutung, benannt nach Péter Frankl, betrifft endliche
vereinigungsabgeschlossene Mengenfamilien. In der folgenden Formulierung ist
\(M\) eine Menge von Mengen.
Für ein Element \(a\) verwenden wir die zuvor eingeführte Zerlegung in die
Teilmengenfamilien \(\FamContaining{M}{a}\) und \(\FamAvoiding{M}{a}\).

\FormulaDefDeltaK[Halbhäufiges Element einer endlichen Familie]{%
  \HalfOcc{a}{M}\coloneqq
  a\in\bigcup M\land
  \exists F\colon \FamAvoiding{M}{a}\inj \FamContaining{M}{a}
}{Halbhäufiges Element}{%
  \DeltaRow{Mengen}{M\dsep X\dsep a}
  \DeltaRow{Funktionen}{F}
  \DeltaRow{Teilmengenfamilien}{\FamContaining{M}{a}\dsep \FamAvoiding{M}{a}}
  \DeltaRow{Neue Symbole}{\HalfOcc}
}

Für endliche \(M\) bedeutet \(\HalfOcc{a}{M}\), dass höchstens so
viele Mengen aus \(M\) das Element \(a\) vermeiden wie enthalten. Da
\(\FamContaining{M}{a}\) und \(\FamAvoiding{M}{a}\) zusammen \(M\)
zerlegen, ist dies die
injektive Form der Aussage, dass \(a\) in mindestens der Hälfte der Mengen
aus \(M\) enthalten ist.

\FormulaDefDeltaK[Frankl-Eigenschaft einer Familie]{%
  \Frankl{M}\coloneqq \exists a\,\HalfOcc{a}{M}
}{Frankl-Eigenschaft}{%
  \DeltaRow{Mengen}{M\dsep a}
  \DeltaRow{Neue Symbole}{\Frankl}
}

\FormulaDefDeltaK[Voraussetzung der Frankl-Vermutung]{%
  \FranklFam{M}
}{Frankl-Familie}{%
  \DeltaRow{Mengen}{M\dsep n}
  \DeltaRow{\textbf{Axiome}}{}
  \DeltaRow{Nichtleerheit}
           {M\neq\varnothing}
           [\parbox[t]{3cm}{\FormulaRefAuto{A\neq\varnothing \eqvdash \exists x\,(x \in A)}}]
  \DeltaRow{\makecell[l]{Nichtleerer\\ Träger}}
           {\bigcup M\neq\varnothing}
           [\parbox[t]{3cm}{%
             \FormulaRefAuto{\bigcup A \coloneq \iota C\Bigl(\forall x\;\bigl(x \in C \;\leftrightarrow\;\exists B\in A\,(x \in B)\bigr)\Bigr)}\\
             \FormulaRefAuto{A\neq\varnothing \eqvdash \exists x\,(x \in A)}}]
  \DeltaRow{Endlichkeit}
           {\Finite{M}}
           [\parbox[t]{3cm}{%
             \FormulaRefAuto{FiniteEqCardEquiv}[thm]}]
  \DeltaRow{\makecell[l]{Vereinigungs-\\ abschluss}}
           {\UnionClosedFam(M)}
           [\parbox[t]{3cm}{\FormulaRefAuto{Vereinigungsabgeschlossene Familie}}]
  \DeltaRow{\textbf{Neue Symbole}}{}
  \DeltaPrem{Frankl-Familien}{\FranklFam{M}}
}

\FormulaThmDeltaK[Frankl-Familien sind nichtleer]{%
  \FranklFam{M}\vdash M\neq\varnothing
}{FranklFamNonempty}{%
  \DeltaRow{Mengen}{M}
}
\begin{tabproof}
  \proofstep{1}{\FranklFam{M}}{\rA}
  \proofstep{1}{M\neq\varnothing}{\FormulaRefAuto{Frankl-Familie}{1}}
\end{tabproof}

\FormulaThmDeltaK[Frankl-Familien haben nichtleeren Träger]{%
  \FranklFam{M}\vdash \bigcup M\neq\varnothing
}{FranklFamCarrierNonempty}{%
  \DeltaRow{Mengen}{M}
}
\begin{tabproof}
  \proofstep{1}{\FranklFam{M}}{\rA}
  \proofstep{1}{\bigcup M\neq\varnothing}{\FormulaRefAuto{Frankl-Familie}{1}}
\end{tabproof}

\FormulaThmDeltaK[Frankl-Familien sind endlich]{%
  \FranklFam{M}\vdash \Finite{M}
}{FranklFamFinite}{%
  \DeltaRow{Mengen}{M}
}
\begin{tabproof}
  \proofstep{1}{\FranklFam{M}}{\rA}
  \proofstep{1}{\Finite{M}}{\FormulaRefAuto{Frankl-Familie}{1}}
\end{tabproof}

\FormulaThmDeltaK[Frankl-Familien sind vereinigungsabgeschlossen]{%
  \FranklFam{M}\vdash \UnionClosedFam(M)
}{FranklFamUnionClosed}{%
  \DeltaRow{Mengen}{M}
}
\begin{tabproof}
  \proofstep{1}{\FranklFam{M}}{\rA}
  \proofstep{1}{\UnionClosedFam(M)}{\FormulaRefAuto{Frankl-Familie}{1}}
\end{tabproof}

Wegen der Bedingung \(a\in\bigcup M\) in der Definition von
\(\HalfOcc{a}{M}\) muss die Trägervereinigung nichtleer sein. Die
Frankl-Vermutung lautet in dieser Form daher:
\[
  \FranklFam{M}\vdash \Frankl{M}.
\]
Ohne die Zusatzannahme \(\bigcup M\neq\varnothing\) wäre bereits
\(M=\{\varnothing\}\) ein Gegenbeispiel zur Formulierung, denn dann gibt es
kein \(a\in\bigcup M\). Die Aussage wird hier nicht als bewiesener Satz
verwendet. Wir zeigen zunächst einige einfache Spezialfälle.

\subsection{Fall einer enthaltenen Einermenge}

\FormulaThmDelta[Frankl-Vermutung bei enthaltener Einermenge]{%
  \Finite{M}\dsep \UnionClosedFam(M)\dsep \{a\}\in M\vdash \Frankl{M}
}{%
  \DeltaRow{Mengen}{M\dsep a}
  \DeltaRow{Funktionen}{\FranklAdj{M}{a}}
}
\begin{tabproof}
  \proofstep{1}{\Finite{M}}{\rA}
  \proofstep{2}{\UnionClosedFam(M)}{\rA}
  \proofstep{3}{\{a\}\in M}{\rA}
  \proofstep{2,3}{\FranklAdj{M}{a}\colon \FamAvoiding{M}{a}\inj \FamContaining{M}{a}}{%
    \FormulaRefAuto{\UnionClosedFam(M)\dsep \{a\}\in M
    \vdash \FranklAdj{M}{a}\colon \FamAvoiding{M}{a}\inj \FamContaining{M}{a}}{2,3}}
  \proofstep{2,3}{\exists F\colon \FamAvoiding{M}{a}\inj \FamContaining{M}{a}}{\rEI{4}}
  \proofstep{}{a\in\{a\}}{\FormulaRefAuto{a\in\{a\}}}
  \proofstep{3}{a\in\bigcup M}{%
    \FormulaRefAuto{B\in A\dsep x\in B\vdash x \in \bigcup A}{3,6}}
  \proofstep{2,3}{\HalfOcc{a}{M}}{%
    \FormulaRefAuto{\HalfOcc{a}{M}\coloneqq
    a\in\bigcup M\land
    \exists F\colon \FamAvoiding{M}{a}\inj \FamContaining{M}{a}}{\rAI{7,5}}}
  \proofstep{2,3}{\exists a\,\HalfOcc{a}{M}}{\rEI{8}}
  \proofstep{1,2,3}{\Frankl{M}}{%
    \FormulaRefAuto{\Frankl{M}\coloneqq \exists a\,\HalfOcc{a}{M}}{9}}
\end{tabproof}

\subsection{Fall eines nichtleeren Durchschnitts}

Mit der Notation für Durchschnitte von Mengenfamilien kann der zweite
Spezialfall direkt als \(\bigcap M\neq\varnothing\) formuliert werden.

\FormulaThmDelta[Leere vermeidende Teilfamilie liefert Frankl]{%
  a\in\bigcup M\dsep \FamAvoiding{M}{a}=\varnothing
  \vdash \Frankl{M}
}{%
  \DeltaRow{Mengen}{M\dsep a}
  \DeltaRow{Funktionen}{F}
  \DeltaRow{Teilmengenfamilien}{%
    \FamAvoiding{M}{a}\dsep \FamContaining{M}{a}}
}
\begin{tabproof}
  \proofstep{1}{a\in\bigcup M}{\rA}
  \proofstep{2}{\FamAvoiding{M}{a}=\varnothing}{\rA}
  \proofstep{}{\varnothing\colon\varnothing\inj \FamContaining{M}{a}}{%
    \FormulaRefAuto{\varnothing\colon \varnothing\inj M}}
  \proofstep{2}{\varnothing\colon
    \FamAvoiding{M}{a}\inj \FamContaining{M}{a}}{%
    \rIE{\FormulaRefAuto{a=b\vdash b=a}{2},3}}
  \proofstep{2}{%
    \exists F\colon \FamAvoiding{M}{a}\inj
    \FamContaining{M}{a}}{\rEI{4}}
  \proofstep{1,2}{\HalfOcc{a}{M}}{%
    \FormulaRefAuto{\HalfOcc{a}{M}\coloneqq
    a\in\bigcup M\land
    \exists F\colon \FamAvoiding{M}{a}\inj
    \FamContaining{M}{a}}{\rAI{1,5}}}
  \proofstep{1,2}{\exists a\,\HalfOcc{a}{M}}{\rEI{6}}
  \proofstep{1,2}{\Frankl{M}}{%
    \FormulaRefAuto{\Frankl{M}\coloneqq
    \exists a\,\HalfOcc{a}{M}}{7}}
\end{tabproof}

\FormulaThmDelta[Frankl-Vermutung bei nichtleerem Durchschnitt]{%
  \FranklFam{M}\dsep
  \bigcap M\neq\varnothing\vdash \Frankl{M}
}{%
  \DeltaRow{Mengen}{M\dsep X\dsep a}
  \DeltaRow{Funktionen}{F}
}
\begin{tabproof}
  \proofstep{1}{\FranklFam{M}}{\rA}
  \proofstep{2}{\bigcap M\neq\varnothing}{\rA}
  \proofstep{1}{M\neq\varnothing}{%
    \FormulaRefAuto{Frankl-Familie}{1}}
  \proofstep{2}{\exists a\in\bigcap M}{%
    \FormulaRefAuto{B\neq\varnothing\vdash \exists x\in B}{2}}

  \proofstep{5}{a\in\bigcap M}{\rA}
  \proofstep{1,5}{a\in\bigcup M}{%
    \FormulaRefAuto{M\neq\varnothing \dsep a\in\bigcap M
    \vdash a\in\bigcup M}{3,5}}

  \proofstep{7}{X\in \FamAvoiding{M}{a}}{\rA}
  \proofstep{7}{X\in M\land a\notin X}{%
    \FormulaRefAuto{\FamAvoiding{M}{a}\coloneqq \{X\in M\mid a\notin X\}}{7}}
  \proofstep{7}{X\in M}{\rAEa{8}}
  \proofstep{7}{a\notin X}{\rAEb{8}}
  \proofstep{1,5,7}{a\in X}{%
    \FormulaRefAuto{M\neq\varnothing \dsep a\in\bigcap M \dsep
    X\in M \vdash a\in X}{3,5,9}}
  \proofstep{1,5,7}{\bot}{\rBI{11,10}}
  \proofstep{1,5}{X\notin \FamAvoiding{M}{a}}{\rCI{7,12}}
  \proofstep{1,5}{\forall X\,X\notin \FamAvoiding{M}{a}}{\rUI{13}}
  \proofstep{1,5}{\FamAvoiding{M}{a}=\varnothing}{%
    \FormulaRefAuto{\forall x\, x\notin A\vdash A=\varnothing}{14}}

  \proofstep{1,5}{\Frankl{M}}{%
    \FormulaRefAuto{a\in\bigcup M\dsep
    \FamAvoiding{M}{a}=\varnothing
    \vdash \Frankl{M}}{6,15}}

  \proofstep{1,2}{\Frankl{M}}{\rEE{4,5,16}}
\end{tabproof}


\subsection{Fall einer enthaltenen Zweiermenge}

Wir betrachten nun den Fall, dass die Familie ein Glied \(\{a,b\}\)
enthält. Lokal zerlegen wir \(M\) nach dem Auftreten von \(a\) und \(b\):
\[
\begin{array}{c|cc}
 & b\text{ wird vermieden} & b\text{ wird enthalten}\\
\midrule
a\text{ vermieden} &
M_{00}=\FamAvoiding{M}{a}\cap\FamAvoiding{M}{b} &
M_{01}=\FamAvoiding{M}{a}\cap\FamContaining{M}{b}\\
a\text{ enthalten} &
M_{10}=\FamContaining{M}{a}\cap\FamAvoiding{M}{b} &
M_{11}=\FamContaining{M}{a}\cap\FamContaining{M}{b}
\end{array}
\]
Die Abbildung \(X\mapsto X\cup\{a,b\}\) injiziert \(M_{00}\) in
\(M_{11}\). Die gemischten Klassen \(M_{01}\) und \(M_{10}\) sind als
Teilmengen von \(M\) endlich. Gibt es eine Injektion
\(M_{01}\to M_{10}\), so werden die beiden Zweige zu einer Injektion
\(\FamAvoiding{M}{a}\to\FamContaining{M}{a}\) zusammengeführt; damit ist
\(a\) halbhäufig. Gibt es keine solche Injektion, liefert der Vergleich
endlicher Mengen eine Injektion \(M_{10}\to M_{01}\). Dann werden die
Zweige analog zu einer Injektion
\(\FamAvoiding{M}{b}\to\FamContaining{M}{b}\) zusammengeführt; damit ist
\(b\) halbhäufig.

\paragraph{Lokale Klassenrechnungen}

\FormulaThmDeltaK[Zerlegung der \(a\)-vermeidenden Klassen im Zweierfall]{%
  \begin{aligned}[t]
  \FamAvoiding{M}{a}
  &=
  \bigl(\FamAvoiding{M}{a}\cap\FamAvoiding{M}{b}\bigr)
  \cup
  \bigl(\FamAvoiding{M}{a}\cap\FamContaining{M}{b}\bigr)
  \end{aligned}
}{FranklPairAvoidADecomp}{%
  \DeltaRow{Mengen}{M\dsep X\dsep a\dsep b}
}
\begin{tabproof}
  \proofstep{}{\FamAvoiding{M}{a}\subseteq M}{%
    \FormulaRefAuto{\FamAvoiding{M}{a}\subseteq M}}
  \proofstep{}{
    \FamAvoiding{M}{a}
    =
    \bigl(\FamAvoiding{M}{a}\cap\FamAvoiding{M}{b}\bigr)
    \cup
    \bigl(\FamAvoiding{M}{a}\cap(M\setminus\FamAvoiding{M}{b})\bigr)}{%
    \FormulaRefAuto{SubsetRelativeComplementDecomp}{1}}
  \proofstep{}{M\setminus\FamAvoiding{M}{b}=\FamContaining{M}{b}}{%
    \FormulaRefAuto{FamContainingAsAvoidingRelativeComplement}}
  \proofstep{}{
    \FamAvoiding{M}{a}
    =
    \bigl(\FamAvoiding{M}{a}\cap\FamAvoiding{M}{b}\bigr)
    \cup
    \bigl(\FamAvoiding{M}{a}\cap\FamContaining{M}{b}\bigr)}{%
    \rIE{3,2}}
\end{tabproof}

\FormulaThmDeltaK[Disjunktheit der \(a\)-vermeidenden Zweige]{%
  \bigl(\FamAvoiding{M}{a}\cap\FamAvoiding{M}{b}\bigr)
  \cap
  \bigl(\FamAvoiding{M}{a}\cap\FamContaining{M}{b}\bigr)
  =\varnothing
}{FranklPairAvoidADisjoint}{%
  \DeltaRow{Mengen}{M\dsep X\dsep a\dsep b}
}
\begin{tabproof}
  \proofstep{}{
    \bigl(\FamAvoiding{M}{a}\cap\FamAvoiding{M}{b}\bigr)
    \cap
    \bigl(\FamAvoiding{M}{a}\cap(M\setminus\FamAvoiding{M}{b})\bigr)
    =\varnothing}{%
    \FormulaRefAuto{RelativeComplementDecompDisjoint}}
  \proofstep{}{M\setminus\FamAvoiding{M}{b}=\FamContaining{M}{b}}{%
    \FormulaRefAuto{FamContainingAsAvoidingRelativeComplement}}
  \proofstep{}{
    \bigl(\FamAvoiding{M}{a}\cap\FamAvoiding{M}{b}\bigr)
    \cap
    \bigl(\FamAvoiding{M}{a}\cap\FamContaining{M}{b}\bigr)
    =\varnothing}{%
    \rIE{2,1}}
\end{tabproof}

\FormulaThmDeltaK[Zerlegung der \(a\)-enthaltenden Zielklassen]{%
  \begin{aligned}[t]
  \FamContaining{M}{a}
  &=
  \bigl(\FamContaining{M}{a}\cap\FamContaining{M}{b}\bigr)
  \cup
  \bigl(\FamContaining{M}{a}\cap\FamAvoiding{M}{b}\bigr)
  \end{aligned}
}{FranklPairContainADecomp}{%
  \DeltaRow{Mengen}{M\dsep X\dsep a\dsep b}
}
\begin{tabproofwide}
  \proofstepwidestar{\FamContaining{M}{a}\subseteq M}{%
    \FormulaRefAuto{\FamContaining{M}{a}\subseteq M}}
  \proofstepwidestar{\FamAvoiding{M}{b}\subseteq M}{%
    \FormulaRefAuto{\FamAvoiding{M}{a}\subseteq M}}
  \proofstepwidestar{%
    M\setminus(M\setminus\FamAvoiding{M}{b})
    =
    \FamAvoiding{M}{b}}{%
    \FormulaRefAuto{RelCompInvolutive}{2}}
  \proofstepwidestar{M\setminus\FamAvoiding{M}{b}=\FamContaining{M}{b}}{%
    \FormulaRefAuto{FamContainingAsAvoidingRelativeComplement}}
  \proofstepwidestar{M\setminus\FamContaining{M}{b}=\FamAvoiding{M}{b}}{%
    \rIE{4,3}}
  \proofstepwide{\FamContaining{M}{a}}{=}{%
    \begin{aligned}[t]
    &\bigl(\FamContaining{M}{a}\cap\FamContaining{M}{b}\bigr)\\
    &\cup
    \bigl(\FamContaining{M}{a}\cap(M\setminus\FamContaining{M}{b})\bigr)
    \end{aligned}
  }{%
    \FormulaRefAuto{SubsetRelativeComplementDecomp}{1}}
  \proofstepwide{\FamContaining{M}{a}}{=}{%
    \begin{aligned}[t]
    &\bigl(\FamContaining{M}{a}\cap\FamContaining{M}{b}\bigr)\\
    &\cup
    \bigl(\FamContaining{M}{a}\cap\FamAvoiding{M}{b}\bigr)
    \end{aligned}
  }{\rIE{5,6}}
\end{tabproofwide}

\FormulaThmDeltaK[Disjunktheit der \(a\)-enthaltenden Zielzweige]{%
  \bigl(\FamContaining{M}{a}\cap\FamContaining{M}{b}\bigr)
  \cap
  \bigl(\FamContaining{M}{a}\cap\FamAvoiding{M}{b}\bigr)
  =\varnothing
}{FranklPairContainADisjoint}{%
  \DeltaRow{Mengen}{M\dsep X\dsep a\dsep b}
}
\begin{tabproofwide}
  \proofstepwidestar{\FamAvoiding{M}{b}\subseteq M}{%
    \FormulaRefAuto{\FamAvoiding{M}{a}\subseteq M}}
  \proofstepwidestar{%
    M\setminus(M\setminus\FamAvoiding{M}{b})
    =
    \FamAvoiding{M}{b}}{%
    \FormulaRefAuto{RelCompInvolutive}{1}}
  \proofstepwidestar{M\setminus\FamAvoiding{M}{b}=\FamContaining{M}{b}}{%
    \FormulaRefAuto{FamContainingAsAvoidingRelativeComplement}}
  \proofstepwidestar{M\setminus\FamContaining{M}{b}=\FamAvoiding{M}{b}}{%
    \rIE{3,2}}
  \proofstepwidestar{%
    \begin{aligned}[t]
    \Bigl[
    \bigl(\FamContaining{M}{a}\cap\FamContaining{M}{b}\bigr)
    \cap
    \bigl(\FamContaining{M}{a}\cap(M\setminus\FamContaining{M}{b})\bigr)
    \Bigr]
    &=\varnothing
    \end{aligned}
  }{%
    \FormulaRefAuto{RelativeComplementDecompDisjoint}}
  \proofstepwidestar{%
    \begin{aligned}[t]
    \Bigl[
    \bigl(\FamContaining{M}{a}\cap\FamContaining{M}{b}\bigr)
    \cap
    \bigl(\FamContaining{M}{a}\cap\FamAvoiding{M}{b}\bigr)
    \Bigr]
    &=\varnothing
    \end{aligned}
  }{\rIE{4,5}}
\end{tabproofwide}

\FormulaThmDeltaK[Endlichkeit der Klasse \(M_{01}\)]{%
  \Finite{M}\vdash
  \Finite{\FamAvoiding{M}{a}\cap\FamContaining{M}{b}}
}{FranklPairMixed01Finite}{%
  \DeltaRow{Mengen}{M\dsep a\dsep b}
}
\begin{tabproof}
  \proofstep{1}{\Finite{M}}{\rA}
  \proofstep{}{\FamAvoiding{M}{a}\subseteq M}{%
    \FormulaRefAuto{\FamAvoiding{M}{a}\subseteq M}}
  \proofstep{}{\FamAvoiding{M}{a}\cap\FamContaining{M}{b}\subseteq M}{%
    \FormulaRefAuto{A\subseteq M \vdash A\cap B\subseteq M}{2}}
  \proofstep{1}{\Finite{\FamAvoiding{M}{a}\cap\FamContaining{M}{b}}}{%
    \FormulaRefAuto{FiniteSubset}[thm]{3,1}}
\end{tabproof}

\FormulaThmDeltaK[Gegeninjektion im negativen Mischklassenfall]{%
  \begin{aligned}[t]
  &\Finite{M}\dsep
  \neg\exists F\colon
  \bigl(\FamAvoiding{M}{a}\cap\FamContaining{M}{b}\bigr)
  \inj
  \bigl(\FamContaining{M}{a}\cap\FamAvoiding{M}{b}\bigr)
  \vdash{}\\
  &\exists G\colon
  \bigl(\FamContaining{M}{a}\cap\FamAvoiding{M}{b}\bigr)
  \inj
  \bigl(\FamAvoiding{M}{a}\cap\FamContaining{M}{b}\bigr)
  \end{aligned}
}{FranklPairMixedReverseInjection}{%
  \DeltaRow{Mengen}{M\dsep a\dsep b}
  \DeltaRow{Funktionen}{F\dsep G}
}
\begin{tabproof}
  \proofstep{1}{\Finite{M}}{\rA}
  \proofstep{2}{\neg\exists F\colon
    \bigl(\FamAvoiding{M}{a}\cap\FamContaining{M}{b}\bigr)
    \inj
    \bigl(\FamContaining{M}{a}\cap\FamAvoiding{M}{b}\bigr)}{\rA}
  \proofstep{1}{\Finite{\FamAvoiding{M}{a}\cap\FamContaining{M}{b}}}{%
    \FormulaRefAuto{FranklPairMixed01Finite}{1}}
  \proofstep{1}{\Finite{\FamAvoiding{M}{b}\cap\FamContaining{M}{a}}}{%
    \FormulaRefAuto{FranklPairMixed01Finite}{1}}
  \proofstep{}{
    \FamAvoiding{M}{b}\cap\FamContaining{M}{a}
    =
    \FamContaining{M}{a}\cap\FamAvoiding{M}{b}}{%
    \FormulaRefAuto{A \cap B = B \cap A}}
  \proofstep{1}{\Finite{\FamContaining{M}{a}\cap\FamAvoiding{M}{b}}}{%
    \rIE{5,4}}
  \proofstep{1,2}{\exists G\colon
    \bigl(\FamContaining{M}{a}\cap\FamAvoiding{M}{b}\bigr)
    \inj
    \bigl(\FamAvoiding{M}{a}\cap\FamContaining{M}{b}\bigr)}{%
    \FormulaRefAuto{\Finite{M}\dsep \Finite{N}\dsep
    \neg\exists F\colon M\inj N
    \vdash \exists G\colon N\inj M}{3,6,2}}
\end{tabproof}

\paragraph{Die Zweieradjunktion}

\FormulaDefDeltaK[Zweieradjunktionsabbildung]{%
  \UnionClosedFam(M)\dsep \{a,b\}\in M\vdash
  \begin{aligned}[t]
    &U_{M,a,b}\colon
      \bigl(\FamAvoiding{M}{a}\cap\FamAvoiding{M}{b}\bigr)
      \\[-2pt]
    &\hspace{4em}\to
      \bigl(\FamContaining{M}{a}\cap\FamContaining{M}{b}\bigr),\\[-2pt]
    &\forall X\in\FamAvoiding{M}{a}\cap\FamAvoiding{M}{b}\,
      \bigl(U_{M,a,b}(X)\coloneqq X\cup\{a,b\}\bigr)
  \end{aligned}
}{FranklPairAdjDef}{%
  \DeltaRow{Mengen}{M\dsep X\dsep a\dsep b}
  \DeltaRow{Funktionen}{U_{M,a,b}}
  \DeltaRow{\textbf{Neues Symbol}}{}
  \DeltaPrem{Zweieradjunktionsabbildung}{U_{M,a,b}}
}
\begin{tabproof}
  \proofstep{1}{\UnionClosedFam(M)}{\rA}
  \proofstep{2}{\{a,b\}\in M}{\rA}
  \proofstep{3}{X\in\FamAvoiding{M}{a}\cap\FamAvoiding{M}{b}}{\rA}
  \proofstep{3}{X\in\FamAvoiding{M}{a}}{%
    \FormulaRefAuto{x\in A\cap B\vdash x\in A}{3}}
  \proofstep{3}{X\in M}{%
    \rAEa{\FormulaRefAuto{\FamAvoiding{M}{a}\coloneqq
      \{X\in M\mid a\notin X\}}{4}}}
  \proofstep{1,2,3}{X\cup\{a,b\}\in M}{%
    \FormulaRefAuto{Vereinigungsabgeschlossene Familie}{1,5,2}}
  \proofstep{}{a\in\{a,b\}}{\FormulaRefAuto{a \in \{a,b\}}}
  \proofstep{}{b\in\{a,b\}}{\FormulaRefAuto{b \in \{a,b\}}}
  \proofstep{}{a\in X\cup\{a,b\}}{%
    \FormulaRefAuto{z \in B \vdash z \in A \cup B}{7}}
  \proofstep{}{b\in X\cup\{a,b\}}{%
    \FormulaRefAuto{z \in B \vdash z \in A \cup B}{8}}
  \proofstep{1,2,3}{X\cup\{a,b\}\in\FamContaining{M}{a}}{%
    \FormulaRefAuto{\FamContaining{M}{a}\coloneqq
      \{X\in M\mid a\in X\}}{\rAI{6,9}}}
  \proofstep{1,2,3}{X\cup\{a,b\}\in\FamContaining{M}{b}}{%
    \FormulaRefAuto{\FamContaining{M}{a}\coloneqq
      \{X\in M\mid a\in X\}}{\rAI{6,10}}}
  \proofstep{1,2,3}{X\cup\{a,b\}\in
    \FamContaining{M}{a}\cap\FamContaining{M}{b}}{%
    \FormulaRefAuto{x \in A \cap B \eqvdash x \in A\land x \in B}{%
      \rAI{11,12}}}
  \proofstep{1,2}{%
    \forall X\in\FamAvoiding{M}{a}\cap\FamAvoiding{M}{b}\,
      X\cup\{a,b\}\in
        \FamContaining{M}{a}\cap\FamContaining{M}{b}}{%
    \rUI{\rRI{3,13}}}
\end{tabproof}

\FormulaThmDeltaK[Zweieradjunktionsabbildung ist Funktion]{%
  \UnionClosedFam(M)\dsep \{a,b\}\in M\vdash
  U_{M,a,b}\colon
  \bigl(\FamAvoiding{M}{a}\cap\FamAvoiding{M}{b}\bigr)
  \to
  \bigl(\FamContaining{M}{a}\cap\FamContaining{M}{b}\bigr)
}{FranklPairAdjFunction}{%
  \DeltaRow{Mengen}{M\dsep X\dsep a\dsep b}
  \DeltaRow{Funktionen}{U_{M,a,b}}
}
\begin{tabproof}
  \proofstep{1}{\UnionClosedFam(M)}{\rA}
  \proofstep{2}{\{a,b\}\in M}{\rA}
  \proofstep{1,2}{U_{M,a,b}\colon
    \bigl(\FamAvoiding{M}{a}\cap\FamAvoiding{M}{b}\bigr)
    \to
    \bigl(\FamContaining{M}{a}\cap\FamContaining{M}{b}\bigr)}{%
    \FormulaRefAuto{FranklPairAdjDef}[def]{1,2}}
\end{tabproof}

\FormulaThmDeltaK[Auswertung der Zweieradjunktionsabbildung]{%
  \UnionClosedFam(M)\dsep \{a,b\}\in M\dsep
  X\in \FamAvoiding{M}{a}\cap\FamAvoiding{M}{b}
  \vdash U_{M,a,b}(X)=X\cup\{a,b\}
}{FranklPairAdjEval}{%
  \DeltaRow{Mengen}{M\dsep X\dsep a\dsep b}
  \DeltaRow{Funktionen}{U_{M,a,b}}
}
\begin{tabproof}
  \proofstep{1}{\UnionClosedFam(M)}{\rA}
  \proofstep{2}{\{a,b\}\in M}{\rA}
  \proofstep{3}{X\in \FamAvoiding{M}{a}\cap\FamAvoiding{M}{b}}{\rA}
  \proofstep{1,2}{%
    \forall Y\in\FamAvoiding{M}{a}\cap\FamAvoiding{M}{b}\,
      U_{M,a,b}(Y)=Y\cup\{a,b\}}{%
    \FormulaRefAuto{FranklPairAdjDef}[def]{1,2}}
  \proofstep{1,2,3}{U_{M,a,b}(X)=X\cup\{a,b\}}{%
    \rRE{\rUE{4},3}}
\end{tabproof}

\FormulaThmDeltaK[Zweieradjunktionsabbildung ist injektiv]{%
  \UnionClosedFam(M)\dsep \{a,b\}\in M
  \vdash
  U_{M,a,b}\colon
  \bigl(\FamAvoiding{M}{a}\cap\FamAvoiding{M}{b}\bigr)
  \inj
  \bigl(\FamContaining{M}{a}\cap\FamContaining{M}{b}\bigr)
}{FranklPairAdjInjection}{%
  \DeltaRow{Mengen}{M\dsep X\dsep Y\dsep a\dsep b}
  \DeltaRow{Funktionen}{U_{M,a,b}}
}
Zur Abkürzung schreiben wir im folgenden Beweis
\[
  \begin{aligned}
    \mathcal A&\coloneqq
      \FamAvoiding{M}{a}\cap\FamAvoiding{M}{b},&
    \mathcal B&\coloneqq
      \FamContaining{M}{a}\cap\FamContaining{M}{b},\\
    U&\coloneqq U_{M,a,b}.
  \end{aligned}
\]
\begin{tabproof}
  \proofstep{1}{\UnionClosedFam(M)}{\rA}
  \proofstep{2}{\{a,b\}\in M}{\rA}
  \proofstep{1,2}{U\colon\mathcal A\to\mathcal B}{%
    \FormulaRefAuto{FranklPairAdjFunction}{1,2}}
  \proofstep{4}{X\in\mathcal A}{\rA}
  \proofstep{5}{Y\in\mathcal A}{\rA}
  \proofstep{6}{U(X)=U(Y)}{\rA}
  \proofstep{4}{X\in\FamAvoiding{M}{a}}{%
    \FormulaRefAuto{x\in A\cap B\vdash x\in A}{4}}
  \proofstep{4}{X\in\FamAvoiding{M}{b}}{%
    \FormulaRefAuto{x\in A\cap B\vdash x\in B}{4}}
  \proofstep{5}{Y\in\FamAvoiding{M}{a}}{%
    \FormulaRefAuto{x\in A\cap B\vdash x\in A}{5}}
  \proofstep{5}{Y\in\FamAvoiding{M}{b}}{%
    \FormulaRefAuto{x\in A\cap B\vdash x\in B}{5}}
  \proofstep{4}{a\notin X}{%
    \rAEb{\FormulaRefAuto{\FamAvoiding{M}{a}\coloneqq
      \{X\in M\mid a\notin X\}}{7}}}
  \proofstep{4}{b\notin X}{%
    \rAEb{\FormulaRefAuto{\FamAvoiding{M}{a}\coloneqq
      \{X\in M\mid a\notin X\}}{8}}}
  \proofstep{5}{a\notin Y}{%
    \rAEb{\FormulaRefAuto{\FamAvoiding{M}{a}\coloneqq
      \{X\in M\mid a\notin X\}}{9}}}
  \proofstep{5}{b\notin Y}{%
    \rAEb{\FormulaRefAuto{\FamAvoiding{M}{a}\coloneqq
      \{X\in M\mid a\notin X\}}{10}}}
  \proofstep{4}{X\cap\{a,b\}=\varnothing}{%
    \FormulaRefAuto{PairSetDisjointFromTwoAvoided}{11,12}}
  \proofstep{5}{Y\cap\{a,b\}=\varnothing}{%
    \FormulaRefAuto{PairSetDisjointFromTwoAvoided}{13,14}}
  \proofstep{1,2,4}{U(X)=X\cup\{a,b\}}{%
    \FormulaRefAuto{FranklPairAdjEval}{1,2,4}}
  \proofstep{1,2,5}{U(Y)=Y\cup\{a,b\}}{%
    \FormulaRefAuto{FranklPairAdjEval}{1,2,5}}
  \proofstep{1,2,4}{X\cup\{a,b\}=U(X)}{%
    \FormulaRefAuto{a=b\vdash b=a}{17}}
  \proofstep{1,2,4,6}{X\cup\{a,b\}=U(Y)}{%
    \FormulaRefAuto{a = b,\, b = c \vdash a = c}{19,6}}
  \proofstep{1,2,4,5,6}{X\cup\{a,b\}=Y\cup\{a,b\}}{%
    \FormulaRefAuto{a = b,\, b = c \vdash a = c}{20,18}}
  \proofstep{1,2,4,5,6}{X=Y}{%
    \FormulaRefAuto{UnionCancelDisjointRight}{15,16,21}}
  \proofstep{1,2}{U\colon\mathcal A\inj\mathcal B}{%
    \ProofRefStack{%
      \FormulaRefAuto{Injektive Funktion}\\
      \bigl(3,\rUI{\rRI{4,\rRI{5,\rRI{6,22}}}}\bigr)}}
\end{tabproof}

\paragraph{Zusammenführung der Zweige}

\FormulaThmDeltaK[Positiver Mischklassenfall liefert Frankl]{%
  \begin{aligned}[t]
  &\UnionClosedFam(M)\dsep \{a,b\}\in M\dsep a\in\bigcup M\dsep{}\\
  &\exists F\colon
  \bigl(\FamAvoiding{M}{a}\cap\FamContaining{M}{b}\bigr)
  \inj
  \bigl(\FamContaining{M}{a}\cap\FamAvoiding{M}{b}\bigr)
  \vdash \Frankl{M}
  \end{aligned}
}{FranklPairPositiveBranch}{%
  \DeltaRow{Mengen}{M\dsep a\dsep b}
  \DeltaRow{Funktionen}{F\dsep H\dsep U_{M,a,b}}
}
Zur Abkürzung setzen wir
\[
  \begin{aligned}
    \mathcal A_{00}&\coloneqq
      \FamAvoiding{M}{a}\cap\FamAvoiding{M}{b},&
    \mathcal A_{01}&\coloneqq
      \FamAvoiding{M}{a}\cap\FamContaining{M}{b},\\
    \mathcal A_{10}&\coloneqq
      \FamContaining{M}{a}\cap\FamAvoiding{M}{b},&
    \mathcal A_{11}&\coloneqq
      \FamContaining{M}{a}\cap\FamContaining{M}{b}.
  \end{aligned}
\]
\begin{tabproof}
  \proofstep{1}{\UnionClosedFam(M)}{\rA}
  \proofstep{2}{\{a,b\}\in M}{\rA}
  \proofstep{3}{a\in\bigcup M}{\rA}
  \proofstep{4}{\exists F\colon\mathcal A_{01}\inj\mathcal A_{10}}{\rA}
  \proofstep{1,2}{U_{M,a,b}\colon\mathcal A_{00}\inj\mathcal A_{11}}{%
    \FormulaRefAuto{FranklPairAdjInjection}{1,2}}

  \proofstep{6}{F\colon\mathcal A_{01}\inj\mathcal A_{10}}{\rA}
  \proofstep{}{\mathcal A_{00}\cap\mathcal A_{01}=\varnothing}{%
    \FormulaRefAuto{FranklPairAvoidADisjoint}}
  \proofstep{}{\mathcal A_{11}\cap\mathcal A_{10}=\varnothing}{%
    \FormulaRefAuto{FranklPairContainADisjoint}}
  \proofstep{1,2,6}{%
    \exists H\colon
      \mathcal A_{00}\cup\mathcal A_{01}
      \inj
      \mathcal A_{11}\cup\mathcal A_{10}}{%
    \FormulaRefAuto{CaseFunDisjointTargetInjection}{7,8,5,6}}
  \proofstep{10}{%
    H\colon
      \mathcal A_{00}\cup\mathcal A_{01}
      \inj
      \mathcal A_{11}\cup\mathcal A_{10}}{\rA}
  \proofstep{}{\mathcal A_{00}\cup\mathcal A_{01}
    =\FamAvoiding{M}{a}}{%
    \FormulaRefAuto{a=b\vdash b=a}{\FormulaRefAuto{FranklPairAvoidADecomp}}}
  \proofstep{}{\mathcal A_{11}\cup\mathcal A_{10}
    =\FamContaining{M}{a}}{%
    \FormulaRefAuto{a=b\vdash b=a}{\FormulaRefAuto{FranklPairContainADecomp}}}
  \proofstep{10}{H\colon \FamAvoiding{M}{a}\inj\FamContaining{M}{a}}{%
    \rIE{11,\rIE{12,10}}}
  \proofstep{3,10}{\Frankl{M}}{%
    \ProofRefStack{%
      \FormulaRefAuto{\Frankl{M}\coloneqq \exists a\,\HalfOcc{a}{M}}\\
      \FormulaRefAuto{\HalfOcc{a}{M}\coloneqq
        a\in\bigcup M\land
        \exists F\colon \FamAvoiding{M}{a}\inj \FamContaining{M}{a}}\\
      \bigl(\rEI{\rAI{3,\rEI{13}}}\bigr)}}
  \proofstep{1,2,3,6}{\Frankl{M}}{\rEE{9,10,14}}
  \proofstep{1,2,3,4}{\Frankl{M}}{\rEE{4,6,15}}
\end{tabproof}

\FormulaThmDeltaK[Negativer Mischklassenfall liefert Frankl]{%
  \begin{aligned}[t]
  &\Finite{M}\dsep \UnionClosedFam(M)\dsep
  \{a,b\}\in M\dsep b\in\bigcup M\dsep{}\\
  &\neg\exists F\colon
  \bigl(\FamAvoiding{M}{a}\cap\FamContaining{M}{b}\bigr)
  \inj
  \bigl(\FamContaining{M}{a}\cap\FamAvoiding{M}{b}\bigr)
  \vdash \Frankl{M}
  \end{aligned}
}{FranklPairNegativeBranch}{%
  \DeltaRow{Mengen}{M\dsep a\dsep b}
  \DeltaRow{Funktionen}{F\dsep G}
}
\begin{tabproof}
  \proofstep{1}{\Finite{M}}{\rA}
  \proofstep{2}{\UnionClosedFam(M)}{\rA}
  \proofstep{3}{\{a,b\}\in M}{\rA}
  \proofstep{4}{b\in\bigcup M}{\rA}
  \proofstep{5}{\neg\exists F\colon
    \bigl(\FamAvoiding{M}{a}\cap\FamContaining{M}{b}\bigr)
    \inj
    \bigl(\FamContaining{M}{a}\cap\FamAvoiding{M}{b}\bigr)}{\rA}
  \proofstep{1,5}{\exists G\colon
    \bigl(\FamContaining{M}{a}\cap\FamAvoiding{M}{b}\bigr)
    \inj
    \bigl(\FamAvoiding{M}{a}\cap\FamContaining{M}{b}\bigr)}{%
    \FormulaRefAuto{FranklPairMixedReverseInjection}{1,5}}
  \proofstep{7}{G\colon
    \bigl(\FamContaining{M}{a}\cap\FamAvoiding{M}{b}\bigr)
    \inj
    \bigl(\FamAvoiding{M}{a}\cap\FamContaining{M}{b}\bigr)}{\rA}
  \proofstep{}{\{a,b\}=\{b,a\}}{\FormulaRefAuto{\{a,b\} = \{b,a\}}}
  \proofstep{3}{\{b,a\}\in M}{\rIE{8,3}}
  \proofstep{}{
    \FamContaining{M}{a}\cap\FamAvoiding{M}{b}
    =
    \FamAvoiding{M}{b}\cap\FamContaining{M}{a}}{%
    \FormulaRefAuto{A \cap B = B \cap A}}
  \proofstep{}{
    \FamAvoiding{M}{a}\cap\FamContaining{M}{b}
    =
    \FamContaining{M}{b}\cap\FamAvoiding{M}{a}}{%
    \FormulaRefAuto{A \cap B = B \cap A}}
  \proofstep{7}{G\colon
    \bigl(\FamAvoiding{M}{b}\cap\FamContaining{M}{a}\bigr)
    \inj
    \bigl(\FamContaining{M}{b}\cap\FamAvoiding{M}{a}\bigr)}{%
    \rIE{11,\rIE{10,7}}}
  \proofstep{7}{\exists F\colon
    \bigl(\FamAvoiding{M}{b}\cap\FamContaining{M}{a}\bigr)
    \inj
    \bigl(\FamContaining{M}{b}\cap\FamAvoiding{M}{a}\bigr)}{\rEI{12}}
  \proofstep{2,3,4,7}{\Frankl{M}}{%
    \FormulaRefAuto{FranklPairPositiveBranch}{2,9,4,13}}
  \proofstep{1,2,3,4,5}{\Frankl{M}}{\rEE{6,7,14}}
\end{tabproof}

\FormulaThmDeltaK[Frankl-Vermutung bei enthaltener Zweiermenge]{%
  \Finite{M}\dsep \UnionClosedFam(M)\dsep a\neq b\dsep \{a,b\}\in M
  \vdash \Frankl{M}
}{FranklPairMemberFrankl}{%
  \DeltaRow{Mengen}{M\dsep a\dsep b}
  \DeltaRow{Funktionen}{F}
}
\begin{tabproof}
  \proofstep{1}{\Finite{M}}{\rA}
  \proofstep{2}{\UnionClosedFam(M)}{\rA}
  \proofstep{3}{a\neq b}{\rA}
  \proofstep{4}{\{a,b\}\in M}{\rA}
  \proofstep{}{a\in\{a,b\}}{\FormulaRefAuto{a \in \{a,b\}}}
  \proofstep{}{b\in\{a,b\}}{\FormulaRefAuto{b \in \{a,b\}}}
  \proofstep{4}{a\in\bigcup M}{%
    \FormulaRefAuto{B\in A\dsep x\in B\vdash x \in \bigcup A}{4,5}}
  \proofstep{4}{b\in\bigcup M}{%
    \FormulaRefAuto{B\in A\dsep x\in B\vdash x \in \bigcup A}{4,6}}
  \proofstep{}{%
    \begin{aligned}[t]
    &\exists F\colon
    \bigl(\FamAvoiding{M}{a}\cap\FamContaining{M}{b}\bigr)
    \inj
    \bigl(\FamContaining{M}{a}\cap\FamAvoiding{M}{b}\bigr)\\
    &{}\lor
    \neg\exists F\colon
    \bigl(\FamAvoiding{M}{a}\cap\FamContaining{M}{b}\bigr)
    \inj
    \bigl(\FamContaining{M}{a}\cap\FamAvoiding{M}{b}\bigr)
    \end{aligned}}{%
    \FormulaRefAuto{P \lor \neg P}}

  \proofstep{10}{%
    \begin{aligned}[t]
    &\exists F\colon
    \bigl(\FamAvoiding{M}{a}\cap\FamContaining{M}{b}\bigr)
    \inj\\
    &\bigl(\FamContaining{M}{a}\cap\FamAvoiding{M}{b}\bigr)
    \end{aligned}}{\rA}
  \proofstep{2,4,7,10}{\Frankl{M}}{%
    \FormulaRefAuto{FranklPairPositiveBranch}{2,4,7,10}}

  \proofstep{12}{\neg\exists F\colon
    \bigl(\FamAvoiding{M}{a}\cap\FamContaining{M}{b}\bigr)
    \inj
    \bigl(\FamContaining{M}{a}\cap\FamAvoiding{M}{b}\bigr)}{\rA}
  \proofstep{1,2,4,8,12}{\Frankl{M}}{%
    \FormulaRefAuto{FranklPairNegativeBranch}{1,2,4,8,12}}

  \proofstep{1,2,3,4}{\Frankl{M}}{\rOE{9,10,11,12,13}}
\end{tabproof}

\FormulaThmDeltaK[Frankl-Familie mit enthaltener Zweiermenge erfüllt Frankl]{%
  \FranklFam{M}\dsep
  \exists a\,\exists b\,(a\neq b\land \{a,b\}\in M)
  \vdash \Frankl{M}
}{FranklFamPairMemberFrankl}{%
  \DeltaRow{Mengen}{M\dsep a\dsep b}
}
\begin{tabproof}
  \proofstep{1}{\FranklFam{M}}{\rA}
  \proofstep{2}{\exists a\,\exists b\,(a\neq b\land \{a,b\}\in M)}{\rA}
  \proofstep{1}{\Finite{M}}{\FormulaRefAuto{FranklFamFinite}{1}}
  \proofstep{1}{\UnionClosedFam(M)}{\FormulaRefAuto{FranklFamUnionClosed}{1}}

  \proofstep{5}{\exists b\,(a\neq b\land \{a,b\}\in M)}{\rA}
  \proofstep{6}{a\neq b\land \{a,b\}\in M}{\rA}
  \proofstep{6}{a\neq b}{\rAEa{6}}
  \proofstep{6}{\{a,b\}\in M}{\rAEb{6}}
  \proofstep{1,6}{\Frankl{M}}{%
    \FormulaRefAuto{FranklPairMemberFrankl}{3,4,7,8}}
  \proofstep{1,5}{\Frankl{M}}{\rEE{5,6,9}}
  \proofstep{1,2}{\Frankl{M}}{\rEE{2,5,10}}
\end{tabproof}


\subsection{Reduktion auf endliche Mengenglieder über den
Ununterscheidbarkeitsquotienten}

Der Ununterscheidbarkeitsquotient einer Mengenfamilie wurde bereits in
Band~11 entwickelt. Dort stehen auch seine allgemeinen Struktursätze:
\(\OccQuotProj{M}\) ist die surjektive punktweise Projektion auf den
Quotiententräger; die davon induzierte Abbildung
\(\OccQuotMap{M}\) ist eine Bijektion von \(M\) auf
\(\OccQuotFam{M}\), erhält Vereinigungen, und die Spurabbildung
\(\TraceMap{\OccQuotMap{M}}\) bettet den Quotiententräger
\(\OccQuotCarrier{M}\) injektiv in \(\powerset(M)\) ein. Die Symbole
\(\OccClass{M}{a}\), \(\OccQuotCarrier{M}\), \(\OccQuotProj{M}\),
\(\OccQuotMap{M}\) und \(\OccQuotFam{M}\) werden hier daher nur noch
verwendet, nicht erneut definiert.

Für die Frankl-Theorie werden diese allgemeinen Aussagen in zwei Schritten
ausgewertet:
\[
\begin{aligned}
M
&\xrightarrow[\substack{\text{bijektiv und}\\
  \text{vereinigungsverträglich}}]{\ \OccQuotMap{M}\ }
  \OccQuotFam{M},
\\
\TraceMap{\OccQuotMap{M}}
&\colon \OccQuotCarrier{M}\inj\powerset(M).
\end{aligned}
\]
Die erste Zeile überträgt die Familienstruktur. Die zweite macht den
Quotiententräger endlich, sobald \(M\) endlich ist. Danach wird eine
halbhäufige Quotientenklasse durch einen Repräsentanten
\(a\in\bigcup M\) zurück nach \(M\) transportiert.

\subsubsection{Der Quotient einer Frankl-Familie}

Die vier Bestandteile einer Frankl-Familie folgen ohne eigene Hilfssatzkette:
Nichtleerheit und nichtleerer Träger aus den entsprechenden
Struktursätzen des Quotienten, Endlichkeit aus der Bijektion und
Vereinigungsabschluss aus der allgemeinen Unionserhaltung.

\FormulaThmDeltaK[Ununterscheidbarkeitsquotient einer Frankl-Familie ist Frankl-Familie]{%
  \FranklFam{M}\vdash \FranklFam{\OccQuotFam{M}}
}{OccQuotFamFranklFam}{%
  \DeltaRow{Mengen}{M}
  \DeltaRow{Funktionen}{\OccQuotMap{M}}
}
\begin{tabproof}
  \proofstep{1}{\FranklFam{M}}{\rA}
  \proofstep{1}{M\neq\varnothing}{\FormulaRefAuto{FranklFamNonempty}{1}}
  \proofstep{1}{\bigcup M\neq\varnothing}{%
    \FormulaRefAuto{FranklFamCarrierNonempty}{1}}
  \proofstep{1}{\Finite{M}}{\FormulaRefAuto{FranklFamFinite}{1}}
  \proofstep{1}{\UnionClosedFam(M)}{%
    \FormulaRefAuto{FranklFamUnionClosed}{1}}
  \proofstep{1}{\OccQuotFam{M}\neq\varnothing}{%
    \FormulaRefAuto{OccQuotFamNonempty}{2}}
  \proofstep{1}{\bigcup\OccQuotFam{M}\neq\varnothing}{%
    \FormulaRefAuto{OccQuotFamCarrierNonempty}{3}}
  \proofstep{}{\OccQuotMap{M}\colon M\bij\OccQuotFam{M}}{%
    \FormulaRefAuto{OccQuotMapBijective}}
  \proofstep{1}{\Finite{\OccQuotFam{M}}}{%
    \FormulaRefAuto{FiniteBijectionTransport}[thm]{4,8}}
  \proofstep{1}{\UnionClosedFam(\OccQuotFam{M})}{%
    \FormulaRefAuto{OccQuotUnionClosed}{5}}
  \proofstep{1}{\FranklFam{\OccQuotFam{M}}}{%
    \FormulaRefAuto{Frankl-Familie}{6,7,9,10}}
\end{tabproof}

\subsubsection{Endlichkeit und Transport der Frankl-Eigenschaft}

Für die Reduktion bleiben die Endlichkeitsfolgen und die
häufigkeitsbezogenen Aussagen für den Rücktransport zu beweisen. Die dafür
benötigte Injektivität der Spurabbildung ist bereits ein Satz aus Band~11.

\paragraph{Endlichkeit des Quotiententrägers}

Die nächsten Sätze verwenden die injektive Spur, um den Quotiententräger in
eine endliche Potenzmenge einzubetten. Daraus folgt zunächst die Endlichkeit von
\(\OccQuotCarrier{M}\), anschließend die Endlichkeit jedes Familienglieds des
Quotienten. Dieser Schritt ist der eigentliche Grund für die spätere Reduktion
auf Familien mit endlichen Mengengliedern.

\FormulaThmDeltaK[Ununterscheidbarkeitsquotiententräger ist endlich]{%
  \Finite{M}\vdash \Finite{\OccQuotCarrier{M}}
}{OccQuotCarrierFinite}{%
  \DeltaRow{Mengen}{M}
  \DeltaRow{Funktionen}{\TraceMap{\OccQuotMap{M}}}
}
\begin{tabproof}
  \proofstep{1}{\Finite{M}}{\rA}
  \proofstep{1}{\Finite{\powerset(M)}}{%
    \FormulaRefAuto{\Finite{A}\vdash \Finite{\powerset(A)}}{1}}
  \proofstep{}{\TraceMap{\OccQuotMap{M}}\colon
    \OccQuotCarrier{M}\inj\powerset(M)}{%
    \FormulaRefAuto{OccQuotTraceMapInjective}}
  \proofstep{}{\TraceMap{\OccQuotMap{M}}\colon
    \OccQuotCarrier{M}\bij
    \TraceMap{\OccQuotMap{M}}[\OccQuotCarrier{M}]}{%
    \FormulaRefAuto{F\colon A\inj B \vdash F\colon A\bij F[A]}{3}}
  \proofstep{}{\TraceMap{\OccQuotMap{M}}\colon
    \OccQuotCarrier{M}\to\powerset(M)}{%
    \FormulaRefAuto{Injektive Funktion}{3}}
  \proofstep{}{\TraceMap{\OccQuotMap{M}}[\OccQuotCarrier{M}]
    \subseteq\powerset(M)}{%
    \FormulaRefAuto{C\subseteq A \dsep F\colon A\to B \vdash F[C]\subseteq B}{%
    \FormulaRefAuto{A\subseteq A},5}}
  \proofstep{1}{\Finite{\TraceMap{\OccQuotMap{M}}[\OccQuotCarrier{M}]}}{%
    \FormulaRefAuto{FiniteSubset}[thm]{6,2}}
  \proofstep{}{{\TraceMap{\OccQuotMap{M}}}^{-1}\colon
    \TraceMap{\OccQuotMap{M}}[\OccQuotCarrier{M}]
    \bij\OccQuotCarrier{M}}{%
    \FormulaRefAuto{F\colon A\bij B\vdash F^{-1}\colon B\bij A}{4}}
  \proofstep{1}{\Finite{\OccQuotCarrier{M}}}{%
    \FormulaRefAuto{FiniteBijectionTransport}[thm]{7,8}}
\end{tabproof}

\FormulaThmDeltaK[Glieder des Ununterscheidbarkeitsquotienten sind endlich]{%
  \Finite{M}\dsep Y\in\OccQuotFam{M}
  \vdash \Finite{Y}
}{OccQuotMembersFinite}{%
  \DeltaRow{Mengen}{M\dsep Y}
  \DeltaRow{Quotiententräger}{\OccQuotCarrier{M}}
}
\begin{tabproof}
  \proofstep{1}{\Finite{M}}{\rA}
  \proofstep{2}{Y\in\OccQuotFam{M}}{\rA}
  \proofstep{2}{Y\subseteq\OccQuotCarrier{M}}{\FormulaRefAuto{OccQuotMemberSubsetCarrier}{2}}
  \proofstep{1}{\Finite{\OccQuotCarrier{M}}}{\FormulaRefAuto{OccQuotCarrierFinite}{1}}
  \proofstep{1,2}{\Finite{Y}}{\FormulaRefAuto{FiniteSubset}[thm]{3,4}}
\end{tabproof}

\paragraph{Transport der Teilfamilien}

Der folgende Block vergleicht die vermeidenden und enthaltenden Teilfamilien im
Quotienten mit den entsprechenden Teilfamilien von \(M\). Zuerst werden ihre
Mitgliedschaften punktweise übersetzt. Die vermeidende Seite wird danach als
eigener Urbild- und Gleichmächtigkeitsschritt festgehalten; die enthaltende
Seite wird im Rücktransport lokal aus der punktweisen Übersetzung abgeleitet.
Diese Aussagen bereiten den Rücktransport einer Injektion aus dem Quotienten
vor.

\FormulaThmDeltaK[Vermeidende Quotientenmitgliedschaft]{%
  C=\OccClass{M}{a}\dsep a\in\bigcup M\dsep X\in M
  \vdash
  \bigl(\OccQuotMap{M}(X)\in
  \FamAvoiding{\OccQuotFam{M}}{C}
  \leftrightarrow X\in\FamAvoiding{M}{a}\bigr)
}{OccQuotAvoidingMembership}{%
  \DeltaRow{Mengen}{M\dsep X\dsep C\dsep a}
  \DeltaRow{Funktionen}{\OccQuotMap{M}}
  \DeltaRow{Teilmengenfamilien}{%
    \FamAvoiding{\OccQuotFam{M}}{C}\dsep \FamAvoiding{M}{a}}
}
\begin{tabproofsplitwide}
  \proofpartwideindR[\(\rightarrow\)]{%
    C=\OccClass{M}{a}\dsep a\in\bigcup M\dsep X\in M
    \vdash
    \bigl(\OccQuotMap{M}(X)\in
    \FamAvoiding{\OccQuotFam{M}}{C}
    \rightarrow X\in\FamAvoiding{M}{a}\bigr)
  }{%
    C=\OccClass{M}{a}\dsep a\in\bigcup M\dsep X\in M
    \vdash
    \bigl(\OccQuotMap{M}(X)\in
    \FamAvoiding{\OccQuotFam{M}}{C}
    \rightarrow X\in\FamAvoiding{M}{a}\bigr)
  }[OccQuotAvoidingMembershipForward]
    \proofstepwidestar[1]{C=\OccClass{M}{a}}{\rA}
    \proofstepwidestar[2]{a\in\bigcup M}{\rA}
    \proofstepwidestar[3]{X\in M}{\rA}
    \proofstepwidestar[4]{%
    \OccQuotMap{M}(X)\in
    \FamAvoiding{\OccQuotFam{M}}{C}}{\rA}
    \proofstepwidestar[4]{\OccQuotMap{M}(X)\in\OccQuotFam{M}\land
    C\notin\OccQuotMap{M}(X)}{%
    \FormulaRefAuto{\FamAvoiding{M}{a}\coloneqq
    \{X\in M\mid a\notin X\}}{4}}
    \proofstepwidestar[4]{C\notin\OccQuotMap{M}(X)}{\rAEb{5}}
    \proofstepwidestar[7]{a\in X}{\rA}
    \proofstepwidestar[2,3,7]{\OccClass{M}{a}\in\OccQuotMap{M}(X)}{%
    \FormulaRefAuto{OccQuotMembershipTransport}{3,2,7}}
    \proofstepwidestar[1,2,3,7]{C\in\OccQuotMap{M}(X)}{\rIE{1,8}}
    \proofstepwidestar[1,2,3,4,7]{\bot}{\rBI{9,6}}
    \proofstepwidestar[1,2,3,4]{a\notin X}{\rCI{7,10}}
    \proofstepwidestar[1,2,3,4]{X\in\FamAvoiding{M}{a}}{%
    \FormulaRefAuto{\FamAvoiding{M}{a}\coloneqq
    \{X\in M\mid a\notin X\}}{\rAI{3,11}}}
    \proofstepwidestar[1,2,3]{%
    \OccQuotMap{M}(X)\in
    \FamAvoiding{\OccQuotFam{M}}{C}
    \rightarrow X\in\FamAvoiding{M}{a}}{\rRI{4,12}}
  \closeproofpartwideindR

  \proofpartwideindR[\(\leftarrow\)]{%
    C=\OccClass{M}{a}\dsep a\in\bigcup M\dsep X\in M
    \vdash
    \bigl(X\in\FamAvoiding{M}{a}\rightarrow
    \OccQuotMap{M}(X)\in
    \FamAvoiding{\OccQuotFam{M}}{C}\bigr)
  }{%
    C=\OccClass{M}{a}\dsep a\in\bigcup M\dsep X\in M
    \vdash
    \bigl(X\in\FamAvoiding{M}{a}\rightarrow
    \OccQuotMap{M}(X)\in
    \FamAvoiding{\OccQuotFam{M}}{C}\bigr)
  }[OccQuotAvoidingMembershipBackward]
    \proofstepwidestar[1]{C=\OccClass{M}{a}}{\rA}
    \proofstepwidestar[2]{a\in\bigcup M}{\rA}
    \proofstepwidestar[3]{X\in M}{\rA}
    \proofstepwidestar[4]{X\in\FamAvoiding{M}{a}}{\rA}
    \proofstepwidestar[4]{X\in M\land a\notin X}{%
    \FormulaRefAuto{\FamAvoiding{M}{a}\coloneqq
    \{X\in M\mid a\notin X\}}{4}}
    \proofstepwidestar[4]{a\notin X}{\rAEb{5}}
    \proofstepwidestar[3]{\OccQuotMap{M}(X)\in\OccQuotFam{M}}{%
    \FormulaRefAuto{OccQuotSetInFam}{3}}
    \proofstepwidestar[8]{C\in\OccQuotMap{M}(X)}{\rA}
    \proofstepwidestar[1,8]{\OccClass{M}{a}\in\OccQuotMap{M}(X)}{%
    \rIE{\FormulaRefAuto{a=b\vdash b=a}{1},8}}
    \proofstepwidestar[2,3,8]{a\in X}{%
    \FormulaRefAuto{OccQuotMembershipTransport}{3,2,9}}
    \proofstepwidestar[1,2,3,4,8]{\bot}{\rBI{10,6}}
    \proofstepwidestar[1,2,3,4]{C\notin\OccQuotMap{M}(X)}{\rCI{8,11}}
    \proofstepwidestar[1,2,3,4]{%
    \OccQuotMap{M}(X)\in
    \FamAvoiding{\OccQuotFam{M}}{C}}{%
    \FormulaRefAuto{\FamAvoiding{M}{a}\coloneqq
    \{X\in M\mid a\notin X\}}{\rAI{7,12}}}
    \proofstepwidestar[1,2,3]{%
    X\in\FamAvoiding{M}{a}\rightarrow
    \OccQuotMap{M}(X)\in
    \FamAvoiding{\OccQuotFam{M}}{C}}{\rRI{4,13}}
  \closeproofpartwideindR

  \proofpartwide{Schluss}
    \proofstepwidestar[1]{C=\OccClass{M}{a}}{\rA}
    \proofstepwidestar[2]{a\in\bigcup M}{\rA}
    \proofstepwidestar[3]{X\in M}{\rA}
    \proofstepwidestar[1,2,3]{%
    \OccQuotMap{M}(X)\in
    \FamAvoiding{\OccQuotFam{M}}{C}
    \rightarrow X\in\FamAvoiding{M}{a}}{%
    \FormulaRefAuto{OccQuotAvoidingMembershipForward}{1,2,3}}
    \proofstepwidestar[1,2,3]{%
    X\in\FamAvoiding{M}{a}\rightarrow
    \OccQuotMap{M}(X)\in
    \FamAvoiding{\OccQuotFam{M}}{C}}{%
    \FormulaRefAuto{OccQuotAvoidingMembershipBackward}{1,2,3}}
    \proofstepwidestar[1,2,3]{%
    \OccQuotMap{M}(X)\in
    \FamAvoiding{\OccQuotFam{M}}{C}
    \leftrightarrow X\in\FamAvoiding{M}{a}}{\rLRI{4,5}}
  \closeproofpartwide
\end{tabproofsplitwide}

\FormulaThmDeltaK[Enthaltende Quotientenmitgliedschaft]{%
  C=\OccClass{M}{a}\dsep a\in\bigcup M\dsep X\in M
  \vdash
  \bigl(\OccQuotMap{M}(X)\in
  \FamContaining{\OccQuotFam{M}}{C}
  \leftrightarrow X\in\FamContaining{M}{a}\bigr)
}{OccQuotContainingMembership}{%
  \DeltaRow{Mengen}{M\dsep X\dsep C\dsep a}
  \DeltaRow{Funktionen}{\OccQuotMap{M}}
  \DeltaRow{Teilmengenfamilien}{%
    \FamContaining{\OccQuotFam{M}}{C}\dsep \FamContaining{M}{a}}
}
\begin{tabproofsplitwide}
  \proofpartwideindR[\(\rightarrow\)]{%
    C=\OccClass{M}{a}\dsep a\in\bigcup M\dsep X\in M
    \vdash
    \bigl(\OccQuotMap{M}(X)\in
    \FamContaining{\OccQuotFam{M}}{C}
    \rightarrow X\in\FamContaining{M}{a}\bigr)
  }{%
    C=\OccClass{M}{a}\dsep a\in\bigcup M\dsep X\in M
    \vdash
    \bigl(\OccQuotMap{M}(X)\in
    \FamContaining{\OccQuotFam{M}}{C}
    \rightarrow X\in\FamContaining{M}{a}\bigr)
  }[OccQuotContainingMembershipForward]
    \proofstepwidestar[1]{C=\OccClass{M}{a}}{\rA}
    \proofstepwidestar[2]{a\in\bigcup M}{\rA}
    \proofstepwidestar[3]{X\in M}{\rA}
    \proofstepwidestar[4]{%
    \OccQuotMap{M}(X)\in
    \FamContaining{\OccQuotFam{M}}{C}}{\rA}
    \proofstepwidestar[4]{\OccQuotMap{M}(X)\in\OccQuotFam{M}\land
    C\in\OccQuotMap{M}(X)}{%
    \FormulaRefAuto{\FamContaining{M}{a}\coloneqq
    \{X\in M\mid a\in X\}}{4}}
    \proofstepwidestar[4]{C\in\OccQuotMap{M}(X)}{\rAEb{5}}
    \proofstepwidestar[1,4]{\OccClass{M}{a}\in\OccQuotMap{M}(X)}{%
    \rIE{\FormulaRefAuto{a=b\vdash b=a}{1},6}}
    \proofstepwidestar[2,3,4]{a\in X}{%
    \FormulaRefAuto{OccQuotMembershipTransport}{3,2,7}}
    \proofstepwidestar[1,2,3,4]{X\in\FamContaining{M}{a}}{%
    \FormulaRefAuto{\FamContaining{M}{a}\coloneqq
    \{X\in M\mid a\in X\}}{\rAI{3,8}}}
    \proofstepwidestar[1,2,3]{%
    \OccQuotMap{M}(X)\in
    \FamContaining{\OccQuotFam{M}}{C}
    \rightarrow X\in\FamContaining{M}{a}}{\rRI{4,9}}
  \closeproofpartwideindR

  \proofpartwideindR[\(\leftarrow\)]{%
    C=\OccClass{M}{a}\dsep a\in\bigcup M\dsep X\in M
    \vdash
    \bigl(X\in\FamContaining{M}{a}\rightarrow
    \OccQuotMap{M}(X)\in
    \FamContaining{\OccQuotFam{M}}{C}\bigr)
  }{%
    C=\OccClass{M}{a}\dsep a\in\bigcup M\dsep X\in M
    \vdash
    \bigl(X\in\FamContaining{M}{a}\rightarrow
    \OccQuotMap{M}(X)\in
    \FamContaining{\OccQuotFam{M}}{C}\bigr)
  }[OccQuotContainingMembershipBackward]
    \proofstepwidestar[1]{C=\OccClass{M}{a}}{\rA}
    \proofstepwidestar[2]{a\in\bigcup M}{\rA}
    \proofstepwidestar[3]{X\in M}{\rA}
    \proofstepwidestar[4]{X\in\FamContaining{M}{a}}{\rA}
    \proofstepwidestar[4]{X\in M\land a\in X}{%
    \FormulaRefAuto{\FamContaining{M}{a}\coloneqq
    \{X\in M\mid a\in X\}}{4}}
    \proofstepwidestar[4]{a\in X}{\rAEb{5}}
    \proofstepwidestar[3]{\OccQuotMap{M}(X)\in\OccQuotFam{M}}{%
    \FormulaRefAuto{OccQuotSetInFam}{3}}
    \proofstepwidestar[2,3,4]{\OccClass{M}{a}\in\OccQuotMap{M}(X)}{%
    \FormulaRefAuto{OccQuotMembershipTransport}{3,2,6}}
    \proofstepwidestar[1,2,3,4]{C\in\OccQuotMap{M}(X)}{\rIE{1,8}}
    \proofstepwidestar[1,2,3,4]{%
    \OccQuotMap{M}(X)\in
    \FamContaining{\OccQuotFam{M}}{C}}{%
    \FormulaRefAuto{\FamContaining{M}{a}\coloneqq
    \{X\in M\mid a\in X\}}{\rAI{7,9}}}
    \proofstepwidestar[1,2,3]{%
    X\in\FamContaining{M}{a}\rightarrow
    \OccQuotMap{M}(X)\in
    \FamContaining{\OccQuotFam{M}}{C}}{\rRI{4,10}}
  \closeproofpartwideindR

  \proofpartwide{Schluss}
    \proofstepwidestar[1]{C=\OccClass{M}{a}}{\rA}
    \proofstepwidestar[2]{a\in\bigcup M}{\rA}
    \proofstepwidestar[3]{X\in M}{\rA}
    \proofstepwidestar[1,2,3]{%
    \OccQuotMap{M}(X)\in
    \FamContaining{\OccQuotFam{M}}{C}
    \rightarrow X\in\FamContaining{M}{a}}{%
    \FormulaRefAuto{OccQuotContainingMembershipForward}{1,2,3}}
    \proofstepwidestar[1,2,3]{%
    X\in\FamContaining{M}{a}\rightarrow
    \OccQuotMap{M}(X)\in
    \FamContaining{\OccQuotFam{M}}{C}}{%
    \FormulaRefAuto{OccQuotContainingMembershipBackward}{1,2,3}}
    \proofstepwidestar[1,2,3]{%
    \OccQuotMap{M}(X)\in
    \FamContaining{\OccQuotFam{M}}{C}
    \leftrightarrow X\in\FamContaining{M}{a}}{\rLRI{4,5}}
  \closeproofpartwide
\end{tabproofsplitwide}

\FormulaThmDeltaK[Urbild der vermeidenden Quotiententeilfamilie]{%
  C=\OccClass{M}{a}\dsep a\in\bigcup M
  \vdash
  \bigl(\OccQuotMap{M}\bigr)^{-1}[\FamAvoiding{\OccQuotFam{M}}{C}]
  =
  \FamAvoiding{M}{a}
}{OccQuotAvoidingPreimage}{%
  \DeltaRow{Mengen}{M\dsep X\dsep C\dsep a}
  \DeltaRow{Funktionen}{\OccQuotMap{M}}
  \DeltaRow{Teilmengenfamilien}{%
    \FamAvoiding{\OccQuotFam{M}}{C}\dsep \FamAvoiding{M}{a}}
}
\begin{tabproof}
  \proofstep{1}{C=\OccClass{M}{a}}{\rA}
  \proofstep{2}{a\in\bigcup M}{\rA}
  \proofstep{}{\OccQuotMap{M}\colon M\to\OccQuotFam{M}}{%
    \FormulaRefAuto{OccQuotMapFunction}}
  \proofstep{}{%
    \FamAvoiding{\OccQuotFam{M}}{C}\subseteq \OccQuotFam{M}}{%
    \FormulaRefAuto{\FamAvoiding{M}{a}\subseteq M}}
  \proofstep{}{\FamAvoiding{M}{a}\subseteq M}{%
    \FormulaRefAuto{\FamAvoiding{M}{a}\subseteq M}}
  \proofstep{6}{X\in M}{\rA}
  \proofstep{1,2,6}{%
    \OccQuotMap{M}(X)\in
    \FamAvoiding{\OccQuotFam{M}}{C}
    \leftrightarrow X\in\FamAvoiding{M}{a}}{%
    \FormulaRefAuto{OccQuotAvoidingMembership}{1,2,6}}
  \proofstep{1,2}{%
    \forall X\in M\,\bigl(
    \OccQuotMap{M}(X)\in
    \FamAvoiding{\OccQuotFam{M}}{C}
    \leftrightarrow X\in\FamAvoiding{M}{a}\bigr)}{%
    \rUI{\rRI{6,7}}}
  \proofstep{1,2}{%
    \bigl(\OccQuotMap{M}\bigr)^{-1}[\FamAvoiding{\OccQuotFam{M}}{C}]
    =
    \FamAvoiding{M}{a}}{%
    \FormulaRefAuto{PreimageEqFromPointwiseIff}{3,4,5,8}}
\end{tabproof}

\FormulaThmDeltaK[Vermeidende Teilfamilien sind unter dem Quotienten gleichmächtig]{%
  C=\OccClass{M}{a}\dsep a\in\bigcup M
  \vdash
  \FamAvoiding{\OccQuotFam{M}}{C}\EqCard
  \FamAvoiding{M}{a}
}{OccQuotAvoidingEqCard}{%
  \DeltaRow{Mengen}{M\dsep C\dsep a}
  \DeltaRow{Funktionen}{H\dsep \OccQuotMap{M}}
  \DeltaRow{Teilmengenfamilien}{%
    \FamAvoiding{\OccQuotFam{M}}{C}\dsep \FamAvoiding{M}{a}}
}
\begin{tabproof}
  \proofstep{1}{C=\OccClass{M}{a}}{\rA}
  \proofstep{2}{a\in\bigcup M}{\rA}
  \proofstep{1,2}{%
    \bigl(\OccQuotMap{M}\bigr)^{-1}[\FamAvoiding{\OccQuotFam{M}}{C}]
    =
    \FamAvoiding{M}{a}}{%
    \FormulaRefAuto{OccQuotAvoidingPreimage}{1,2}}
  \proofstep{}{\OccQuotMap{M}\colon M\bij\OccQuotFam{M}}{%
    \FormulaRefAuto{OccQuotMapBijective}}
  \proofstep{}{\FamAvoiding{\OccQuotFam{M}}{C}\subseteq
    \OccQuotFam{M}}{%
    \FormulaRefAuto{\FamAvoiding{M}{a}\subseteq M}}
  \proofstep{1,2}{\exists H\colon
    \bigl(\OccQuotMap{M}\bigr)^{-1}[\FamAvoiding{\OccQuotFam{M}}{C}]
    \bij
    \FamAvoiding{\OccQuotFam{M}}{C}}{%
    \rEI{\FormulaRefAuto{F\colon A\bij B \dsep T\subseteq B \vdash F\restriction_{F^{-1}[T]}\colon F^{-1}[T]\bij T}{4,5}}}
  \proofstep{1,2}{%
    \bigl(\OccQuotMap{M}\bigr)^{-1}[\FamAvoiding{\OccQuotFam{M}}{C}]
    \EqCard
    \FamAvoiding{\OccQuotFam{M}}{C}}{%
    \FormulaRefAuto{A \EqCard B \coloneqq \exists F\colon A\bij B}{6}}
  \proofstep{1,2}{%
    \FamAvoiding{M}{a}\EqCard
    \FamAvoiding{\OccQuotFam{M}}{C}}{\rIE{3,7}}
  \proofstep{1,2}{%
    \FamAvoiding{\OccQuotFam{M}}{C}\EqCard
    \FamAvoiding{M}{a}}{%
    \FormulaRefAuto{EqCardSymmetric}{8}}
\end{tabproof}

\FormulaThmDeltaK[Quotienteninjektion zieht sich auf Teilfamilien zurück]{%
  \begin{aligned}[t]
    &C=\OccClass{M}{a}\dsep a\in\bigcup M\dsep
    F\colon \FamAvoiding{\OccQuotFam{M}}{C}\inj
    \FamContaining{\OccQuotFam{M}}{C}\\
    &\vdash
    \exists G\colon \FamAvoiding{M}{a}\inj\FamContaining{M}{a}
  \end{aligned}
}{OccQuotInjectionPullback}{%
  \DeltaRow{Mengen}{M\dsep C\dsep a}
  \DeltaRow{Funktionen}{F\dsep G\dsep \OccQuotMap{M}}
  \DeltaRow{Teilmengenfamilien}{%
    \FamAvoiding{\OccQuotFam{M}}{C}\dsep
    \FamContaining{\OccQuotFam{M}}{C}\dsep
    \FamAvoiding{M}{a}\dsep \FamContaining{M}{a}}
}
\begin{tabproof}
  \proofstep{1}{C=\OccClass{M}{a}}{\rA}
  \proofstep{2}{a\in\bigcup M}{\rA}
  \proofstep{3}{F\colon \FamAvoiding{\OccQuotFam{M}}{C}\inj
    \FamContaining{\OccQuotFam{M}}{C}}{\rA}
  \proofstep{1,2}{%
    \FamAvoiding{\OccQuotFam{M}}{C}\EqCard
    \FamAvoiding{M}{a}}{%
    \FormulaRefAuto{OccQuotAvoidingEqCard}{1,2}}
  \proofstep{}{\OccQuotMap{M}\colon M\to\OccQuotFam{M}}{%
    \FormulaRefAuto{OccQuotMapFunction}}
  \proofstep{}{\FamContaining{\OccQuotFam{M}}{C}\subseteq
    \OccQuotFam{M}}{%
    \FormulaRefAuto{\FamContaining{M}{a}\subseteq M}}
  \proofstep{}{\FamContaining{M}{a}\subseteq M}{%
    \FormulaRefAuto{\FamContaining{M}{a}\subseteq M}}
  \proofstep{8}{X\in M}{\rA}
  \proofstep{1,2,8}{%
    \OccQuotMap{M}(X)\in
    \FamContaining{\OccQuotFam{M}}{C}
    \leftrightarrow X\in\FamContaining{M}{a}}{%
    \FormulaRefAuto{OccQuotContainingMembership}{1,2,8}}
  \proofstep{1,2}{%
    \forall X\in M\,\bigl(
    \OccQuotMap{M}(X)\in
    \FamContaining{\OccQuotFam{M}}{C}
    \leftrightarrow X\in\FamContaining{M}{a}\bigr)}{%
    \rUI{\rRI{8,9}}}
  \proofstep{1,2}{%
    \bigl(\OccQuotMap{M}\bigr)^{-1}[\FamContaining{\OccQuotFam{M}}{C}]
    =
    \FamContaining{M}{a}}{%
    \FormulaRefAuto{PreimageEqFromPointwiseIff}{5,6,7,10}}
  \proofstep{}{\OccQuotMap{M}\colon M\bij\OccQuotFam{M}}{%
    \FormulaRefAuto{OccQuotMapBijective}}
  \proofstep{1,2}{\exists H\colon
    \bigl(\OccQuotMap{M}\bigr)^{-1}[\FamContaining{\OccQuotFam{M}}{C}]
    \bij
    \FamContaining{\OccQuotFam{M}}{C}}{%
    \rEI{\FormulaRefAuto{F\colon A\bij B \dsep T\subseteq B \vdash F\restriction_{F^{-1}[T]}\colon F^{-1}[T]\bij T}{12,6}}}
  \proofstep{1,2}{%
    \bigl(\OccQuotMap{M}\bigr)^{-1}[\FamContaining{\OccQuotFam{M}}{C}]
    \EqCard
    \FamContaining{\OccQuotFam{M}}{C}}{%
    \FormulaRefAuto{A \EqCard B \coloneqq \exists F\colon A\bij B}{13}}
  \proofstep{1,2}{%
    \FamContaining{M}{a}\EqCard
    \FamContaining{\OccQuotFam{M}}{C}}{\rIE{11,14}}
  \proofstep{1,2}{%
    \FamContaining{\OccQuotFam{M}}{C}\EqCard
    \FamContaining{M}{a}}{%
    \FormulaRefAuto{EqCardSymmetric}{15}}
  \proofstep{1,2,3}{%
    \exists G\colon \FamAvoiding{M}{a}\inj
    \FamContaining{M}{a}}{%
    \FormulaRefAuto{EqCardInjectionTransport}{4,16,3}}
\end{tabproof}

\FormulaThmDeltaK[Quotienteninjektion liefert Halbhäufigkeit eines Repräsentanten]{%
  \begin{aligned}[t]
    &C=\OccClass{M}{a}\dsep a\in\bigcup M\dsep
    \exists F\colon \FamAvoiding{\OccQuotFam{M}}{C}\inj
    \FamContaining{\OccQuotFam{M}}{C}\\
    &\vdash \HalfOcc{a}{M}
  \end{aligned}
}{OccQuotRepresentativeHalfOcc}{%
  \DeltaRow{Mengen}{M\dsep C\dsep a}
  \DeltaRow{Funktionen}{F\dsep G}
  \DeltaRow{Teilmengenfamilien}{\FamAvoiding{M}{a}\dsep
    \FamContaining{M}{a}}
}
\begin{tabproofwide}
  \proofstepwidestar[1]{C=\OccClass{M}{a}}{\rA}
  \proofstepwidestar[2]{a\in\bigcup M}{\rA}
  \proofstepwidestar[3]{\exists F\colon
    \FamAvoiding{\OccQuotFam{M}}{C}\inj
    \FamContaining{\OccQuotFam{M}}{C}}{\rA}

  \proofstepwidestar[5]{F\colon
    \FamAvoiding{\OccQuotFam{M}}{C}\inj
    \FamContaining{\OccQuotFam{M}}{C}}{\rA}
  \proofstepwidestar[1,2,5]{%
    \exists G\colon \FamAvoiding{M}{a}\inj\FamContaining{M}{a}}{%
    \FormulaRefAuto{OccQuotInjectionPullback}{1,2,5}}
  \proofstepwidestar[1,2,5]{\HalfOcc{a}{M}}{%
    \FormulaRefAuto{Halbhäufiges Element}{\rAI{2,6}}}
  \proofstepwidestar[1,2,3]{\HalfOcc{a}{M}}{\rEE{3,5,7}}
\end{tabproofwide}

\FormulaThmDeltaK[Halbhäufigkeit transportiert vom Ununterscheidbarkeitsquotienten]{%
  \HalfOcc{C}{\OccQuotFam{M}}\vdash \Frankl{M}
}{OccQuotHalfOccTransport}{%
  \DeltaRow{Mengen}{M\dsep X\dsep Y\dsep C\dsep a}
  \DeltaRow{Funktionen}{F\dsep G\dsep \OccQuotMap{M}}
  \DeltaRow{Teilmengenfamilien}{\FamContaining{M}{a}\dsep \FamAvoiding{M}{a}}
}
\begin{tabproofwide}
  \proofstepwidestar[1]{\HalfOcc{C}{\OccQuotFam{M}}}{\rA}
  \proofstepwidestar[1]{C\in\bigcup\OccQuotFam{M}\land
    \exists F\colon \FamAvoiding{\OccQuotFam{M}}{C}\inj
    \FamContaining{\OccQuotFam{M}}{C}}{%
    \FormulaRefAuto{Halbhäufiges Element}{1}}
  \proofstepwidestar[1]{C\in\bigcup\OccQuotFam{M}}{\rAEa{2}}
  \proofstepwidestar[1]{\exists F\colon
    \FamAvoiding{\OccQuotFam{M}}{C}\inj
    \FamContaining{\OccQuotFam{M}}{C}}{\rAEb{2}}
  \proofstepwidestar[1]{\exists a\in\bigcup M\,C=\OccClass{M}{a}}{%
    \FormulaRefAuto{OccQuotUnionRepresentative}{3}}

  \proofstepwidestar[6]{a\in\bigcup M\land C=\OccClass{M}{a}}{\rA}
  \proofstepwidestar[6]{a\in\bigcup M}{\rAEa{6}}
  \proofstepwidestar[6]{C=\OccClass{M}{a}}{\rAEb{6}}
  \proofstepwidestar[1,6]{\HalfOcc{a}{M}}{%
    \FormulaRefAuto{OccQuotRepresentativeHalfOcc}{8,7,4}}
  \proofstepwidestar[1,6]{\exists a\,\HalfOcc{a}{M}}{\rEI{9}}
  \proofstepwidestar[1,6]{\Frankl{M}}{%
    \FormulaRefAuto{Frankl-Eigenschaft}{10}}
  \proofstepwidestar[1]{\Frankl{M}}{\rEE{5,6,11}}
\end{tabproofwide}

\FormulaThmDeltaK[Frankl-Eigenschaft transportiert vom Ununterscheidbarkeitsquotienten]{%
  \Frankl{\OccQuotFam{M}}\vdash \Frankl{M}
}{OccQuotFranklTransport}{%
  \DeltaRow{Mengen}{M\dsep C\dsep a}
  \DeltaRow{Funktionen}{F\dsep G\dsep \OccQuotMap{M}}
  \DeltaRow{Teilmengenfamilien}{\FamContaining{M}{a}\dsep \FamAvoiding{M}{a}}
}
\begin{tabproof}
  \proofstep{1}{\Frankl{\OccQuotFam{M}}}{\rA}
  \proofstep{1}{\exists C\,\HalfOcc{C}{\OccQuotFam{M}}}{%
    \FormulaRefAuto{Frankl-Eigenschaft}{1}}
  \proofstep{3}{\HalfOcc{C}{\OccQuotFam{M}}}{\rA}
  \proofstep{3}{\Frankl{M}}{%
    \FormulaRefAuto{OccQuotHalfOccTransport}{3}}
  \proofstep{1}{\Frankl{M}}{\rEE{2,3,4}}
\end{tabproof}

Damit ist gezeigt, dass eine Frankl-Lösung im Quotienten wieder eine
Frankl-Lösung in der ursprünglichen Familie liefert. Der Nachweis bleibt dabei
vollständig formal: Die halbhäufige Quotientenklasse wird durch einen
Repräsentanten aus \(\bigcup M\) zurückgeführt.

\subsubsection{Reduktion auf endliche Mengenglieder}

Der letzte Block fasst die vorherigen Konstruktionen zur Reduktion zusammen.
Unter der Annahme, dass alle Frankl-Familien mit endlichen Mengengliedern die
Frankl-Eigenschaft besitzen, wird eine beliebige Frankl-Familie \(M\) durch
ihren Ununterscheidbarkeitsquotienten ersetzt. Für diesen Quotienten sind die
benötigten endlichen Mengenglieder bereits bewiesen, und die Frankl-Eigenschaft
wird anschließend nach \(M\) zurücktransportiert.

\FormulaThmDeltaK[Reduktion auf endliche Mengenglieder]{%
  \begin{aligned}[t]
    &\forall N\,\Bigl(
      \FranklFam{N}\land
      \forall Y\in N\,\Finite{Y}
      \rightarrow \Frankl{N}\Bigr)\\
    &\dsep \FranklFam{M}
    \vdash \Frankl{M}
  \end{aligned}
}{FranklFiniteMembersReduction}{%
  \DeltaRow{Mengen}{M\dsep N\dsep Y}
  \DeltaRow{Ununterscheidbarkeitsquotient}{\OccQuotFam{M}}
}
\begin{tabproof}
  \proofstep{1}{%
    \begin{aligned}[t]
      \forall N\,\Bigl(
      &\FranklFam{N}\land{}\\
      &\forall Y\in N\,\Finite{Y}
      \rightarrow \Frankl{N}\Bigr)
  \end{aligned}}{\rA}
  \proofstep{2}{\FranklFam{M}}{\rA}
  \proofstep{2}{\Finite{M}}{\FormulaRefAuto{FranklFamFinite}{2}}
  \proofstep{2}{\FranklFam{\OccQuotFam{M}}}{%
    \FormulaRefAuto{OccQuotFamFranklFam}{2}}
  \proofstep{2}{\forall Y\in\OccQuotFam{M}\,\Finite{Y}}{%
    \rUI{\rRI{\FormulaRefAuto{OccQuotMembersFinite}{3}}}}
  \proofstep{2}{%
    \begin{aligned}[t]
      &\FranklFam{\OccQuotFam{M}}\land{}\\
      &\forall Y\in\OccQuotFam{M}\,\Finite{Y}
    \end{aligned}}{\rAI{4,5}}
  \proofstep{1}{%
    \begin{aligned}[t]
      &\FranklFam{\OccQuotFam{M}}\land{}\\
      &\forall Y\in\OccQuotFam{M}\,\Finite{Y}
      \rightarrow \Frankl{\OccQuotFam{M}}
    \end{aligned}}{\rUE{1}}
  \proofstep{1,2}{\Frankl{\OccQuotFam{M}}}{\rRE{7,6}}
  \proofstep{1,2}{\Frankl{M}}{\FormulaRefAuto{OccQuotFranklTransport}{8}}
\end{tabproof}

Die Reduktion ist damit vollständig. Der Quotient einer beliebigen
Frankl-Familie ist wieder eine Frankl-Familie, alle seine Mengenglieder sind
endlich, und die Frankl-Eigenschaft kann vom Quotienten zurück auf die
ursprüngliche Familie transportiert werden.


\section{Normalisierung durch Adjunktion der leeren Menge}

Für die spätere Reduktion verwenden wir die Kurzschreibweise
\(M^\circ=M\cup\{\varnothing\}\). Die zugehörigen Erhaltungs- und
Transportsätze werden vor dem globalen Äquivalenzsatz vollständig bewiesen.

\FormulaDefDeltaK[Erweiterung einer Familie um die leere Menge]{%
  M^\circ\coloneqq M\cup\{\varnothing\}%
}{FranklEmptyAdjunction}{%
  \DeltaRow{Mengen}{M}
  \DeltaRow{\textbf{Neues Symbol}}{}
  \DeltaPrem{Erweiterte Familie}{M^\circ}
}

\chapter{Frankls Vermutung in Halbverbandsform}

Die benötigte Theorie der Supremums-Halbverbände mit Null und der echten
beschränkten Supremums-Halbverbände steht nun unabhängig in Band~34. Im
endlichen Fall liefert der dortige Beschränktheitssatz automatisch ein
eindeutiges größtes Element.

\FormulaDefDeltaK[Seltenes Element einer endlichen Familie]{%
  \RareOcc{a}{M}\coloneqq
  a\in\bigcup M\land
  \exists F\colon \FamContaining{M}{a}\inj\FamAvoiding{M}{a}
}{RareOccurrence}{%
  \DeltaRow{Mengen}{M\dsep a}
  \DeltaRow{Funktionen}{F}
  \DeltaRow{Teilmengenfamilien}{%
    \FamContaining{M}{a}\dsep\FamAvoiding{M}{a}}
  \DeltaRow{\textbf{Neues Symbol}}{}
  \DeltaPrem{Seltene Elemente}{\RareOcc{a}{M}}
}

\FormulaDefDeltaK[Frankl-Eigenschaft]{%
  \FranklSemi{A,\JoinOp,0}\coloneqq
  \exists j\Bigl(
    \JoinIrr{A,\JoinOp,0}{j}\land
    \AtMostHalf{\PrincipalFilter{A,\InducedLeq}{j}}{A}
  \Bigr)
}{FranklSemilatticeProperty}{%
  \DeltaRow{Mengen}{A\dsep 0\dsep j}
  \DeltaRow{Funktionen}{\JoinOp\dsep F}
  \DeltaPrem{Halbverband mit Nullelement}{\ZeroJoinSemi{A,\JoinOp,0}}
  \DeltaRow{Hauptfilter}{\PrincipalFilter{A,\InducedLeq}{j}}
  \DeltaRow{\textbf{Neues Symbol}}{}
  \DeltaPrem{Frankl-Halbverbände}{\FranklSemi{A,\JoinOp,0}}
}

Die rein halbverbandstheoretische Frankl-Vermutung lautet damit:

\begin{samepage}
\begin{center}
\fbox{\begin{minipage}{0.88\linewidth}
\textbf{Frankl-Vermutung für Supremums-Halbverbände.}
Ist \((A,\JoinOp,0)\) ein endlicher Supremums-Halbverband mit Nullelement,
so ist er nach
\FormulaRefAuto{FiniteZeroJoinSemilatticeIsBounded}[thm]
automatisch beschränkt. Besitzt \(A\) ein von \(0\) verschiedenes Element,
dann gilt
\[
  \FranklSemi{A,\JoinOp,0}.
\]
Äquivalent existieren ein \(\JoinOp\)-irreduzibles \(j\) und eine Injektion
\[
  \PrincipalFilter{A,\InducedLeq}{j}
  \longrightarrow
  A\setminus\PrincipalFilter{A,\InducedLeq}{j}.
\]
\end{minipage}}
\end{center}
\end{samepage}

Diese Aussage wird bewusst weder als Axiom noch als Theorem registriert. Sie
ist mit Stand vom 23.~Juli~2026 weiterhin offen. Bewiesen werden im Folgenden
ihre exakte Äquivalenz zur zuvor entwickelten Mengenform und einige Spezialfälle.

\begin{remark}
Das Nullelement ist wesentlich. Betrachte den dreielementigen Supremums-Halbverband
\[
  A=\{a,b,1\},\qquad a\JoinOp{}b=1.
\]
Hier sind \(a\) und \(b\) die einzigen \(\JoinOp\)-irreduziblen Elemente. Beide
Hauptfilter haben zwei Elemente, ihre
Komplemente aber nur eines. Es kann daher keine Injektion vom Hauptfilter in
sein Komplement geben. Nach Adjunktion eines Nullelements entsteht das
viergliedrige boolesche Gitter; dort haben die entsprechenden Hauptfilter
genau zwei Elemente.
\end{remark}

\section{Zwei vollständig bewiesene Spezialfälle}

Zunächst isolieren wir eine elementare Abbildung, die auch außerhalb der
Frankl-Frage nützlich ist.


\paragraph{Beweisskizze.}
Ist \(t\) ein größtes irreduzibles Element, so ist sein Hauptfilter die
Einermenge \(\{t\}\). Da \(t\neq0\), liegt \(0\) im Komplement dieses
Hauptfilters; die eben konstruierte konstante Injektion liefert den
Frankl-Zeugen. In einer endlichen nichttrivialen Kette ist das größte
Nichtnullelement irreduzibel, sodass der zweite Spezialfall auf den ersten
zurückgeführt wird.

\FormulaThmDeltaK[\(\JoinOp\)-irreduzibles größtes Element erfüllt Frankl]{%
  \begin{aligned}[t]
    &\ZeroJoinSemi{A,\JoinOp,0}\dsep \JoinIrr{A,\JoinOp,0}{t}\dsep
      \forall x\in A\,x\InducedLeq t\vdash{}\\
    &\FranklSemi{A,\JoinOp,0}
  \end{aligned}
}{JoinIrreducibleTopFrankl}{%
  \DeltaRow{Mengen}{A\dsep 0\dsep t\dsep x}
  \DeltaRow{Funktionen}{\JoinOp\dsep F}
}
\begin{tabproofsplitwide}
  \proofpartwideindR[Der Hauptfilter ist eine Einermenge]{%
    \begin{aligned}[t]
      &\ZeroJoinSemi{A,\JoinOp,0}\dsep \JoinIrr{A,\JoinOp,0}{t}\dsep
        \forall x\in A\,x\InducedLeq t\vdash{}\\[-2pt]
      &\PrincipalFilter{A,\InducedLeq}{t}=\{t\}
    \end{aligned}
  }{\ZeroJoinSemi{A,\JoinOp,0}\dsep \JoinIrr{A,\JoinOp,0}{t}\dsep
    \forall x\in A\,x\InducedLeq t\vdash
    \PrincipalFilter{A,\InducedLeq}{t}=\{t\}}%
  [JoinIrreducibleTopFilterSingleton]
    \proofstepwidestar[1]{\ZeroJoinSemi{A,\JoinOp,0}}{\rA}
    \proofstepwidestar[2]{\JoinIrr{A,\JoinOp,0}{t}}{\rA}
    \proofstepwidestar[3]{\forall x\in A\,x\InducedLeq t}{\rA}
    \proofstepwidestar[2]{t\in A}{\FormulaRefAuto{JoinIrreducible}[def]{2}}
    \proofstepwidestar[1]{\Semi{A,\JoinOp}}{%
      \FormulaRefAuto{ZeroJoinBaseSemilatticeAxiom}[ax]{1}}
    \proofstepwidestar[4,5]{t\InducedLeq t}{%
      \FormulaRefAuto{JoinOrderReflexive}[thm]{5,4}}
    \proofstepwidestar[4,6]{t\in\PrincipalFilter{A,\InducedLeq}{t}}{%
      \FormulaRefAuto{PrincipalFilter}[def]{4,6}}
    \proofstepwidestar[8]{x\in\PrincipalFilter{A,\InducedLeq}{t}}{\rA}
    \proofstepwidestar[8]{x\in A\land t\InducedLeq x}{%
      \FormulaRefAuto{PrincipalFilter}[def]{8}}
    \proofstepwidestar[8]{x\in A}{\rAEa{9}}
    \proofstepwidestar[8]{t\InducedLeq x}{\rAEb{9}}
    \proofstepwidestar[3,8]{x\InducedLeq t}{\rRE{10,\rUE{3}}}
    \proofstepwidestar[1,3,8]{x=t}{%
      \FormulaRefAuto{JoinOrderAntisymmetric}[thm]{5,12,11}}
    \proofstepwidestar[1,3,8]{x\in\{t\}}{%
      \FormulaRefAuto{x\in\{a\}\eqvdash x=a}{13}}
    \proofstepwidestar[1,2,3]{%
      \PrincipalFilter{A,\InducedLeq}{t}\subseteq\{t\}}{%
      \FormulaRefAuto{A\subseteq B \coloneq \forall x\,(x\in A\rightarrow x\in B)}{%
        \rUI{\rRI{8,14}}}}
    \proofstepwidestar[1,2,3]{%
      \{t\}\subseteq\PrincipalFilter{A,\InducedLeq}{t}}{%
      \FormulaRefAuto{a\in A\vdash\{a\}\subseteq A}{7}}
    \proofstepwidestar[1,2,3]{%
      \PrincipalFilter{A,\InducedLeq}{t}=\{t\}}{%
      \FormulaRefAuto{A \subseteq B \land B \subseteq A \eqvdash A = B}{15,16}}
  \closeproofpartwideindR

  \proofpartwideindR[Injektion in das Komplement]{%
    \begin{aligned}[t]
      &\ZeroJoinSemi{A,\JoinOp,0}\dsep \JoinIrr{A,\JoinOp,0}{t}\dsep
        \forall x\in A\,x\InducedLeq t\vdash{}\\[-2pt]
      &\AtMostHalf{\PrincipalFilter{A,\InducedLeq}{t}}{A}
    \end{aligned}
  }{\ZeroJoinSemi{A,\JoinOp,0}\dsep \JoinIrr{A,\JoinOp,0}{t}\dsep
    \forall x\in A\,x\InducedLeq t\vdash
    \AtMostHalf{\PrincipalFilter{A,\InducedLeq}{t}}{A}}%
  [JoinIrreducibleTopAtMostHalf]
    \proofstepwidestar[1]{\ZeroJoinSemi{A,\JoinOp,0}}{\rA}
    \proofstepwidestar[2]{\JoinIrr{A,\JoinOp,0}{t}}{\rA}
    \proofstepwidestar[3]{\forall x\in A\,x\InducedLeq t}{\rA}
    \proofstepwidestar[1]{0\in A}{\FormulaRefAuto{ZeroJoinCarrierAxiom}[ax]{1}}
    \proofstepwidestar[2]{t\neq0}{\FormulaRefAuto{JoinIrreducible}[def]{2}}
    \proofstepwidestar[1,2]{0\in A\setminus\{t\}}{%
      \FormulaRefAuto{x\in A\setminus B\eqvdash x\in A\land x\notin B}{%
        \rAI{4,\FormulaRefAuto{x\notin\{a\}\eqvdash x\neq a}{%
          \FormulaRefAuto{a=b\vdash b=a}{5}}}}}
    \proofstepwidestar[1,2]{\exists F\colon\{t\}\inj A\setminus\{t\}}{%
      \FormulaRefAuto{SingletonInjectionIntoNonempty}[thm]{6}}
    \proofstepwidestar[1,2,3]{%
      \PrincipalFilter{A,\InducedLeq}{t}=\{t\}}{%
      \FormulaRefAuto{JoinIrreducibleTopFilterSingleton}[thm]{1,2,3}}
    \proofstepwidestar[]{\{t\}\subseteq A}{%
      \FormulaRefAuto{a\in A\vdash\{a\}\subseteq A}{%
        \FormulaRefAuto{JoinIrreducible}[def]{2}}}
    \proofstepwidestar[1,2,3]{%
      \AtMostHalf{\PrincipalFilter{A,\InducedLeq}{t}}{A}}{%
      \FormulaRefAuto{AtMostHalf}[def]{\rIE{8,9},\rIE{8,7}}}
  \closeproofpartwideindR

  \proofpartwide{Schluss}
    \proofstepwidestar[1]{\ZeroJoinSemi{A,\JoinOp,0}}{\rA}
    \proofstepwidestar[2]{\JoinIrr{A,\JoinOp,0}{t}}{\rA}
    \proofstepwidestar[3]{\forall x\in A\,x\InducedLeq t}{\rA}
    \proofstepwidestar[1,2,3]{%
      \AtMostHalf{\PrincipalFilter{A,\InducedLeq}{t}}{A}}{%
      \FormulaRefAuto{JoinIrreducibleTopAtMostHalf}[thm]{1,2,3}}
    \proofstepwidestar[1,2,3]{\FranklSemi{A,\JoinOp,0}}{%
      \FormulaRefAuto{FranklSemilatticeProperty}[def]{%
        \rEI{\rAI{2,4}}}}
  \closeproofpartwide
\end{tabproofsplitwide}

\FormulaThmDeltaK[Endliche Ketten erfüllen die Halbverbandsform]{%
  \begin{aligned}[t]
    &\ZeroJoinSemi{A,\JoinOp,0}\dsep \Finite A\dsep
      \TotOrd{A,\InducedLeq}\dsep
      t\in A\dsep{}\\
    &\forall x\in A\,x\InducedLeq t\dsep
      t\neq0\vdash{}\\
    &\FranklSemi{A,\JoinOp,0}
  \end{aligned}
}{FranklJoinChains}{%
  \DeltaRow{Mengen}{A\dsep 0\dsep t\dsep x\dsep y}
  \DeltaRow{Funktionen}{\JoinOp}
}
\begin{tabproofsplitwide}
  \proofpartwideindR[Das größte Kettenelement ist irreduzibel]{%
    \begin{aligned}[t]
      &\ZeroJoinSemi{A,\JoinOp,0}\dsep
        \TotOrd{A,\InducedLeq}\dsep t\in A\dsep{}\\[-2pt]
      &\forall x\in A\,x\InducedLeq t\dsep t\neq0
        \vdash \JoinIrr{A,\JoinOp,0}{t}
    \end{aligned}
  }{\ZeroJoinSemi{A,\JoinOp,0}\dsep
    \TotOrd{A,\InducedLeq}\dsep t\in A\dsep
    \forall x\in A\,x\InducedLeq t\dsep t\neq0
    \vdash \JoinIrr{A,\JoinOp,0}{t}}[FranklJoinChainTopIrreducible]
    \proofstepwidestar[1]{\ZeroJoinSemi{A,\JoinOp,0}}{\rA}
    \proofstepwidestar[2]{\TotOrd{A,\InducedLeq}}{\rA}
    \proofstepwidestar[3]{t\in A}{\rA}
    \proofstepwidestar[4]{\forall x\in A\,x\InducedLeq t}{\rA}
    \proofstepwidestar[5]{t\neq0}{\rA}
    \proofstepwidestar[1]{\Semi{A,\JoinOp}}{%
      \FormulaRefAuto{ZeroJoinBaseSemilatticeAxiom}[ax]{1}}
    \proofstepwidestar[7]{x\in A}{\rA}
    \proofstepwidestar[8]{y\in A}{\rA}
    \proofstepwidestar[9]{x\JoinOp{}y=t}{\rA}
    \proofstepwidestar[2,7,8]{%
      x\InducedLeq y\lor y\InducedLeq x}{%
      \FormulaRefAuto{x\in A \dsep y\in A \vdash
        x\leq y \lor y\leq x}[ax]{2,7,8}}
    \proofstepwidestar[11]{x\InducedLeq y}{\rA}
    \proofstepwidestar[1,7,8,11]{x\JoinOp{}y=y}{%
      \FormulaRefAuto{JoinOrderEvaluation}[thm]{6,11}}
    \proofstepwidestar[7,8,9,11]{y=t}{%
      \FormulaRefAuto{a=b,\,a=c\vdash b=c}{12,9}}
    \proofstepwidestar[7,8,9,11]{x=t\lor y=t}{\rOIb{13}}
    \proofstepwidestar[15]{y\InducedLeq x}{\rA}
    \proofstepwidestar[1,7,8,15]{y\JoinOp{}x=x}{%
      \FormulaRefAuto{JoinOrderEvaluation}[thm]{6,15}}
    \proofstepwidestar[1,7,8]{y\JoinOp{}x=x\JoinOp{}y}{%
      \FormulaRefAuto{SemilatticeCommutativityAxiom}[ax]{6,8,7}}
    \proofstepwidestar[1,7,8,15]{x\JoinOp{}y=x}{%
      \rIE{17,16}}
    \proofstepwidestar[7,8,9,15]{x=t}{%
      \FormulaRefAuto{a=b,\,a=c\vdash b=c}{18,9}}
    \proofstepwidestar[7,8,9,15]{x=t\lor y=t}{\rOIa{19}}
    \proofstepwidestar[2,7,8,9]{x=t\lor y=t}{%
      \rOE{10,11,14,15,20}}
    \proofstepwidestar[1,2,3,4,5]{\JoinIrr{A,\JoinOp,0}{t}}{%
      \FormulaRefAuto{JoinIrreducible}[def]{3,5,7,8,9,21}}
  \closeproofpartwideindR

  \proofpartwide{Schluss}
    \proofstepwidestar[1]{\ZeroJoinSemi{A,\JoinOp,0}}{\rA}
    \proofstepwidestar[2]{\Finite A}{\rA}
    \proofstepwidestar[3]{\TotOrd{A,\InducedLeq}}{\rA}
    \proofstepwidestar[4]{t\in A}{\rA}
    \proofstepwidestar[5]{\forall x\in A\,x\InducedLeq t}{\rA}
    \proofstepwidestar[6]{t\neq0}{\rA}
    \proofstepwidestar[1,3,4,5,6]{\JoinIrr{A,\JoinOp,0}{t}}{%
      \FormulaRefAuto{FranklJoinChainTopIrreducible}[thm]{1,3,4,5,6}}
    \proofstepwidestar[1,3,4,5,6]{\FranklSemi{A,\JoinOp,0}}{%
      \FormulaRefAuto{JoinIrreducibleTopFrankl}[thm]{1,7,5}}
  \closeproofpartwide
\end{tabproofsplitwide}

Der Endlichkeitsparameter wird im letzten Beweis nur benötigt, um die
Kettenaussage als Spezialfall der endlichen Frankl-Vermutung einzuordnen. Ist
ein größtes Element bereits gegeben, funktioniert die Injektion unabhängig
von der Endlichkeit.


\section{Vom Halbverband zur Vereinigungsfamilie}

\paragraph{Beweisskizze.}
Zu \(x\in A\) betrachten wir zunächst das Irreduziblenprofil
\(\IrrBelow{A,\JoinOp,0}{x}\) und anschließend sein Komplement
\(\CanonicalProfileMap{A,\JoinOp,0}(x)\). Irreduzible Elemente trennen
verschiedene Elemente; daher ist \(\CanonicalProfileMap{A,\JoinOp,0}\)
bijektiv. Paarinfima werden auf Durchschnitte der
Irreduziblenprofile und damit, nach de~Morgan, auf Vereinigungen der
Komplementprofile abgebildet. So entsteht eine endliche Frankl-Familie
\(\mathcal U_{A,\JoinOp,0}\). Ein dort halbhäufiges Profilelement wird entlang
der Bijektion zurücktransportiert und liefert den gesuchten kleinen
Hauptfilter:
\[
  \begin{aligned}
    A&\xrightarrow[\text{bijektiv}]{\CanonicalProfileMap{A,\JoinOp,0}}
      \mathcal U_{A,\JoinOp,0},\\
    j\notin\CanonicalProfileMap{A,\JoinOp,0}(x)
      &\quad\Longleftrightarrow\quad j\InducedLeq x.
  \end{aligned}
\]

\FormulaDefDeltaK[\(\JoinOp\)-Irreduzible unter einem Element]{%
  \IrrBelow{A,\JoinOp,0}{x}
  \coloneqq
  \{j\in\JoinIrrSet{A,\JoinOp,0}\mid j\InducedLeq x\}
}{IrreduciblesBelow}{%
  \DeltaRow{Mengen}{A\dsep j\dsep x\dsep 0}
  \DeltaRow{Funktionen}{\JoinOp}
  \DeltaRow{\textbf{Neues Symbol}}{}
  \DeltaPrem{Irreduziblenprofil}{\IrrBelow{A,\JoinOp,0}{x}}
}

\FormulaDefDeltaK[Kanonische Profilfamilien]{%
  \begin{aligned}[t]
    \CanonicalMeetFamily{A,\JoinOp,0}
      &\coloneqq\{\IrrBelow{A,\JoinOp,0}{x}\mid x\in A\},\\
    \mathcal U_{A,\JoinOp,0}
      &\coloneqq
      \{\JoinIrrSet{A,\JoinOp,0}\setminus\IrrBelow{A,\JoinOp,0}{x}\mid x\in A\}
  \end{aligned}
}{CanonicalProfileFamilies}{%
  \DeltaRow{Mengen}{A\dsep x\dsep 0}
  \DeltaRow{Funktionen}{\JoinOp}
  \DeltaRow{\textbf{Neue Symbole}}{}
  \DeltaPrem{Schnittprofilfamilie}{\CanonicalMeetFamily{A,\JoinOp,0}}
  \DeltaPrem{Komplementprofilfamilie}{\mathcal U_{A,\JoinOp,0}}
}

Die erste Profilfamilie ist schnittabgeschlossen. Ist \(m\) das Paarinfimum
von \(x\) und \(y\), so folgt unmittelbar aus dessen universeller Eigenschaft
\[
  \IrrBelow{A,\JoinOp,0}{m}
  =
  \IrrBelow{A,\JoinOp,0}{x}\cap\IrrBelow{A,\JoinOp,0}{y}.
\]
Die Komplementprofilfamilie ist daher nach de~Morgan
vereinigungsabgeschlossen. Entscheidend ist außerdem, dass die Profile alle
Elemente unterscheiden.

\FormulaThmDeltaK[Charakterisierung der Komplementprofilfamilie]{%
  \begin{aligned}[t]
    Y\in\mathcal U_{A,\JoinOp,0}
    \quad\Longleftrightarrow\quad
    \exists x\in A\,\bigl(
      Y={}&\JoinIrrSet{A,\JoinOp,0}\setminus
      \IrrBelow{A,\JoinOp,0}{x}\bigr)
  \end{aligned}
}{CanonicalProfileFamilyMembership}{%
  \DeltaRow{Mengen}{A\dsep Y\dsep x\dsep 0}
  \DeltaRow{Funktionen}{\JoinOp}
}
\begin{tabproofwide}
  \proofstepwide{Y\in\mathcal U_{A,\JoinOp,0}}{\leftrightarrow}{%
    \exists x\in A\,\bigl(
      Y=\JoinIrrSet{A,\JoinOp,0}\setminus
      \IrrBelow{A,\JoinOp,0}{x}\bigr)}{%
    \FormulaRefAuto{CanonicalProfileFamilies}[def]}
\end{tabproofwide}

\FormulaDefDeltaK[Kanonische Profilabbildung]{%
  \ZeroJoinSemi{A,\JoinOp,0}\vdash
  \begin{aligned}[t]
    &\CanonicalProfileMap{A,\JoinOp,0}
      \colon A\to\mathcal U_{A,\JoinOp,0},\\[-2pt]
    &\forall x\in A\,
      \bigl(
        \CanonicalProfileMap{A,\JoinOp,0}(x)\coloneqq
        \JoinIrrSet{A,\JoinOp,0}\setminus\IrrBelow{A,\JoinOp,0}{x}
      \bigr)
  \end{aligned}
}{CanonicalProfileMapDefinition}{%
  \DeltaRow{Mengen}{A\dsep x\dsep 0}
  \DeltaRow{Funktionen}{\JoinOp\dsep\CanonicalProfileMap{A,\JoinOp,0}}
  \DeltaRow{\textbf{Neues Symbol}}{}
  \DeltaPrem{Profilabbildung}{\CanonicalProfileMap{A,\JoinOp,0}}
}
\begin{tabproofwide}
  \proofstepwidestar[1]{\ZeroJoinSemi{A,\JoinOp,0}}{\rA}
  \proofstepwidestar[2]{x\in A}{\rA}
  \proofstepwidestar[2]{%
    \exists u\in A\,\bigl(
      \JoinIrrSet{A,\JoinOp,0}\setminus\IrrBelow{A,\JoinOp,0}{x}
      =\JoinIrrSet{A,\JoinOp,0}\setminus\IrrBelow{A,\JoinOp,0}{u}\bigr)}{%
    \rEI{\rAI{2,\rII}}}
  \proofstepwidestar[1,2]{%
    \JoinIrrSet{A,\JoinOp,0}\setminus\IrrBelow{A,\JoinOp,0}{x}
      \in\mathcal U_{A,\JoinOp,0}}{%
    \FormulaRefAuto{CanonicalProfileFamilyMembership}[thm]{3}}
  \proofstepwidestar[1]{%
    \forall x\in A\,
      \JoinIrrSet{A,\JoinOp,0}\setminus\IrrBelow{A,\JoinOp,0}{x}
        \in\mathcal U_{A,\JoinOp,0}}{%
    \rUI{\rRI{2,4}}}
\end{tabproofwide}

\FormulaThmDeltaK[Die kanonische Profilabbildung ist eine Funktion]{%
  \ZeroJoinSemi{A,\JoinOp,0}\vdash
  \CanonicalProfileMap{A,\JoinOp,0}
    \colon A\to\mathcal U_{A,\JoinOp,0}
}{CanonicalProfileFunction}{%
  \DeltaRow{Mengen}{A\dsep x\dsep 0}
  \DeltaRow{Funktionen}{\JoinOp\dsep\CanonicalProfileMap{A,\JoinOp,0}}
}
\begin{tabproofwide}
  \proofstepwidestar[1]{\ZeroJoinSemi{A,\JoinOp,0}}{\rA}
  \proofstepwidestar[1]{%
    \CanonicalProfileMap{A,\JoinOp,0}
      \colon A\to\mathcal U_{A,\JoinOp,0}}{%
    \FormulaRefAuto{CanonicalProfileMapDefinition}[def]{1}}
\end{tabproofwide}

\FormulaThmDeltaK[Auswertung der kanonischen Profilabbildung]{%
  \begin{aligned}[t]
    &\ZeroJoinSemi{A,\JoinOp,0}\dsep x\in A\vdash{}\\
    &\CanonicalProfileMap{A,\JoinOp,0}(x)=
      \JoinIrrSet{A,\JoinOp,0}\setminus\IrrBelow{A,\JoinOp,0}{x}
  \end{aligned}
}{CanonicalProfileEvaluation}{%
  \DeltaRow{Mengen}{A\dsep x\dsep 0}
  \DeltaRow{Funktionen}{\JoinOp\dsep\CanonicalProfileMap{A,\JoinOp,0}}
}
\begin{tabproofwide}
  \proofstepwidestar[1]{\ZeroJoinSemi{A,\JoinOp,0}}{\rA}
  \proofstepwidestar[2]{x\in A}{\rA}
  \proofstepwidestar[1]{%
    \forall u\in A\,
      \CanonicalProfileMap{A,\JoinOp,0}(u)=
        \JoinIrrSet{A,\JoinOp,0}\setminus\IrrBelow{A,\JoinOp,0}{u}}{%
    \FormulaRefAuto{CanonicalProfileMapDefinition}[def]{1}}
  \proofstepwidestar[1,2]{%
    \CanonicalProfileMap{A,\JoinOp,0}(x)=
      \JoinIrrSet{A,\JoinOp,0}\setminus\IrrBelow{A,\JoinOp,0}{x}}{%
    \rRE{\rUE{3},2}}
\end{tabproofwide}

\FormulaThmDeltaK[Die kanonische Profilabbildung ist surjektiv]{%
  \ZeroJoinSemi{A,\JoinOp,0}\vdash
  \CanonicalProfileMap{A,\JoinOp,0}
    \colon A\sur\mathcal U_{A,\JoinOp,0}
}{CanonicalProfileSurjective}{%
  \DeltaRow{Mengen}{A\dsep Y\dsep x\dsep 0}
  \DeltaRow{Funktionen}{\JoinOp\dsep\CanonicalProfileMap{A,\JoinOp,0}}
}
\begin{tabproofwide}
  \proofstepwidestar[1]{\ZeroJoinSemi{A,\JoinOp,0}}{\rA}
  \proofstepwidestar[1]{%
    \CanonicalProfileMap{A,\JoinOp,0}
      \colon A\to\mathcal U_{A,\JoinOp,0}}{%
    \FormulaRefAuto{CanonicalProfileFunction}[thm]{1}}
  \proofstepwidestar[3]{Y\in\mathcal U_{A,\JoinOp,0}}{\rA}
  \proofstepwidestar[3]{%
    \exists x\in A\,\bigl(
      Y=\JoinIrrSet{A,\JoinOp,0}\setminus
        \IrrBelow{A,\JoinOp,0}{x}\bigr)}{%
    \FormulaRefAuto{CanonicalProfileFamilyMembership}[thm]{3}}
  \proofstepwidestar[5]{%
    x\in A\land
    Y=\JoinIrrSet{A,\JoinOp,0}\setminus
      \IrrBelow{A,\JoinOp,0}{x}}{\rA}
  \proofstepwidestar[5]{x\in A}{\rAEa{5}}
  \proofstepwidestar[1,5]{%
    \CanonicalProfileMap{A,\JoinOp,0}(x)=
      \JoinIrrSet{A,\JoinOp,0}\setminus\IrrBelow{A,\JoinOp,0}{x}}{%
    \FormulaRefAuto{CanonicalProfileEvaluation}[thm]{1,6}}
  \proofstepwidestar[1,5]{%
    \CanonicalProfileMap{A,\JoinOp,0}(x)=Y}{%
    \FormulaRefAuto{a=b,\,c=b\vdash a=c}{7,\rAEb{5}}}
  \proofstepwidestar[1,3]{%
    \exists x\in A\,\bigl(
      \CanonicalProfileMap{A,\JoinOp,0}(x)=Y\bigr)}{%
    \rEE{4,5,\rEI{\rAI{6,8}}}}
  \proofstepwidestar[1]{%
    \forall Y\in\mathcal U_{A,\JoinOp,0}\,
      \exists x\in A\,\bigl(
        \CanonicalProfileMap{A,\JoinOp,0}(x)=Y\bigr)}{%
    \rUI{\rRI{3,9}}}
  \proofstepwidestar[1]{%
    \CanonicalProfileMap{A,\JoinOp,0}
      \colon A\sur\mathcal U_{A,\JoinOp,0}}{%
    \FormulaRefAuto{Surjektive Funktion}[def]{2,10}}
\end{tabproofwide}

\FormulaThmDeltaK[Irreduziblenprofile liegen im Irreduziblenträger]{%
  \ZeroJoinSemi{A,\JoinOp,0}\dsep x\in A\vdash
  \IrrBelow{A,\JoinOp,0}{x}\subseteq\JoinIrrSet{A,\JoinOp,0}
}{IrreducibleProfileSubset}{%
  \DeltaRow{Mengen}{A\dsep j\dsep x\dsep 0}
  \DeltaRow{Funktionen}{\JoinOp}
}
\begin{tabproofwide}
  \proofstepwidestar[1]{\ZeroJoinSemi{A,\JoinOp,0}}{\rA}
  \proofstepwidestar[2]{x\in A}{\rA}
  \proofstepwidestar[3]{j\in\IrrBelow{A,\JoinOp,0}{x}}{\rA}
  \proofstepwidestar[3]{%
    j\in\JoinIrrSet{A,\JoinOp,0}\land
      j\InducedLeq x}{%
    \FormulaRefAuto{IrreduciblesBelow}[def]{3}}
  \proofstepwidestar[3]{j\in\JoinIrrSet{A,\JoinOp,0}}{\rAEa{4}}
  \proofstepwidestar[]{%
    \IrrBelow{A,\JoinOp,0}{x}\subseteq\JoinIrrSet{A,\JoinOp,0}}{%
    \FormulaRefAuto{A \subseteq B \coloneq
      \forall x\,(x\in A\rightarrow x\in B)}{\rUI{\rRI{3,5}}}}
\end{tabproofwide}

\FormulaThmDeltaK[Profilinklusion reflektiert die Halbverbandsordnung]{%
  \begin{aligned}[t]
    &\Finite A\dsep\ZeroJoinSemi{A,\JoinOp,0}\dsep x,y\in A\dsep{}\\
    &\IrrBelow{A,\JoinOp,0}{x}\subseteq\IrrBelow{A,\JoinOp,0}{y}
      \vdash x\InducedLeq y
  \end{aligned}
}{IrreducibleProfileReflectsOrder}{%
  \DeltaRow{Mengen}{A\dsep j\dsep x\dsep y\dsep 0}
  \DeltaRow{Funktionen}{\JoinOp}
}
\begin{tabproofwide}
  \proofstepwidestar[1]{\Finite A}{\rA}
  \proofstepwidestar[2]{\ZeroJoinSemi{A,\JoinOp,0}}{\rA}
  \proofstepwidestar[3]{x\in A}{\rA}
  \proofstepwidestar[4]{y\in A}{\rA}
  \proofstepwidestar[5]{%
    \IrrBelow{A,\JoinOp,0}{x}\subseteq\IrrBelow{A,\JoinOp,0}{y}}{\rA}
  \proofstepwidestar[6]{x\not\InducedLeq y}{\rA}
  \proofstepwidestar[1,2,3,4,6]{%
    \exists j\in\JoinIrrSet{A,\JoinOp,0}\,\bigl(
      j\InducedLeq x\land
      j\not\InducedLeq y\bigr)}{%
    \FormulaRefAuto{JoinIrreduciblesSeparate}[thm]{1,2,3,4,6}}
  \proofstepwidestar[8]{%
    j\in\JoinIrrSet{A,\JoinOp,0}\land\bigl(
      j\InducedLeq x\land
      j\not\InducedLeq y\bigr)}{\rA}
  \proofstepwidestar[8]{j\in\JoinIrrSet{A,\JoinOp,0}}{\rAEa{8}}
  \proofstepwidestar[8]{j\InducedLeq x}{\rAEa{\rAEb{8}}}
  \proofstepwidestar[8]{j\not\InducedLeq y}{\rAEb{\rAEb{8}}}
  \proofstepwidestar[8]{j\in\IrrBelow{A,\JoinOp,0}{x}}{%
    \FormulaRefAuto{IrreduciblesBelow}[def]{\rAI{9,10}}}
  \proofstepwidestar[5,8]{j\in\IrrBelow{A,\JoinOp,0}{y}}{%
    \FormulaRefAuto{A\subseteq B,\,x\in A\vdash x\in B}{5,12}}
  \proofstepwidestar[5,8]{j\InducedLeq y}{%
    \rAEb{\FormulaRefAuto{IrreduciblesBelow}[def]{13}}}
  \proofstepwidestar[5,8]{\bot}{\rBI{11,14}}
  \proofstepwidestar[1,2,3,4,5,6]{\bot}{\rEE{7,8,15}}
  \proofstepwidestar[1,2,3,4,5]{x\InducedLeq y}{%
    \rCE{6,16}}
\end{tabproofwide}

\FormulaThmDeltaK[Die kanonische Profilabbildung ist injektiv]{%
  \Finite A\dsep\ZeroJoinSemi{A,\JoinOp,0}\vdash
  \CanonicalProfileMap{A,\JoinOp,0}
    \colon A\inj\mathcal U_{A,\JoinOp,0}
}{CanonicalProfileInjective}{%
  \DeltaRow{Mengen}{A\dsep x\dsep y\dsep 0}
  \DeltaRow{Funktionen}{\JoinOp\dsep\CanonicalProfileMap{A,\JoinOp,0}}
}
Für den Beweis verwenden wir die Abkürzungen
\[
  \begin{aligned}
    \Phi&\coloneqq\CanonicalProfileMap{A,\JoinOp,0},&
    \mathcal U&\coloneqq\mathcal U_{A,\JoinOp,0},\\
    J&\coloneqq\JoinIrrSet{A,\JoinOp,0},&
    I_z&\coloneqq\IrrBelow{A,\JoinOp,0}{z}.
  \end{aligned}
\]
\begin{tabproofwide}
  \proofstepwidestar[1]{\Finite A}{\rA}
  \proofstepwidestar[2]{\ZeroJoinSemi{A,\JoinOp,0}}{\rA}
  \proofstepwidestar[2]{\Phi\colon A\to\mathcal U}{%
    \FormulaRefAuto{CanonicalProfileFunction}[thm]{2}}
  \proofstepwidestar[4]{x\in A}{\rA}
  \proofstepwidestar[5]{y\in A}{\rA}
  \proofstepwidestar[6]{\Phi(x)=\Phi(y)}{\rA}
  \proofstepwidestar[2,4]{\Phi(x)=J\setminus I_x}{%
    \FormulaRefAuto{CanonicalProfileEvaluation}[thm]{2,4}}
  \proofstepwidestar[2,5]{\Phi(y)=J\setminus I_y}{%
    \FormulaRefAuto{CanonicalProfileEvaluation}[thm]{2,5}}
  \proofstepwidestar[2,4,5,6]{J\setminus I_x=\Phi(y)}{%
    \FormulaRefAuto{a=b,\,a=c\vdash b=c}{7,6}}
  \proofstepwidestar[2,4,5,6]{J\setminus I_x=J\setminus I_y}{%
    \FormulaRefAuto{a=b,\,b=c\vdash a=c}{9,8}}
  \proofstepwidestar[2,4]{I_x\subseteq J}{%
    \FormulaRefAuto{IrreducibleProfileSubset}[thm]{2,4}}
  \proofstepwidestar[2,5]{I_y\subseteq J}{%
    \FormulaRefAuto{IrreducibleProfileSubset}[thm]{2,5}}
  \proofstepwidestar[2,4,5,6]{I_x=I_y}{%
    \FormulaRefAuto{RelCompEqualityCancel}[thm]{12,11,10}}
  \proofstepwidestar[2,4,5,6]{I_x\subseteq I_y}{%
    \FormulaRefAuto{A=B\vdash A\subseteq B}{13}}
  \proofstepwidestar[1,2,4,5,6]{x\InducedLeq y}{%
    \FormulaRefAuto{IrreducibleProfileReflectsOrder}[thm]{1,2,4,5,14}}
  \proofstepwidestar[2,4,5,6]{I_y\subseteq I_x}{%
    \FormulaRefAuto{A=B\vdash B\subseteq A}{13}}
  \proofstepwidestar[1,2,4,5,6]{y\InducedLeq x}{%
    \FormulaRefAuto{IrreducibleProfileReflectsOrder}[thm]{1,2,5,4,16}}
  \proofstepwidestar[1,2,4,5,6]{x=y}{%
    \ProofRefStack{%
      \FormulaRefAuto{JoinOrderAntisymmetric}[thm]\\
      \bigl(\FormulaRefAuto{ZeroJoinBaseSemilatticeAxiom}[ax]{2},
        15,17\bigr)}}
  \proofstepwidestar[1,2]{\Phi\colon A\inj\mathcal U}{%
    \ProofRefStack{%
      \FormulaRefAuto{Injektive Funktion}[def]\\
      \bigl(3,\rUI{\rRI{4,\rRI{5,\rRI{6,18}}}}\bigr)}}
\end{tabproofwide}

\FormulaThmDeltaK[Die kanonische Profilabbildung ist bijektiv]{%
  \Finite A\dsep\ZeroJoinSemi{A,\JoinOp,0}\vdash
  \CanonicalProfileMap{A,\JoinOp,0}
    \colon A\bij\mathcal U_{A,\JoinOp,0}
}{CanonicalProfileBijection}{%
  \DeltaRow{Mengen}{A\dsep 0}
  \DeltaRow{Funktionen}{\JoinOp\dsep\CanonicalProfileMap{A,\JoinOp,0}}
}
\begin{tabproofwide}
  \proofstepwidestar[1]{\Finite A}{\rA}
  \proofstepwidestar[2]{\ZeroJoinSemi{A,\JoinOp,0}}{\rA}
  \proofstepwidestar[1,2]{%
    \CanonicalProfileMap{A,\JoinOp,0}
      \colon A\inj\mathcal U_{A,\JoinOp,0}}{%
    \FormulaRefAuto{CanonicalProfileInjective}[thm]{1,2}}
  \proofstepwidestar[2]{%
    \CanonicalProfileMap{A,\JoinOp,0}
      \colon A\sur\mathcal U_{A,\JoinOp,0}}{%
    \FormulaRefAuto{CanonicalProfileSurjective}[thm]{2}}
  \proofstepwidestar[1,2]{%
    \CanonicalProfileMap{A,\JoinOp,0}
      \colon A\bij\mathcal U_{A,\JoinOp,0}}{%
    \FormulaRefAuto{F\colon A\inj B\dsep F\colon A\sur B
      \vdash F\colon A\bij B}{3,4}}
\end{tabproofwide}

\FormulaThmDeltaK[Repräsentant eines kanonischen Profils]{%
  \begin{aligned}[t]
    &\ZeroJoinSemi{A,\JoinOp,0}\dsep
      X\in\mathcal U_{A,\JoinOp,0}\vdash{}\\
    &\exists x\in A\,\bigl(
      X=\CanonicalProfileMap{A,\JoinOp,0}(x)\bigr)
  \end{aligned}
}{CanonicalProfileRepresentative}{%
  \DeltaRow{Mengen}{A\dsep X\dsep x\dsep 0}
  \DeltaRow{Funktionen}{\JoinOp\dsep\CanonicalProfileMap{A,\JoinOp,0}}
}
\begin{tabproofwide}
  \proofstepwidestar[1]{\ZeroJoinSemi{A,\JoinOp,0}}{\rA}
  \proofstepwidestar[2]{X\in\mathcal U_{A,\JoinOp,0}}{\rA}
  \proofstepwidestar[1]{%
    \CanonicalProfileMap{A,\JoinOp,0}
      \colon A\sur\mathcal U_{A,\JoinOp,0}}{%
    \FormulaRefAuto{CanonicalProfileSurjective}[thm]{1}}
  \proofstepwidestar[1,2]{%
    \exists x\in A\,\bigl(
      \CanonicalProfileMap{A,\JoinOp,0}(x)=X\bigr)}{%
    \FormulaRefAuto{Surjektivität}{3,2}}
  \proofstepwidestar[5]{%
    x\in A\land\CanonicalProfileMap{A,\JoinOp,0}(x)=X}{\rA}
  \proofstepwidestar[5]{x\in A}{\rAEa{5}}
  \proofstepwidestar[5]{%
    X=\CanonicalProfileMap{A,\JoinOp,0}(x)}{%
    \FormulaRefAuto{a=b\vdash b=a}{\rAEb{5}}}
  \proofstepwidestar[5]{%
    \exists x\in A\,\bigl(
      X=\CanonicalProfileMap{A,\JoinOp,0}(x)\bigr)}{%
    \rEI{\rAI{6,7}}}
  \proofstepwidestar[1,2]{%
    \exists x\in A\,\bigl(
      X=\CanonicalProfileMap{A,\JoinOp,0}(x)\bigr)}{%
    \rEE{4,5,8}}
\end{tabproofwide}

\FormulaThmDeltaK[Relatives de-Morgan-Gesetz]{%
  A\setminus(B\cap C)
  =
  (A\setminus B)\cup(A\setminus C)
}{RelativeComplementIntersectionDeMorgan}{%
  \DeltaRow{Mengen}{A\dsep B\dsep C\dsep x}
}
\begin{tabproofwide}
  \proofstepwide{x\in A\setminus(B\cap C)}{\leftrightarrow}{%
    x\in A\land x\notin B\cap C}{%
    \FormulaRefAuto{x \in A \setminus B \eqvdash
      x \in A \land x \notin B}}
  \proofstepwide{} {\leftrightarrow}{%
    x\in A\land\neg(x\in B\land x\in C)}{%
    \ProofRefStack{%
      \FormulaRefAuto{x \in A \cap B \eqvdash
        x \in A \land x \in B}\\
      \FormulaRefAuto{P \leftrightarrow Q \eqvdash
        \neg P \leftrightarrow \neg Q}\,(2)}}
  \proofstepwide{} {\leftrightarrow}{%
    x\in A\land(x\notin B\lor x\notin C)}{%
    \ProofRefStack{%
      \FormulaRefAuto{\neg(P \land Q) \eqvdash
        \neg P \lor \neg Q}\,(2)}}
  \proofstepwide{} {\leftrightarrow}{%
    \begin{aligned}[t]
      &(x\in A\land x\notin B)\lor{}\\[-2pt]
      &(x\in A\land x\notin C)
    \end{aligned}}{%
    \FormulaRefAuto{P \land (Q\lor R) \eqvdash
      (P \land Q) \lor (P\land R)}}
  \proofstepwide{} {\leftrightarrow}{%
    x\in A\setminus B\lor(x\in A\land x\notin C)}{%
    \ProofRefStack{%
      \FormulaRefAuto{x \in A \setminus B \eqvdash
        x \in A \land x \notin B}\,(1)}}
  \proofstepwide{} {\leftrightarrow}{%
    x\in A\setminus B\lor x\in A\setminus C}{%
    \ProofRefStack{%
      \FormulaRefAuto{x \in A \setminus B \eqvdash
        x \in A \land x \notin B}\,(1)}}
  \proofstepwide{} {\leftrightarrow}{%
    x\in(A\setminus B)\cup(A\setminus C)}{%
    \FormulaRefAuto{x \in A \cup B \eqvdash
      x \in A \lor x \in B}}
  \proofstepwide{A\setminus(B\cap C)}{=}{%
    (A\setminus B)\cup(A\setminus C)}{%
    \ProofRefStack{%
      \FormulaRefAuto{\forall x\, (x \in A \leftrightarrow x \in B)
        \eqvdash A = B}\\
      \bigl(\rUI{\rChain{1,2,3,4,5,6,7}}\bigr)}}
\end{tabproofwide}

\FormulaThmDeltaK[Komplementprofile eines Schnitts vereinigen sich]{%
  \begin{aligned}[t]
    &D=B\cap C\dsep X=U\setminus D\dsep
      Y=U\setminus B\dsep Z=U\setminus C\vdash{}\\
    &X=Y\cup Z
  \end{aligned}
}{RelativeComplementProfileUnionTransport}{%
  \DeltaRow{Mengen}{B\dsep C\dsep D\dsep U\dsep X\dsep Y\dsep Z}
}
\begin{tabproofwide}
  \proofstepwidestar[1]{D=B\cap C}{\rA}
  \proofstepwidestar[2]{X=U\setminus D}{\rA}
  \proofstepwidestar[3]{Y=U\setminus B}{\rA}
  \proofstepwidestar[4]{Z=U\setminus C}{\rA}
  \proofstepwidestar[]{%
    U\setminus(B\cap C)
      =(U\setminus B)\cup(U\setminus C)}{%
    \FormulaRefAuto{RelativeComplementIntersectionDeMorgan}[thm]}
  \proofstepwidestar[1]{%
    U\setminus D=U\setminus(B\cap C)}{\rIE{1,\rII}}
  \proofstepwidestar[1,2]{%
    X=U\setminus(B\cap C)}{%
    \FormulaRefAuto{a=b,\,b=c\vdash a=c}{2,6}}
  \proofstepwidestar[1,2]{%
    X=(U\setminus B)\cup(U\setminus C)}{%
    \FormulaRefAuto{a=b,\,b=c\vdash a=c}{7,5}}
  \proofstepwidestar[3,4]{%
    Y\cup Z=(U\setminus B)\cup(U\setminus C)}{%
    \FormulaRefAuto{a=b\dsep c=d \dsep P(b,d)\vdash P(a,c)}{%
      3,4,\rII}}
  \proofstepwidestar[1,2,3,4]{X=Y\cup Z}{%
    \FormulaRefAuto{a=b,\,c=b\vdash a=c}{8,9}}
\end{tabproofwide}

\FormulaThmDeltaK[Paarinfimum und Schnitt der Irreduziblenprofile]{%
  \begin{aligned}[t]
    &\ZeroJoinSemi{A,\JoinOp,0}\dsep x,y\in A\dsep
      m\in\LB_{A,\InducedLeq}(\{x,y\})\dsep{}\\
    &\forall d\in\LB_{A,\InducedLeq}(\{x,y\})\,
      d\InducedLeq m\vdash{}\\
    &\IrrBelow{A,\JoinOp,0}{m}
      =\IrrBelow{A,\JoinOp,0}{x}\cap
       \IrrBelow{A,\JoinOp,0}{y}
  \end{aligned}
}{CanonicalProfilePairInfimumIntersection}{%
  \DeltaRow{Mengen}{A\dsep d\dsep j\dsep m\dsep x\dsep y\dsep 0}
  \DeltaRow{Funktionen}{\JoinOp}
}
\begin{tabproofsplitwide}
  \proofpartwideindR[Erste Teilmengenrichtung]{%
    \begin{aligned}[t]
      &\ZeroJoinSemi{A,\JoinOp,0}\dsep x,y\in A\dsep
        m\in\LB_{A,\InducedLeq}(\{x,y\})\dsep{}\\[-2pt]
      &\forall d\in\LB_{A,\InducedLeq}(\{x,y\})\,
        d\InducedLeq m\vdash{}\\[-2pt]
      &\IrrBelow{A,\JoinOp,0}{m}\subseteq
        \IrrBelow{A,\JoinOp,0}{x}\cap
        \IrrBelow{A,\JoinOp,0}{y}
    \end{aligned}
  }{%
    \ZeroJoinSemi{A,\JoinOp,0}\dsep x,y\in A\dsep
    m\in\LB_{A,\InducedLeq}(\{x,y\})\dsep
    \forall d\in\LB_{A,\InducedLeq}(\{x,y\})\,d\InducedLeq m\vdash
    \IrrBelow{A,\JoinOp,0}{m}\subseteq
      \IrrBelow{A,\JoinOp,0}{x}\cap\IrrBelow{A,\JoinOp,0}{y}
  }[CanonicalProfilePairInfimumLowerInclusion]
    \proofstepwidestar[1]{\ZeroJoinSemi{A,\JoinOp,0}}{\rA}
    \proofstepwidestar[2]{x\in A}{\rA}
    \proofstepwidestar[3]{y\in A}{\rA}
    \proofstepwidestar[4]{m\in\LB_{A,\InducedLeq}(\{x,y\})}{\rA}
    \proofstepwidestar[5]{%
      \forall d\in\LB_{A,\InducedLeq}(\{x,y\})\,d\InducedLeq m}{\rA}
    \proofstepwidestar[1]{\Semi{A,\JoinOp}}{%
      \FormulaRefAuto{ZeroJoinBaseSemilatticeAxiom}[ax]{1}}
    \proofstepwidestar[1]{\PartOrd{A,\InducedLeq}}{%
      \FormulaRefAuto{JoinSemilatticeInducesOrder}[thm]{6}}
    \proofstepwidestar[1,2,3]{%
      m\in\LB_{A,\InducedLeq}(\{x,y\})\leftrightarrow
      \bigl(m\in A\land m\InducedLeq x\land m\InducedLeq y\bigr)}{%
      \FormulaRefAuto{PairLowerBoundSetCharacterization}[thm]{7,2,3}}
    \proofstepwidestar[1,2,3,4]{%
      m\in A\land m\InducedLeq x\land m\InducedLeq y}{%
      \FormulaRefAuto{P\leftrightarrow Q, P\vdash Q}{8,4}}
    \proofstepwidestar[1,2,3,4]{m\InducedLeq x}{%
      \FormulaRefAuto{P\land Q\land R\vdash Q}{9}}
    \proofstepwidestar[1,2,3,4]{m\InducedLeq y}{%
      \FormulaRefAuto{P\land Q\land R\vdash R}{9}}
    \proofstepwidestar[12]{j\in\IrrBelow{A,\JoinOp,0}{m}}{\rA}
    \proofstepwidestar[12]{%
      j\in\JoinIrrSet{A,\JoinOp,0}\land j\InducedLeq m}{%
      \FormulaRefAuto{IrreduciblesBelow}[def]{12}}
    \proofstepwidestar[12]{j\in\JoinIrrSet{A,\JoinOp,0}}{\rAEa{13}}
    \proofstepwidestar[12]{j\InducedLeq m}{\rAEb{13}}
    \proofstepwidestar[1,2,3,4,12]{j\InducedLeq x}{%
      \FormulaRefAuto{JoinOrderTransitive}[thm]{6,15,10}}
    \proofstepwidestar[1,2,3,4,12]{j\InducedLeq y}{%
      \FormulaRefAuto{JoinOrderTransitive}[thm]{6,15,11}}
    \proofstepwidestar[1,2,3,4,12]{j\in\IrrBelow{A,\JoinOp,0}{x}}{%
      \FormulaRefAuto{IrreduciblesBelow}[def]{\rAI{14,16}}}
    \proofstepwidestar[1,2,3,4,12]{j\in\IrrBelow{A,\JoinOp,0}{y}}{%
      \FormulaRefAuto{IrreduciblesBelow}[def]{\rAI{14,17}}}
    \proofstepwidestar[1,2,3,4,12]{%
      j\in\IrrBelow{A,\JoinOp,0}{x}\cap\IrrBelow{A,\JoinOp,0}{y}}{%
      \FormulaRefAuto{x \in A \cap B \eqvdash x \in A \land x \in B}{%
        \rAI{18,19}}}
    \proofstepwidestar[1,2,3,4]{%
      \begin{aligned}
        \IrrBelow{A,\JoinOp,0}{m}&\subseteq
          \IrrBelow{A,\JoinOp,0}{x}\cap{}\\[-2pt]
        &\quad\IrrBelow{A,\JoinOp,0}{y}
      \end{aligned}}{%
      \ProofRefStack{%
        \FormulaRefAuto{A \subseteq B \coloneq
          \forall x\,(x\in A\rightarrow x\in B)}\\
        \bigl(\rUI{\rRI{12,20}}\bigr)}}
  \closeproofpartwideindR

  \proofpartwideindR[Zweite Teilmengenrichtung]{%
    \begin{aligned}[t]
      &\ZeroJoinSemi{A,\JoinOp,0}\dsep x,y\in A\dsep
        m\in\LB_{A,\InducedLeq}(\{x,y\})\dsep{}\\[-2pt]
      &\forall d\in\LB_{A,\InducedLeq}(\{x,y\})\,
        d\InducedLeq m\vdash{}\\[-2pt]
      &\IrrBelow{A,\JoinOp,0}{x}\cap
        \IrrBelow{A,\JoinOp,0}{y}\subseteq
        \IrrBelow{A,\JoinOp,0}{m}
    \end{aligned}
  }{%
    \ZeroJoinSemi{A,\JoinOp,0}\dsep x,y\in A\dsep
    m\in\LB_{A,\InducedLeq}(\{x,y\})\dsep
    \forall d\in\LB_{A,\InducedLeq}(\{x,y\})\,d\InducedLeq m\vdash
    \IrrBelow{A,\JoinOp,0}{x}\cap\IrrBelow{A,\JoinOp,0}{y}
      \subseteq\IrrBelow{A,\JoinOp,0}{m}
  }[CanonicalProfilePairInfimumUpperInclusion]
    \proofstepwidestar[1]{\ZeroJoinSemi{A,\JoinOp,0}}{\rA}
    \proofstepwidestar[2]{x\in A}{\rA}
    \proofstepwidestar[3]{y\in A}{\rA}
    \proofstepwidestar[4]{m\in\LB_{A,\InducedLeq}(\{x,y\})}{\rA}
    \proofstepwidestar[5]{%
      \forall d\in\LB_{A,\InducedLeq}(\{x,y\})\,d\InducedLeq m}{\rA}
    \proofstepwidestar[1]{\Semi{A,\JoinOp}}{%
      \FormulaRefAuto{ZeroJoinBaseSemilatticeAxiom}[ax]{1}}
    \proofstepwidestar[1]{\PartOrd{A,\InducedLeq}}{%
      \FormulaRefAuto{JoinSemilatticeInducesOrder}[thm]{6}}
    \proofstepwidestar[8]{%
      j\in\IrrBelow{A,\JoinOp,0}{x}\cap\IrrBelow{A,\JoinOp,0}{y}}{\rA}
    \proofstepwidestar[8]{j\in\IrrBelow{A,\JoinOp,0}{x}}{%
      \FormulaRefAuto{x \in A \cap B \vdash x \in A}{8}}
    \proofstepwidestar[8]{j\in\IrrBelow{A,\JoinOp,0}{y}}{%
      \FormulaRefAuto{x \in A \cap B \vdash x \in B}{8}}
    \proofstepwidestar[8]{%
      j\in\JoinIrrSet{A,\JoinOp,0}\land j\InducedLeq x}{%
      \FormulaRefAuto{IrreduciblesBelow}[def]{9}}
    \proofstepwidestar[8]{%
      j\in\JoinIrrSet{A,\JoinOp,0}\land j\InducedLeq y}{%
      \FormulaRefAuto{IrreduciblesBelow}[def]{10}}
    \proofstepwidestar[8]{j\in\JoinIrrSet{A,\JoinOp,0}}{\rAEa{11}}
    \proofstepwidestar[8]{j\InducedLeq x}{\rAEb{11}}
    \proofstepwidestar[8]{j\InducedLeq y}{\rAEb{12}}
    \proofstepwidestar[8]{\JoinIrr{A,\JoinOp,0}{j}}{%
      \FormulaRefAuto{JoinIrreducibleSet}[def]{13}}
    \proofstepwidestar[8]{j\in A}{%
      \FormulaRefAuto{JoinIrreducible}[def]{16}}
    \proofstepwidestar[1,2,3]{%
      j\in\LB_{A,\InducedLeq}(\{x,y\})\leftrightarrow
      \bigl(j\in A\land j\InducedLeq x\land j\InducedLeq y\bigr)}{%
      \FormulaRefAuto{PairLowerBoundSetCharacterization}[thm]{7,2,3}}
    \proofstepwidestar[8]{%
      j\in A\land j\InducedLeq x\land j\InducedLeq y}{%
      \rAI{17,\rAI{14,15}}}
    \proofstepwidestar[1,2,3,8]{j\in\LB_{A,\InducedLeq}(\{x,y\})}{%
      \FormulaRefAuto{P\leftrightarrow Q, Q\vdash P}{18,19}}
    \proofstepwidestar[1,2,3,5,8]{j\InducedLeq m}{%
      \rRE{20,\rUE{5}}}
    \proofstepwidestar[1,2,3,5,8]{j\in\IrrBelow{A,\JoinOp,0}{m}}{%
      \FormulaRefAuto{IrreduciblesBelow}[def]{\rAI{13,21}}}
    \proofstepwidestar[1,2,3,5]{%
      \begin{aligned}
        \IrrBelow{A,\JoinOp,0}{x}\cap
          \IrrBelow{A,\JoinOp,0}{y}&\subseteq{}\\[-2pt]
        &\quad\IrrBelow{A,\JoinOp,0}{m}
      \end{aligned}}{%
      \ProofRefStack{%
        \FormulaRefAuto{A \subseteq B \coloneq
          \forall x\,(x\in A\rightarrow x\in B)}\\
        \bigl(\rUI{\rRI{8,22}}\bigr)}}
  \closeproofpartwideindR

  \proofpartwide{Schluss}
    \proofstepwidestar[1]{\ZeroJoinSemi{A,\JoinOp,0}}{\rA}
    \proofstepwidestar[2]{x\in A}{\rA}
    \proofstepwidestar[3]{y\in A}{\rA}
    \proofstepwidestar[4]{m\in\LB_{A,\InducedLeq}(\{x,y\})}{\rA}
    \proofstepwidestar[5]{%
      \forall d\in\LB_{A,\InducedLeq}(\{x,y\})\,d\InducedLeq m}{\rA}
    \proofstepwidestar[1,2,3,4,5]{%
      \begin{aligned}
        \IrrBelow{A,\JoinOp,0}{m}&\subseteq
          \IrrBelow{A,\JoinOp,0}{x}\cap{}\\[-2pt]
        &\quad\IrrBelow{A,\JoinOp,0}{y}
      \end{aligned}}{%
      \ProofRefStack{%
        \FormulaRefAuto{CanonicalProfilePairInfimumLowerInclusion}[thm]\\
        (1,2,3,4,5)}}
    \proofstepwidestar[1,2,3,4,5]{%
      \begin{aligned}
        \IrrBelow{A,\JoinOp,0}{x}\cap
          \IrrBelow{A,\JoinOp,0}{y}&\subseteq{}\\[-2pt]
        &\quad\IrrBelow{A,\JoinOp,0}{m}
      \end{aligned}}{%
      \ProofRefStack{%
        \FormulaRefAuto{CanonicalProfilePairInfimumUpperInclusion}[thm]\\
        (1,2,3,4,5)}}
    \proofstepwidestar[1,2,3,4,5]{%
      \begin{aligned}
        \IrrBelow{A,\JoinOp,0}{m}&=
          \IrrBelow{A,\JoinOp,0}{x}\cap{}\\[-2pt]
        &\quad\IrrBelow{A,\JoinOp,0}{y}
      \end{aligned}}{%
      \FormulaRefAuto{A \subseteq B, B \subseteq A \vdash A = B}{6,7}}
  \closeproofpartwide
\end{tabproofsplitwide}

\FormulaThmDeltaK[Profil eines Paarinfimums als Vereinigung]{%
  \begin{aligned}[t]
    &\ZeroJoinSemi{A,\JoinOp,0}\dsep x,y\in A\dsep
      m\in\LB_{A,\InducedLeq}(\{x,y\})\dsep{}\\
    &\forall d\in\LB_{A,\InducedLeq}(\{x,y\})\,
      d\InducedLeq m\vdash{}\\
    &\CanonicalProfileMap{A,\JoinOp,0}(m)
      =\CanonicalProfileMap{A,\JoinOp,0}(x)
       \cup\CanonicalProfileMap{A,\JoinOp,0}(y)
  \end{aligned}
}{CanonicalProfilePairInfimumUnion}{%
  \DeltaRow{Mengen}{A\dsep d\dsep m\dsep x\dsep y\dsep 0}
  \DeltaRow{Funktionen}{\JoinOp\dsep\CanonicalProfileMap{A,\JoinOp,0}}
}
\begin{tabproofwide}
  \proofstepwidestar[1]{\ZeroJoinSemi{A,\JoinOp,0}}{\rA}
  \proofstepwidestar[2]{x\in A}{\rA}
  \proofstepwidestar[3]{y\in A}{\rA}
  \proofstepwidestar[4]{m\in\LB_{A,\InducedLeq}(\{x,y\})}{\rA}
  \proofstepwidestar[5]{%
    \forall d\in\LB_{A,\InducedLeq}(\{x,y\})\,d\InducedLeq m}{\rA}
  \proofstepwidestar[1]{\Semi{A,\JoinOp}}{%
    \FormulaRefAuto{ZeroJoinBaseSemilatticeAxiom}[ax]{1}}
  \proofstepwidestar[1]{\PartOrd{A,\InducedLeq}}{%
    \FormulaRefAuto{JoinSemilatticeInducesOrder}[thm]{6}}
  \proofstepwidestar[2,3]{\{x,y\}\subseteq A}{%
    \FormulaRefAuto{a \in A,\, b \in A \vdash
      \{a,b\} \subseteq A}{2,3}}
  \proofstepwidestar[1,2,3]{%
    \LB_{A,\InducedLeq}(\{x,y\})\subseteq A}{%
    \FormulaRefAuto{LowerBoundSetSubsetCarrier}[thm]{7,8}}
  \proofstepwidestar[1,2,3,4]{m\in A}{%
    \FormulaRefAuto{A\subseteq B,\,x\in A\vdash x\in B}{9,4}}
  \proofstepwidestar[1,2,3,4,5]{%
    \begin{aligned}
      \IrrBelow{A,\JoinOp,0}{m}&=
        \IrrBelow{A,\JoinOp,0}{x}\cap{}\\[-2pt]
      &\quad\IrrBelow{A,\JoinOp,0}{y}
    \end{aligned}}{%
    \ProofRefStack{%
      \FormulaRefAuto{CanonicalProfilePairInfimumIntersection}[thm]\\
      (1,2,3,4,5)}}
  \proofstepwidestar[1,2,3,4]{%
    \CanonicalProfileMap{A,\JoinOp,0}(m)=
    \JoinIrrSet{A,\JoinOp,0}\setminus\IrrBelow{A,\JoinOp,0}{m}}{%
    \FormulaRefAuto{CanonicalProfileEvaluation}[thm]{1,10}}
  \proofstepwidestar[1,2]{%
    \CanonicalProfileMap{A,\JoinOp,0}(x)=
    \JoinIrrSet{A,\JoinOp,0}\setminus\IrrBelow{A,\JoinOp,0}{x}}{%
    \FormulaRefAuto{CanonicalProfileEvaluation}[thm]{1,2}}
  \proofstepwidestar[1,3]{%
    \CanonicalProfileMap{A,\JoinOp,0}(y)=
    \JoinIrrSet{A,\JoinOp,0}\setminus\IrrBelow{A,\JoinOp,0}{y}}{%
    \FormulaRefAuto{CanonicalProfileEvaluation}[thm]{1,3}}
  \proofstepwidestar[1,2,3,4,5]{%
    \begin{aligned}
      \CanonicalProfileMap{A,\JoinOp,0}(m)
        &=\CanonicalProfileMap{A,\JoinOp,0}(x)\\[-2pt]
        &\quad\cup\CanonicalProfileMap{A,\JoinOp,0}(y)
    \end{aligned}}{%
    \FormulaRefAuto{RelativeComplementProfileUnionTransport}[thm]{%
      11,12,13,14}}
\end{tabproofwide}

\FormulaThmDeltaK[Abgeschlossenheit der kanonischen Profilfamilie]{%
  \Finite A\dsep\ZeroJoinSemi{A,\JoinOp,0}
  \vdash\UnionClosedFam(\mathcal U_{A,\JoinOp,0})
}{CanonicalProfileFamilyUnionClosed}{%
  \DeltaRow{Mengen}{A\dsep X\dsep Y\dsep m\dsep x\dsep y\dsep 0}
  \DeltaRow{Funktionen}{\JoinOp\dsep\CanonicalProfileMap{A,\JoinOp,0}}
}
Im folgenden Beweis schreiben wir
\[
  \Phi\coloneqq\CanonicalProfileMap{A,\JoinOp,0},
  \qquad
  \mathcal U\coloneqq\mathcal U_{A,\JoinOp,0},
  \qquad
  L_{x,y}\coloneqq\LB_{A,\InducedLeq}(\{x,y\}).
\]
\begin{tabproofsplitwide}
  \proofpartwideindR[Vereinigung zweier kanonischer Profile]{%
    \begin{aligned}[t]
      &\Finite A\dsep\ZeroJoinSemi{A,\JoinOp,0}\dsep
        X,Y\in\mathcal U_{A,\JoinOp,0}\vdash{}\\[-2pt]
      &X\cup Y\in\mathcal U_{A,\JoinOp,0}
    \end{aligned}
  }{%
    \begin{aligned}[t]
      &\Finite A\dsep\ZeroJoinSemi{A,\JoinOp,0}\dsep{}\\[-2pt]
      &X,Y\in\mathcal U_{A,\JoinOp,0}\vdash
        X\cup Y\in\mathcal U_{A,\JoinOp,0}
    \end{aligned}
  }[CanonicalProfileUnionWitness]
    \proofstepwidestar[1]{\Finite A}{\rA}
    \proofstepwidestar[2]{\ZeroJoinSemi{A,\JoinOp,0}}{\rA}
    \proofstepwidestar[3]{X\in\mathcal U}{\rA}
    \proofstepwidestar[4]{Y\in\mathcal U}{\rA}
    \proofstepwidestar[2,3]{%
      \exists x\in A\,\bigl(X=\Phi(x)\bigr)}{%
      \FormulaRefAuto{CanonicalProfileRepresentative}[thm]{2,3}}
    \proofstepwidestar[5]{%
      x\in A\land X=\Phi(x)}{\rA}
    \proofstepwidestar[2,4]{%
      \exists y\in A\,\bigl(Y=\Phi(y)\bigr)}{%
      \FormulaRefAuto{CanonicalProfileRepresentative}[thm]{2,4}}
    \proofstepwidestar[7]{%
      y\in A\land Y=\Phi(y)}{\rA}
    \proofstepwidestar[1,2,5,7]{%
      \begin{aligned}[t]
        &\exists m\in A\,\Bigl(
          m\in L_{x,y}\land{}\\[-2pt]
        &\quad\bigl(\forall d\in L_{x,y}\,
          d\InducedLeq m\bigr)\land{}\\[-2pt]
        &\quad
          \PairInf{A}{\InducedLeq}{x}{y}{m}\Bigr)
      \end{aligned}}{%
      \FormulaRefAuto{FiniteZeroJoinSemilatticePairInfima}[thm]{%
        1,2,\rAEa{6},\rAEa{8}}}
    \proofstepwidestar[9]{%
      \begin{aligned}[t]
        &m\in A\land\Bigl(
          m\in L_{x,y}\land{}\\[-2pt]
        &\quad\bigl(\forall d\in L_{x,y}\,
          d\InducedLeq m\bigr)\land{}\\[-2pt]
        &\quad
          \PairInf{A}{\InducedLeq}{x}{y}{m}\Bigr)
      \end{aligned}}{\rA}
    \proofstepwidestar[2,9]{%
      \begin{aligned}
        \Phi(m)&=\Phi(x)\cup\Phi(y)
      \end{aligned}}{%
      \ProofRefStack{%
        \FormulaRefAuto{CanonicalProfilePairInfimumUnion}[thm]\\
        \bigl(2,\rAEa{6},\rAEa{8},\rAEa{\rAEb{10}},\bigr.\\[-2pt]
        \bigl.\rAEa{\rAEb{\rAEb{10}}}\bigr)}}
    \proofstepwidestar[5,7]{%
      \begin{aligned}
        X\cup Y&=\Phi(x)\cup\Phi(y)
      \end{aligned}}{%
      \FormulaRefAuto{a=b\dsep c=d \dsep P(b,d)\vdash P(a,c)}{%
        \rAEb{6},\rAEb{8},\rII}}
    \proofstepwidestar[2,5,7,9]{%
      X\cup Y=\Phi(m)}{%
      \FormulaRefAuto{a=b,\,c=b\vdash a=c}{12,11}}
    \proofstepwidestar[2]{\Phi\colon A\to\mathcal U}{%
      \FormulaRefAuto{CanonicalProfileFunction}[thm]{2}}
    \proofstepwidestar[2,9]{\Phi(m)\in\mathcal U}{%
      \FormulaRefAuto{F\colon A\to B\dsep x\in A\vdash F(x)\in B}{%
        14,\rAEa{10}}}
    \proofstepwidestar[2,5,7,9]{X\cup Y\in\mathcal U}{\rIE{13,15}}
    \proofstepwidestar[1,2,5,7]{X\cup Y\in\mathcal U}{\rEE{9,10,16}}
    \proofstepwidestar[1,2,5,4]{X\cup Y\in\mathcal U}{\rEE{7,8,17}}
    \proofstepwidestar[1,2,3,4]{X\cup Y\in\mathcal U}{\rEE{5,6,18}}
  \closeproofpartwideindR

  \proofpartwide{Schluss}
    \proofstepwidestar[1]{\Finite A}{\rA}
    \proofstepwidestar[2]{\ZeroJoinSemi{A,\JoinOp,0}}{\rA}
    \proofstepwidestar[3]{X\in\mathcal U_{A,\JoinOp,0}}{\rA}
    \proofstepwidestar[4]{Y\in\mathcal U_{A,\JoinOp,0}}{\rA}
    \proofstepwidestar[1,2,3,4]{%
      X\cup Y\in\mathcal U_{A,\JoinOp,0}}{%
      \FormulaRefAuto{CanonicalProfileUnionWitness}[thm]{1,2,3,4}}
    \proofstepwidestar[1,2]{\UnionClosedFam(\mathcal U_{A,\JoinOp,0})}{%
      \FormulaRefAuto{Vereinigungsabgeschlossene Familie}{%
        \rUI{\rRI{3,\rRI{4,5}}}}}
  \closeproofpartwide
\end{tabproofsplitwide}

\FormulaThmDeltaK[Das Nullprofil ist leer]{%
  \ZeroJoinSemi{A,\JoinOp,0}\vdash
  \IrrBelow{A,\JoinOp,0}{0}=\varnothing
}{CanonicalZeroProfileEmpty}{%
  \DeltaRow{Mengen}{A\dsep j\dsep 0}
  \DeltaRow{Funktionen}{\JoinOp}
}
\begin{tabproofwide}
  \proofstepwidestar[1]{\ZeroJoinSemi{A,\JoinOp,0}}{\rA}
  \proofstepwidestar[2]{j\in\IrrBelow{A,\JoinOp,0}{0}}{\rA}
  \proofstepwidestar[2]{%
    j\in\JoinIrrSet{A,\JoinOp,0}\land
      j\InducedLeq 0}{%
    \FormulaRefAuto{IrreduciblesBelow}[def]{2}}
  \proofstepwidestar[2]{j\in\JoinIrrSet{A,\JoinOp,0}}{\rAEa{3}}
  \proofstepwidestar[2]{j\InducedLeq 0}{\rAEb{3}}
  \proofstepwidestar[2]{\JoinIrr{A,\JoinOp,0}{j}}{%
    \FormulaRefAuto{JoinIrreducibleSet}[def]{4}}
  \proofstepwidestar[2]{j\in A}{\FormulaRefAuto{JoinIrreducible}[def]{6}}
  \proofstepwidestar[2]{j\neq0}{\FormulaRefAuto{JoinIrreducible}[def]{6}}
  \proofstepwidestar[1,2]{0\InducedLeq j}{%
    \FormulaRefAuto{ZeroJoinLeast}[thm]{1,7}}
  \proofstepwidestar[1]{\Semi{A,\JoinOp}}{%
    \FormulaRefAuto{ZeroJoinBaseSemilatticeAxiom}[ax]{1}}
  \proofstepwidestar[1,2]{j=0}{%
    \FormulaRefAuto{JoinOrderAntisymmetric}[thm]{10,5,9}}
  \proofstepwidestar[1,2]{\bot}{\rBI{11,8}}
  \proofstepwidestar[1]{j\notin\IrrBelow{A,\JoinOp,0}{0}}{\rCI{2,12}}
  \proofstepwidestar[1]{%
    \forall j\,j\notin\IrrBelow{A,\JoinOp,0}{0}}{\rUI{13}}
  \proofstepwidestar[1]{\IrrBelow{A,\JoinOp,0}{0}=\varnothing}{%
    \FormulaRefAuto{\forall x\, x\notin A\vdash A=\varnothing}{14}}
\end{tabproofwide}

\FormulaThmDeltaK[Die kanonische Profilfamilie ist nichtleer]{%
  \ZeroJoinSemi{A,\JoinOp,0}\vdash
  \mathcal U_{A,\JoinOp,0}\neq\varnothing
}{CanonicalProfileFamilyNonempty}{%
  \DeltaRow{Mengen}{A\dsep 0}
  \DeltaRow{Funktionen}{\JoinOp\dsep\CanonicalProfileMap{A,\JoinOp,0}}
}
\begin{tabproofwide}
  \proofstepwidestar[1]{\ZeroJoinSemi{A,\JoinOp,0}}{\rA}
  \proofstepwidestar[1]{0\in A}{\FormulaRefAuto{ZeroJoinCarrierAxiom}[ax]{1}}
  \proofstepwidestar[1]{%
    \CanonicalProfileMap{A,\JoinOp,0}\colon
      A\to\mathcal U_{A,\JoinOp,0}}{%
    \FormulaRefAuto{CanonicalProfileFunction}[thm]{1}}
  \proofstepwidestar[1]{%
    \CanonicalProfileMap{A,\JoinOp,0}(0)\in\mathcal U_{A,\JoinOp,0}}{%
    \FormulaRefAuto{F\colon A\to B\dsep x\in A\vdash F(x)\in B}{3,2}}
  \proofstepwidestar[1]{%
    \exists X\,X\in\mathcal U_{A,\JoinOp,0}}{\rEI{4}}
  \proofstepwidestar[1]{\mathcal U_{A,\JoinOp,0}\neq\varnothing}{%
    \FormulaRefAuto{A\neq\varnothing \eqvdash \exists x\,(x \in A)}{5}}
\end{tabproofwide}


\FormulaThmDeltaK[Der Träger der kanonischen Profilfamilie ist nichtleer]{%
  \begin{aligned}[t]
    &\Finite A\dsep\ZeroJoinSemi{A,\JoinOp,0}\dsep
      \exists x\in A\,(x\neq0)\vdash{}\\
    &\bigcup\mathcal U_{A,\JoinOp,0}\neq\varnothing
  \end{aligned}
}{CanonicalProfileCarrierNonempty}{%
  \DeltaRow{Mengen}{A\dsep j\dsep x\dsep 0}
  \DeltaRow{Funktionen}{\JoinOp\dsep\CanonicalProfileMap{A,\JoinOp,0}}
}
\begin{tabproofwide}
  \proofstepwidestar[1]{\Finite A}{\rA}
  \proofstepwidestar[2]{\ZeroJoinSemi{A,\JoinOp,0}}{\rA}
  \proofstepwidestar[3]{\exists x\in A\,(x\neq0)}{\rA}
  \proofstepwidestar[1,2,3]{\JoinIrrSet{A,\JoinOp,0}\neq\varnothing}{%
    \FormulaRefAuto{NontrivialFiniteZeroJoinHasIrreducible}[thm]{1,2,3}}
  \proofstepwidestar[1,2,3]{%
    \exists j\,j\in\JoinIrrSet{A,\JoinOp,0}}{%
    \FormulaRefAuto{A\neq\varnothing \eqvdash \exists x\,(x \in A)}{4}}
  \proofstepwidestar[6]{j\in\JoinIrrSet{A,\JoinOp,0}}{\rA}
  \proofstepwidestar[2]{0\in A}{\FormulaRefAuto{ZeroJoinCarrierAxiom}[ax]{2}}
  \proofstepwidestar[2]{%
    \CanonicalProfileMap{A,\JoinOp,0}(0)=
      \JoinIrrSet{A,\JoinOp,0}\setminus\IrrBelow{A,\JoinOp,0}{0}}{%
    \FormulaRefAuto{CanonicalProfileEvaluation}[thm]{2,7}}
  \proofstepwidestar[2]{\IrrBelow{A,\JoinOp,0}{0}=\varnothing}{%
    \FormulaRefAuto{CanonicalZeroProfileEmpty}[thm]{2}}
  \proofstepwidestar[]{j\notin\varnothing}{\FormulaRefAuto{x \not\in \varnothing}}
  \proofstepwidestar[2,6]{j\notin\IrrBelow{A,\JoinOp,0}{0}}{%
    \rIE{\FormulaRefAuto{a=b\vdash b=a}{9},10}}
  \proofstepwidestar[2,6]{%
    j\in\JoinIrrSet{A,\JoinOp,0}\setminus\IrrBelow{A,\JoinOp,0}{0}}{%
    \FormulaRefAuto{x \in A\dsep x\notin B\vdash x \in A\setminus B}{6,11}}
  \proofstepwidestar[2,6]{%
    j\in\CanonicalProfileMap{A,\JoinOp,0}(0)}{%
    \rIE{\FormulaRefAuto{a=b\vdash b=a}{8},12}}
  \proofstepwidestar[2]{%
    \CanonicalProfileMap{A,\JoinOp,0}(0)\in\mathcal U_{A,\JoinOp,0}}{%
    \FormulaRefAuto{F\colon A\to B\dsep x\in A\vdash F(x)\in B}{%
      \FormulaRefAuto{CanonicalProfileFunction}[thm]{2},7}}
  \proofstepwidestar[2,6]{j\in\bigcup\mathcal U_{A,\JoinOp,0}}{%
    \FormulaRefAuto{B\in A\dsep x\in B\vdash x \in \bigcup A}{14,13}}
  \proofstepwidestar[2,6]{%
    \exists j\,j\in\bigcup\mathcal U_{A,\JoinOp,0}}{\rEI{15}}
  \proofstepwidestar[1,2,3]{%
    \exists j\,j\in\bigcup\mathcal U_{A,\JoinOp,0}}{\rEE{5,6,16}}
  \proofstepwidestar[1,2,3]{\bigcup\mathcal U_{A,\JoinOp,0}\neq\varnothing}{%
    \FormulaRefAuto{A\neq\varnothing \eqvdash \exists x\,(x \in A)}{17}}
\end{tabproofwide}

\FormulaThmDeltaK[Kanonische Familie ist eine Frankl-Familie]{%
  \begin{aligned}[t]
    &\Finite A\dsep\ZeroJoinSemi{A,\JoinOp,0}\dsep
      \exists x\in A\,(x\neq0)\vdash{}\\
    &\FranklFam{\mathcal U_{A,\JoinOp,0}}
  \end{aligned}
}{CanonicalUnionFamilyIsFranklFamily}{%
  \DeltaRow{Mengen}{A\dsep X\dsep Y\dsep x\dsep y\dsep m\dsep 0}
  \DeltaRow{Funktionen}{\JoinOp\dsep\CanonicalProfileMap{A,\JoinOp,0}}
}
\begin{tabproofwide}
  \proofstepwidestar[1]{\Finite A}{\rA}
  \proofstepwidestar[2]{\ZeroJoinSemi{A,\JoinOp,0}}{\rA}
  \proofstepwidestar[3]{\exists x\in A\,(x\neq0)}{\rA}
  \proofstepwidestar[1,2]{%
    \CanonicalProfileMap{A,\JoinOp,0}
      \colon A\bij\mathcal U_{A,\JoinOp,0}}{%
    \FormulaRefAuto{CanonicalProfileBijection}[thm]{1,2}}
  \proofstepwidestar[1,2]{\Finite{\mathcal U_{A,\JoinOp,0}}}{%
    \FormulaRefAuto{FiniteBijectionTransport}[thm]{1,4}}
  \proofstepwidestar[2]{\mathcal U_{A,\JoinOp,0}\neq\varnothing}{%
    \FormulaRefAuto{CanonicalProfileFamilyNonempty}[thm]{2}}
  \proofstepwidestar[1,2]{\UnionClosedFam(\mathcal U_{A,\JoinOp,0})}{%
    \FormulaRefAuto{CanonicalProfileFamilyUnionClosed}[thm]{1,2}}
  \proofstepwidestar[1,2,3]{%
    \bigcup\mathcal U_{A,\JoinOp,0}\neq\varnothing}{%
    \FormulaRefAuto{CanonicalProfileCarrierNonempty}[thm]{1,2,3}}
  \proofstepwidestar[1,2,3]{\FranklFam{\mathcal U_{A,\JoinOp,0}}}{%
    \FormulaRefAuto{Frankl-Familie}[def]{6,8,5,7}}
\end{tabproofwide}

\FormulaThmDeltaK[Mitglieder der kanonischen Profilfamilie liegen im Irreduziblenträger]{%
  X\in\mathcal U_{A,\JoinOp,0}\vdash
  X\subseteq\JoinIrrSet{A,\JoinOp,0}
}{CanonicalProfileMemberSubsetIrr}{%
  \DeltaRow{Mengen}{A\dsep X\dsep x\dsep 0}
  \DeltaRow{Funktionen}{\JoinOp}
}
\begin{tabproofwide}
  \proofstepwidestar[1]{X\in\mathcal U_{A,\JoinOp,0}}{\rA}
  \proofstepwidestar[1]{%
    \exists x\in A\,\bigl(
      X=\JoinIrrSet{A,\JoinOp,0}\setminus
        \IrrBelow{A,\JoinOp,0}{x}\bigr)}{%
    \FormulaRefAuto{CanonicalProfileFamilyMembership}[thm]{1}}
  \proofstepwidestar[3]{%
    x\in A\land
    X=\JoinIrrSet{A,\JoinOp,0}\setminus
      \IrrBelow{A,\JoinOp,0}{x}}{\rA}
  \proofstepwidestar[]{%
    \JoinIrrSet{A,\JoinOp,0}\setminus
      \IrrBelow{A,\JoinOp,0}{x}
      \subseteq\JoinIrrSet{A,\JoinOp,0}}{%
    \FormulaRefAuto{A\setminus B\subseteq A}}
  \proofstepwidestar[3]{X\subseteq\JoinIrrSet{A,\JoinOp,0}}{%
    \rIE{\rAEb{3},4}}
  \proofstepwidestar[1]{X\subseteq\JoinIrrSet{A,\JoinOp,0}}{%
    \rEE{2,3,5}}
\end{tabproofwide}

\FormulaThmDeltaK[Elemente des kanonischen Familienträgers sind irreduzibel]{%
  j\in\bigcup\mathcal U_{A,\JoinOp,0}\vdash
  j\in\JoinIrrSet{A,\JoinOp,0}
}{CanonicalCarrierElementIrreducible}{%
  \DeltaRow{Mengen}{A\dsep X\dsep j\dsep 0}
  \DeltaRow{Funktionen}{\JoinOp}
}
\begin{tabproofwide}
  \proofstepwidestar[1]{j\in\bigcup\mathcal U_{A,\JoinOp,0}}{\rA}
  \proofstepwidestar[1]{%
    \exists X\in\mathcal U_{A,\JoinOp,0}\,(j\in X)}{%
    \FormulaRefAuto{x \in \bigcup A \;\eqvdash\;
      \exists B\in A\,(x \in B)}{1}}
  \proofstepwidestar[3]{%
    X\in\mathcal U_{A,\JoinOp,0}\land j\in X}{\rA}
  \proofstepwidestar[3]{X\subseteq\JoinIrrSet{A,\JoinOp,0}}{%
    \FormulaRefAuto{CanonicalProfileMemberSubsetIrr}[thm]{\rAEa{3}}}
  \proofstepwidestar[3]{j\in\JoinIrrSet{A,\JoinOp,0}}{%
    \FormulaRefAuto{A\subseteq B,\,x\in A\vdash x\in B}{4,\rAEb{3}}}
  \proofstepwidestar[1]{j\in\JoinIrrSet{A,\JoinOp,0}}{%
    \rEE{2,3,5}}
\end{tabproofwide}


\FormulaThmDeltaK[Vermeidende Profile entsprechen dem Hauptfilter]{%
  \begin{aligned}[t]
    &\ZeroJoinSemi{A,\JoinOp,0}\dsep
      j\in\JoinIrrSet{A,\JoinOp,0}\dsep x\in A\vdash{}\\
    &\CanonicalProfileMap{A,\JoinOp,0}(x)
      \in\FamAvoiding{\mathcal U_{A,\JoinOp,0}}{j}
      \leftrightarrow
      x\in\PrincipalFilter{A,\InducedLeq}{j}
  \end{aligned}
}{CanonicalAvoidingMembership}{%
  \DeltaRow{Mengen}{A\dsep j\dsep x\dsep 0}
  \DeltaRow{Funktionen}{\JoinOp\dsep\CanonicalProfileMap{A,\JoinOp,0}}
}
\begin{tabproofsplitwide}
  \proofpartwide{\(\vdash\)}
    \proofstepwidestar[1]{\ZeroJoinSemi{A,\JoinOp,0}}{\rA}
    \proofstepwidestar[2]{j\in\JoinIrrSet{A,\JoinOp,0}}{\rA}
    \proofstepwidestar[3]{x\in A}{\rA}
    \proofstepwidestar[4]{%
      \CanonicalProfileMap{A,\JoinOp,0}(x)
        \in\FamAvoiding{\mathcal U_{A,\JoinOp,0}}{j}}{\rA}
    \proofstepwidestar[4]{%
      \CanonicalProfileMap{A,\JoinOp,0}(x)\in\mathcal U_{A,\JoinOp,0}
      \land j\notin\CanonicalProfileMap{A,\JoinOp,0}(x)}{%
      \FormulaRefAuto{Vermeidende Teilfamilie}[def]{4}}
    \proofstepwidestar[4]{%
      j\notin\CanonicalProfileMap{A,\JoinOp,0}(x)}{\rAEb{5}}
    \proofstepwidestar[1,3]{%
      \CanonicalProfileMap{A,\JoinOp,0}(x)=
      \JoinIrrSet{A,\JoinOp,0}\setminus\IrrBelow{A,\JoinOp,0}{x}}{%
      \FormulaRefAuto{CanonicalProfileEvaluation}[thm]{1,3}}
    \proofstepwidestar[1,3,4]{%
      j\notin\JoinIrrSet{A,\JoinOp,0}\setminus
        \IrrBelow{A,\JoinOp,0}{x}}{\rIE{7,6}}
    \proofstepwidestar[2]{%
      j\notin\JoinIrrSet{A,\JoinOp,0}\setminus
        \IrrBelow{A,\JoinOp,0}{x}
      \leftrightarrow j\in\IrrBelow{A,\JoinOp,0}{x}}{%
      \FormulaRefAuto{RelCompNotMembership}[thm]{2}}
    \proofstepwidestar[1,2,3,4]{j\in\IrrBelow{A,\JoinOp,0}{x}}{%
      \FormulaRefAuto{P\leftrightarrow Q, P\vdash Q}{9,8}}
    \proofstepwidestar[1,2,3,4]{j\InducedLeq x}{%
      \rAEb{\FormulaRefAuto{IrreduciblesBelow}[def]{10}}}
    \proofstepwidestar[1,2,3,4]{%
      x\in\PrincipalFilter{A,\InducedLeq}{j}}{%
      \FormulaRefAuto{PrincipalFilter}[def]{\rAI{3,11}}}
  \closeproofpartwide

  \proofpartwide{\(\dashv\)}
    \proofstepwidestar[1]{\ZeroJoinSemi{A,\JoinOp,0}}{\rA}
    \proofstepwidestar[2]{j\in\JoinIrrSet{A,\JoinOp,0}}{\rA}
    \proofstepwidestar[3]{x\in A}{\rA}
    \proofstepwidestar[4]{%
      x\in\PrincipalFilter{A,\InducedLeq}{j}}{\rA}
    \proofstepwidestar[4]{j\InducedLeq x}{%
      \rAEb{\FormulaRefAuto{PrincipalFilter}[def]{4}}}
    \proofstepwidestar[2,4]{j\in\IrrBelow{A,\JoinOp,0}{x}}{%
      \FormulaRefAuto{IrreduciblesBelow}[def]{\rAI{2,5}}}
    \proofstepwidestar[2]{%
      j\notin\JoinIrrSet{A,\JoinOp,0}\setminus
        \IrrBelow{A,\JoinOp,0}{x}
      \leftrightarrow j\in\IrrBelow{A,\JoinOp,0}{x}}{%
      \FormulaRefAuto{RelCompNotMembership}[thm]{2}}
    \proofstepwidestar[2,4]{%
      j\notin\JoinIrrSet{A,\JoinOp,0}\setminus
        \IrrBelow{A,\JoinOp,0}{x}}{%
      \FormulaRefAuto{P\leftrightarrow Q, Q\vdash P}{7,6}}
    \proofstepwidestar[1,3]{%
      \CanonicalProfileMap{A,\JoinOp,0}(x)=
      \JoinIrrSet{A,\JoinOp,0}\setminus\IrrBelow{A,\JoinOp,0}{x}}{%
      \FormulaRefAuto{CanonicalProfileEvaluation}[thm]{1,3}}
    \proofstepwidestar[1,2,3,4]{%
      j\notin\CanonicalProfileMap{A,\JoinOp,0}(x)}{%
      \rIE{\FormulaRefAuto{a=b\vdash b=a}{9},8}}
    \proofstepwidestar[1]{%
      \CanonicalProfileMap{A,\JoinOp,0}\colon
        A\to\mathcal U_{A,\JoinOp,0}}{%
      \FormulaRefAuto{CanonicalProfileFunction}[thm]{1}}
    \proofstepwidestar[1,3]{%
      \CanonicalProfileMap{A,\JoinOp,0}(x)\in\mathcal U_{A,\JoinOp,0}}{%
      \FormulaRefAuto{F\colon A\to B\dsep x\in A\vdash F(x)\in B}{11,3}}
    \proofstepwidestar[1,2,3,4]{%
      \CanonicalProfileMap{A,\JoinOp,0}(x)
        \in\FamAvoiding{\mathcal U_{A,\JoinOp,0}}{j}}{%
      \FormulaRefAuto{Vermeidende Teilfamilie}[def]{\rAI{12,10}}}
  \closeproofpartwide
\end{tabproofsplitwide}

\FormulaThmDeltaK[Enthaltende Profile entsprechen dem Filterkomplement]{%
  \begin{aligned}[t]
    &\ZeroJoinSemi{A,\JoinOp,0}\dsep
      j\in\JoinIrrSet{A,\JoinOp,0}\dsep x\in A\vdash{}\\
    &\CanonicalProfileMap{A,\JoinOp,0}(x)
      \in\FamContaining{\mathcal U_{A,\JoinOp,0}}{j}
      \leftrightarrow
      x\in A\setminus\PrincipalFilter{A,\InducedLeq}{j}
  \end{aligned}
}{CanonicalContainingMembership}{%
  \DeltaRow{Mengen}{A\dsep j\dsep x\dsep 0}
  \DeltaRow{Funktionen}{\JoinOp\dsep\CanonicalProfileMap{A,\JoinOp,0}}
}
\begin{tabproofsplitwide}
  \proofpartwide{\(\vdash\)}
    \proofstepwidestar[1]{\ZeroJoinSemi{A,\JoinOp,0}}{\rA}
    \proofstepwidestar[2]{j\in\JoinIrrSet{A,\JoinOp,0}}{\rA}
    \proofstepwidestar[3]{x\in A}{\rA}
    \proofstepwidestar[4]{%
      \CanonicalProfileMap{A,\JoinOp,0}(x)
        \in\FamContaining{\mathcal U_{A,\JoinOp,0}}{j}}{\rA}
    \proofstepwidestar[4]{%
      j\in\CanonicalProfileMap{A,\JoinOp,0}(x)}{%
      \rAEb{\FormulaRefAuto{Enthaltende Teilfamilie}[def]{4}}}
    \proofstepwidestar[6]{%
      x\in\PrincipalFilter{A,\InducedLeq}{j}}{\rA}
    \proofstepwidestar[1,2,3,6]{%
      \CanonicalProfileMap{A,\JoinOp,0}(x)
        \in\FamAvoiding{\mathcal U_{A,\JoinOp,0}}{j}}{%
      \FormulaRefAuto{CanonicalAvoidingMembership}[thm]{1,2,3,6}}
    \proofstepwidestar[7]{%
      j\notin\CanonicalProfileMap{A,\JoinOp,0}(x)}{%
      \rAEb{\FormulaRefAuto{Vermeidende Teilfamilie}[def]{7}}}
    \proofstepwidestar[4,6]{\bot}{\rBI{5,8}}
    \proofstepwidestar[4]{%
      x\notin\PrincipalFilter{A,\InducedLeq}{j}}{%
      \rCI{6,9}}
    \proofstepwidestar[3,4]{%
      x\in A\setminus\PrincipalFilter{A,\InducedLeq}{j}}{%
      \FormulaRefAuto{x \in A\dsep x\notin B\vdash x \in A\setminus B}{3,10}}
  \closeproofpartwide

  \proofpartwide{\(\dashv\)}
    \proofstepwidestar[1]{\ZeroJoinSemi{A,\JoinOp,0}}{\rA}
    \proofstepwidestar[2]{j\in\JoinIrrSet{A,\JoinOp,0}}{\rA}
    \proofstepwidestar[3]{x\in A}{\rA}
    \proofstepwidestar[4]{%
      x\in A\setminus\PrincipalFilter{A,\InducedLeq}{j}}{\rA}
    \proofstepwidestar[4]{%
      x\notin\PrincipalFilter{A,\InducedLeq}{j}}{%
      \FormulaRefAuto{x\in A\setminus B\vdash x\notin B}{4}}
    \proofstepwidestar[1]{%
      \CanonicalProfileMap{A,\JoinOp,0}\colon
        A\to\mathcal U_{A,\JoinOp,0}}{%
      \FormulaRefAuto{CanonicalProfileFunction}[thm]{1}}
    \proofstepwidestar[1,3]{%
      \CanonicalProfileMap{A,\JoinOp,0}(x)\in\mathcal U_{A,\JoinOp,0}}{%
      \FormulaRefAuto{F\colon A\to B\dsep x\in A\vdash F(x)\in B}{6,3}}
    \proofstepwidestar[8]{%
      j\notin\CanonicalProfileMap{A,\JoinOp,0}(x)}{\rA}
    \proofstepwidestar[1,3,8]{%
      \CanonicalProfileMap{A,\JoinOp,0}(x)
        \in\FamAvoiding{\mathcal U_{A,\JoinOp,0}}{j}}{%
      \FormulaRefAuto{Vermeidende Teilfamilie}[def]{\rAI{7,8}}}
    \proofstepwidestar[1,2,3,8]{%
      x\in\PrincipalFilter{A,\InducedLeq}{j}}{%
      \FormulaRefAuto{CanonicalAvoidingMembership}[thm]{1,2,3,9}}
    \proofstepwidestar[4,8]{\bot}{\rBI{10,5}}
    \proofstepwidestar[4]{%
      \neg\bigl(j\notin\CanonicalProfileMap{A,\JoinOp,0}(x)\bigr)}{%
      \rCI{8,11}}
    \proofstepwidestar[4]{%
      j\in\CanonicalProfileMap{A,\JoinOp,0}(x)}{%
      \FormulaRefAuto{rDN}{12}}
    \proofstepwidestar[1,3,4]{%
      \CanonicalProfileMap{A,\JoinOp,0}(x)
        \in\FamContaining{\mathcal U_{A,\JoinOp,0}}{j}}{%
      \FormulaRefAuto{Enthaltende Teilfamilie}[def]{\rAI{7,13}}}
  \closeproofpartwide
\end{tabproofsplitwide}

\FormulaThmDeltaK[Urbild der vermeidenden kanonischen Profile]{%
  \begin{aligned}[t]
    &\ZeroJoinSemi{A,\JoinOp,0}\dsep
      j\in\JoinIrrSet{A,\JoinOp,0}\vdash{}\\
    &\bigl(\CanonicalProfileMap{A,\JoinOp,0}\bigr)^{-1}
      [\FamAvoiding{\mathcal U_{A,\JoinOp,0}}{j}]
      =\PrincipalFilter{A,\InducedLeq}{j}
  \end{aligned}
}{CanonicalAvoidingPreimage}{%
  \DeltaRow{Mengen}{A\dsep j\dsep x\dsep 0}
  \DeltaRow{Funktionen}{\JoinOp\dsep\CanonicalProfileMap{A,\JoinOp,0}}
}
\begin{tabproofwide}
  \proofstepwidestar[1]{\ZeroJoinSemi{A,\JoinOp,0}}{\rA}
  \proofstepwidestar[2]{j\in\JoinIrrSet{A,\JoinOp,0}}{\rA}
  \proofstepwidestar[1]{%
    \CanonicalProfileMap{A,\JoinOp,0}\colon
      A\to\mathcal U_{A,\JoinOp,0}}{%
    \FormulaRefAuto{CanonicalProfileFunction}[thm]{1}}
  \proofstepwidestar[]{%
    \FamAvoiding{\mathcal U_{A,\JoinOp,0}}{j}
      \subseteq\mathcal U_{A,\JoinOp,0}}{%
    \FormulaRefAuto{\FamAvoiding{M}{a}\subseteq M}}
  \proofstepwidestar[]{%
    \PrincipalFilter{A,\InducedLeq}{j}\subseteq A}{%
    \FormulaRefAuto{PrincipalFilterSubsetCarrier}[thm]}
  \proofstepwidestar[3]{x\in A}{\rA}
  \proofstepwidestar[1,2,3]{%
    \CanonicalProfileMap{A,\JoinOp,0}(x)
      \in\FamAvoiding{\mathcal U_{A,\JoinOp,0}}{j}
      \leftrightarrow
      x\in\PrincipalFilter{A,\InducedLeq}{j}}{%
    \FormulaRefAuto{CanonicalAvoidingMembership}[thm]{1,2,6}}
  \proofstepwidestar[1,2]{%
    \forall x\in A\,\bigl(
      \CanonicalProfileMap{A,\JoinOp,0}(x)
        \in\FamAvoiding{\mathcal U_{A,\JoinOp,0}}{j}
      \leftrightarrow
      x\in\PrincipalFilter{A,\InducedLeq}{j}\bigr)}{%
    \rUI{\rRI{6,7}}}
  \proofstepwidestar[1,2]{%
    \bigl(\CanonicalProfileMap{A,\JoinOp,0}\bigr)^{-1}
      [\FamAvoiding{\mathcal U_{A,\JoinOp,0}}{j}]
      =\PrincipalFilter{A,\InducedLeq}{j}}{%
    \FormulaRefAuto{PreimageEqFromPointwiseIff}[thm]{3,4,5,8}}
\end{tabproofwide}

\FormulaThmDeltaK[Urbild der enthaltenden kanonischen Profile]{%
  \begin{aligned}[t]
    &\ZeroJoinSemi{A,\JoinOp,0}\dsep
      j\in\JoinIrrSet{A,\JoinOp,0}\vdash{}\\
    &\bigl(\CanonicalProfileMap{A,\JoinOp,0}\bigr)^{-1}
      [\FamContaining{\mathcal U_{A,\JoinOp,0}}{j}]
      =A\setminus\PrincipalFilter{A,\InducedLeq}{j}
  \end{aligned}
}{CanonicalContainingPreimage}{%
  \DeltaRow{Mengen}{A\dsep j\dsep x\dsep 0}
  \DeltaRow{Funktionen}{\JoinOp\dsep\CanonicalProfileMap{A,\JoinOp,0}}
}
\begin{tabproofwide}
  \proofstepwidestar[1]{\ZeroJoinSemi{A,\JoinOp,0}}{\rA}
  \proofstepwidestar[2]{j\in\JoinIrrSet{A,\JoinOp,0}}{\rA}
  \proofstepwidestar[1]{%
    \CanonicalProfileMap{A,\JoinOp,0}\colon
      A\to\mathcal U_{A,\JoinOp,0}}{%
    \FormulaRefAuto{CanonicalProfileFunction}[thm]{1}}
  \proofstepwidestar[]{%
    \FamContaining{\mathcal U_{A,\JoinOp,0}}{j}
      \subseteq\mathcal U_{A,\JoinOp,0}}{%
    \FormulaRefAuto{\FamContaining{M}{a}\subseteq M}}
  \proofstepwidestar[]{%
    A\setminus\PrincipalFilter{A,\InducedLeq}{j}\subseteq A}{%
    \FormulaRefAuto{A\setminus B\subseteq A}}
  \proofstepwidestar[3]{x\in A}{\rA}
  \proofstepwidestar[1,2,3]{%
    \CanonicalProfileMap{A,\JoinOp,0}(x)
      \in\FamContaining{\mathcal U_{A,\JoinOp,0}}{j}
      \leftrightarrow
      x\in A\setminus\PrincipalFilter{A,\InducedLeq}{j}}{%
    \FormulaRefAuto{CanonicalContainingMembership}[thm]{1,2,6}}
  \proofstepwidestar[1,2]{%
    \forall x\in A\,\bigl(
      \CanonicalProfileMap{A,\JoinOp,0}(x)
        \in\FamContaining{\mathcal U_{A,\JoinOp,0}}{j}
      \leftrightarrow
      x\in A\setminus\PrincipalFilter{A,\InducedLeq}{j}\bigr)}{%
    \rUI{\rRI{6,7}}}
  \proofstepwidestar[1,2]{%
    \bigl(\CanonicalProfileMap{A,\JoinOp,0}\bigr)^{-1}
      [\FamContaining{\mathcal U_{A,\JoinOp,0}}{j}]
      =A\setminus\PrincipalFilter{A,\InducedLeq}{j}}{%
    \FormulaRefAuto{PreimageEqFromPointwiseIff}[thm]{3,4,5,8}}
\end{tabproofwide}

\FormulaThmDeltaK[Injektionstransport entlang einer Bijektion]{%
  c\colon A\bij B\dsep H\subseteq B\dsep\AtMostHalf{H}{B}
  \vdash\exists F\colon c^{-1}[H]\inj c^{-1}[B\setminus H]
}{BijectionAtMostHalfTransport}{%
  \DeltaRow{Mengen}{A\dsep B\dsep H}
  \DeltaRow{Funktionen}{c\dsep F\dsep G}
}
\begin{tabproofwide}
  \proofstepwidestar[1]{c\colon A\bij B}{\rA}
  \proofstepwidestar[2]{H\subseteq B}{\rA}
  \proofstepwidestar[3]{\AtMostHalf{H}{B}}{\rA}
  \proofstepwidestar[1,2]{%
    c\restriction_{c^{-1}[H]}\colon c^{-1}[H]\bij H}{%
    \FormulaRefAuto{F\colon A\bij B \dsep T\subseteq B \vdash
      F\restriction_{F^{-1}[T]}\colon F^{-1}[T]\bij T}{1,2}}
  \proofstepwidestar[]{B\setminus H\subseteq B}{%
    \FormulaRefAuto{A\setminus B\subseteq A}}
  \proofstepwidestar[1]{%
    c\restriction_{c^{-1}[B\setminus H]}\colon
    c^{-1}[B\setminus H]\bij B\setminus H}{%
    \FormulaRefAuto{F\colon A\bij B \dsep T\subseteq B \vdash
      F\restriction_{F^{-1}[T]}\colon F^{-1}[T]\bij T}{1,5}}
  \proofstepwidestar[1]{%
    \left(c\restriction_{c^{-1}[B\setminus H]}\right)^{-1}
    \colon B\setminus H\bij c^{-1}[B\setminus H]}{%
    \FormulaRefAuto{F\colon A\bij B\vdash F^{-1}\colon B\bij A}{6}}
  \proofstepwidestar[2,3]{\exists G\colon H\inj B\setminus H}{%
    \FormulaRefAuto{AtMostHalf}[def]{2,3}}
  \proofstepwidestar[9]{G\colon H\inj B\setminus H}{\rA}
  \proofstepwidestar[1,2]{%
    c\restriction_{c^{-1}[H]}\colon c^{-1}[H]\inj H}{%
    \FormulaRefAuto{F\colon A\bij B\vdash F\colon A\inj B}{4}}
  \proofstepwidestar[1]{%
    \left(c\restriction_{c^{-1}[B\setminus H]}\right)^{-1}
      \colon B\setminus H\inj c^{-1}[B\setminus H]}{%
    \FormulaRefAuto{F\colon A\bij B\vdash F\colon A\inj B}{7}}
  \proofstepwidestar[1,2,9]{%
    G\circ(c\restriction_{c^{-1}[H]})
      \colon c^{-1}[H]\inj B\setminus H}{%
    \FormulaRefAuto{F\colon A\inj B \dsep G\colon B\inj C
      \vdash G\circ F\colon A\inj C}{10,9}}
  \proofstepwidestar[1,2,9]{%
    \left(c\restriction_{c^{-1}[B\setminus H]}\right)^{-1}
      \circ G\circ(c\restriction_{c^{-1}[H]})
      \colon c^{-1}[H]\inj c^{-1}[B\setminus H]}{%
    \FormulaRefAuto{F\colon A\inj B \dsep G\colon B\inj C
      \vdash G\circ F\colon A\inj C}{12,11}}
  \proofstepwidestar[1,2,9]{%
    \exists F\colon c^{-1}[H]\inj c^{-1}[B\setminus H]}{\rEI{13}}
  \proofstepwidestar[1,2,3]{%
    \exists F\colon c^{-1}[H]\inj c^{-1}[B\setminus H]}{%
    \rEE{8,9,14}}
\end{tabproofwide}

\FormulaThmDeltaK[Transport eines halbhäufigen Profilelements]{%
  \begin{aligned}[t]
    &\Finite A\dsep\ZeroJoinSemi{A,\JoinOp,0}\dsep
      \exists x\in A\,(x\neq0)\dsep
      \HalfOcc{j}{\mathcal U_{A,\JoinOp,0}}\vdash{}\\
    &\FranklSemi{A,\JoinOp,0}
  \end{aligned}
}{CanonicalHalfOccurrenceTransport}{%
  \DeltaRow{Mengen}{A\dsep j\dsep x\dsep 0}
  \DeltaRow{Funktionen}{%
    \JoinOp\dsep\CanonicalProfileMap{A,\JoinOp,0}\dsep F\dsep G}
}

\paragraph{Beweisskizze.}
Für die kanonische Bijektion
\(\CanonicalProfileMap{A,\JoinOp,0}\) gilt für jedes \(x\in A\)
\[
  \begin{aligned}
  j\notin \CanonicalProfileMap{A,\JoinOp,0}(x)
  &\ \Longleftrightarrow\
    j\in\IrrBelow{A,\JoinOp,0}{x}\\
  &\ \Longleftrightarrow\
    x\in\PrincipalFilter{A,\InducedLeq}{j}.
  \end{aligned}
\]
Da jedes Mitglied von \(\mathcal U_{A,\JoinOp,0}\) eine Teilmenge von
\(\JoinIrrSet{A,\JoinOp,0}\) ist, liefert
\(j\in\bigcup\mathcal U_{A,\JoinOp,0}\) zugleich die benötigte
\(\JoinOp\)-Irreduzibilität. Die in \(\HalfOcc{j}{\mathcal U_{A,\JoinOp,0}}\) enthaltene
Injektion wird nun entlang der kanonischen Profilabbildung transportiert.

\begin{tabproofsplitwide}
  \proofpartwideindR[Halbhäufigkeit der vermeidenden Teilfamilie]{%
    \begin{aligned}[t]
      &\HalfOcc{j}{\mathcal U_{A,\JoinOp,0}}\vdash{}\\[-2pt]
      &\AtMostHalf{\FamAvoiding{\mathcal U_{A,\JoinOp,0}}{j}}{
        \mathcal U_{A,\JoinOp,0}}
    \end{aligned}
  }{\HalfOcc{j}{\mathcal U_{A,\JoinOp,0}}\vdash
    \AtMostHalf{\FamAvoiding{\mathcal U_{A,\JoinOp,0}}{j}}{
      \mathcal U_{A,\JoinOp,0}}}[CanonicalAvoidingAtMostHalf]
    \proofstepwidestar[1]{\HalfOcc{j}{\mathcal U_{A,\JoinOp,0}}}{\rA}
    \proofstepwidestar[]{%
      \FamAvoiding{\mathcal U_{A,\JoinOp,0}}{j}
        \subseteq\mathcal U_{A,\JoinOp,0}}{%
      \FormulaRefAuto{\FamAvoiding{M}{a}\subseteq M}}
    \proofstepwidestar[1]{%
      \exists F\colon
        \FamAvoiding{\mathcal U_{A,\JoinOp,0}}{j}
        \inj\FamContaining{\mathcal U_{A,\JoinOp,0}}{j}}{%
      \FormulaRefAuto{Halbhäufiges Element}[def]{1}}
    \proofstepwidestar[]{%
      \mathcal U_{A,\JoinOp,0}\setminus
        \FamAvoiding{\mathcal U_{A,\JoinOp,0}}{j}
        =\FamContaining{\mathcal U_{A,\JoinOp,0}}{j}}{%
      \FormulaRefAuto{FamContainingAsAvoidingRelativeComplement}[thm]}
    \proofstepwidestar[3]{%
      F\colon\FamAvoiding{\mathcal U_{A,\JoinOp,0}}{j}
        \inj\FamContaining{\mathcal U_{A,\JoinOp,0}}{j}}{\rA}
    \proofstepwidestar[]{%
      \FamContaining{\mathcal U_{A,\JoinOp,0}}{j}
        =\mathcal U_{A,\JoinOp,0}\setminus
          \FamAvoiding{\mathcal U_{A,\JoinOp,0}}{j}}{%
      \FormulaRefAuto{a=b\vdash b=a}{4}}
    \proofstepwidestar[3]{%
      \begin{aligned}
        F\colon\FamAvoiding{\mathcal U_{A,\JoinOp,0}}{j}
          \inj{}&\mathcal U_{A,\JoinOp,0}\setminus\\[-2pt]
          &\FamAvoiding{\mathcal U_{A,\JoinOp,0}}{j}
      \end{aligned}}{\rIE{6,5}}
    \proofstepwidestar[3]{%
      \exists F\colon
        \FamAvoiding{\mathcal U_{A,\JoinOp,0}}{j}
        \inj\mathcal U_{A,\JoinOp,0}\setminus
          \FamAvoiding{\mathcal U_{A,\JoinOp,0}}{j}}{\rEI{7}}
    \proofstepwidestar[1]{%
      \exists F\colon
        \FamAvoiding{\mathcal U_{A,\JoinOp,0}}{j}
        \inj\mathcal U_{A,\JoinOp,0}\setminus
          \FamAvoiding{\mathcal U_{A,\JoinOp,0}}{j}}{\rEE{3,5,8}}
    \proofstepwidestar[1]{%
      \AtMostHalf{\FamAvoiding{\mathcal U_{A,\JoinOp,0}}{j}}{
        \mathcal U_{A,\JoinOp,0}}}{%
      \FormulaRefAuto{AtMostHalf}[def]{2,9}}
  \closeproofpartwideindR

  \proofpartwideindR[Transport in den Hauptfilter]{%
    \begin{aligned}[t]
      &\Finite A\dsep\ZeroJoinSemi{A,\JoinOp,0}\dsep
        j\in\JoinIrrSet{A,\JoinOp,0}\dsep
        \HalfOcc{j}{\mathcal U_{A,\JoinOp,0}}\vdash{}\\[-2pt]
      &\exists G\colon
        \PrincipalFilter{A,\InducedLeq}{j}
        \inj A\setminus\PrincipalFilter{A,\InducedLeq}{j}
    \end{aligned}
  }{\Finite A\dsep\ZeroJoinSemi{A,\JoinOp,0}\dsep
    j\in\JoinIrrSet{A,\JoinOp,0}\dsep
    \HalfOcc{j}{\mathcal U_{A,\JoinOp,0}}\vdash
    \exists G\colon\PrincipalFilter{A,\InducedLeq}{j}
    \inj A\setminus\PrincipalFilter{A,\InducedLeq}{j}}%
  [CanonicalHalfOccurrenceFilterInjection]
    \proofstepwidestar[1]{\Finite A}{\rA}
    \proofstepwidestar[2]{\ZeroJoinSemi{A,\JoinOp,0}}{\rA}
    \proofstepwidestar[3]{j\in\JoinIrrSet{A,\JoinOp,0}}{\rA}
    \proofstepwidestar[4]{\HalfOcc{j}{\mathcal U_{A,\JoinOp,0}}}{\rA}
    \proofstepwidestar[1,2]{%
      \CanonicalProfileMap{A,\JoinOp,0}
        \colon A\bij\mathcal U_{A,\JoinOp,0}}{%
      \FormulaRefAuto{CanonicalProfileBijection}[thm]{1,2}}
    \proofstepwidestar[]{%
      \FamAvoiding{\mathcal U_{A,\JoinOp,0}}{j}
        \subseteq\mathcal U_{A,\JoinOp,0}}{%
      \FormulaRefAuto{\FamAvoiding{M}{a}\subseteq M}}
    \proofstepwidestar[4]{%
      \AtMostHalf{\FamAvoiding{\mathcal U_{A,\JoinOp,0}}{j}}{
        \mathcal U_{A,\JoinOp,0}}}{%
      \FormulaRefAuto{CanonicalAvoidingAtMostHalf}[thm]{4}}
    \proofstepwidestar[1,2,4]{%
      \begin{aligned}
        &\exists G\colon
        \bigl(\CanonicalProfileMap{A,\JoinOp,0}\bigr)^{-1}
          [\FamAvoiding{\mathcal U_{A,\JoinOp,0}}{j}]\inj{}\\[-2pt]
        &\quad\bigl(\CanonicalProfileMap{A,\JoinOp,0}\bigr)^{-1}
          [\mathcal U_{A,\JoinOp,0}\setminus
            \FamAvoiding{\mathcal U_{A,\JoinOp,0}}{j}]
      \end{aligned}}{%
      \FormulaRefAuto{BijectionAtMostHalfTransport}[thm]{5,6,7}}
    \proofstepwidestar[]{%
      \mathcal U_{A,\JoinOp,0}\setminus
        \FamAvoiding{\mathcal U_{A,\JoinOp,0}}{j}
        =\FamContaining{\mathcal U_{A,\JoinOp,0}}{j}}{%
      \FormulaRefAuto{FamContainingAsAvoidingRelativeComplement}[thm]}
    \proofstepwidestar[1,2,4]{%
      \begin{aligned}
        &\exists G\colon
        \bigl(\CanonicalProfileMap{A,\JoinOp,0}\bigr)^{-1}
          [\FamAvoiding{\mathcal U_{A,\JoinOp,0}}{j}]\inj{}\\[-2pt]
        &\quad\bigl(\CanonicalProfileMap{A,\JoinOp,0}\bigr)^{-1}
          [\FamContaining{\mathcal U_{A,\JoinOp,0}}{j}]
      \end{aligned}}{\rIE{9,8}}
    \proofstepwidestar[2,3]{%
      \bigl(\CanonicalProfileMap{A,\JoinOp,0}\bigr)^{-1}
        [\FamAvoiding{\mathcal U_{A,\JoinOp,0}}{j}]
        =\PrincipalFilter{A,\InducedLeq}{j}}{%
      \FormulaRefAuto{CanonicalAvoidingPreimage}[thm]{2,3}}
    \proofstepwidestar[1,2,3,4]{%
      \exists G\colon
        \PrincipalFilter{A,\InducedLeq}{j}\inj
        \bigl(\CanonicalProfileMap{A,\JoinOp,0}\bigr)^{-1}
          [\FamContaining{\mathcal U_{A,\JoinOp,0}}{j}]}{\rIE{11,10}}
    \proofstepwidestar[2,3]{%
      \bigl(\CanonicalProfileMap{A,\JoinOp,0}\bigr)^{-1}
        [\FamContaining{\mathcal U_{A,\JoinOp,0}}{j}]
        =A\setminus\PrincipalFilter{A,\InducedLeq}{j}}{%
      \FormulaRefAuto{CanonicalContainingPreimage}[thm]{2,3}}
    \proofstepwidestar[1,2,3,4]{%
      \exists G\colon
        \PrincipalFilter{A,\InducedLeq}{j}
        \inj A\setminus\PrincipalFilter{A,\InducedLeq}{j}}{%
      \rIE{13,12}}
  \closeproofpartwideindR

  \proofpartwide{Schluss}
    \proofstepwidestar[1]{\Finite A}{\rA}
    \proofstepwidestar[2]{\ZeroJoinSemi{A,\JoinOp,0}}{\rA}
    \proofstepwidestar[3]{\exists x\in A\,(x\neq0)}{\rA}
    \proofstepwidestar[4]{\HalfOcc{j}{\mathcal U_{A,\JoinOp,0}}}{\rA}
    \proofstepwidestar[4]{j\in\bigcup\mathcal U_{A,\JoinOp,0}}{%
      \FormulaRefAuto{Halbhäufiges Element}[def]{4}}
    \proofstepwidestar[5]{j\in\JoinIrrSet{A,\JoinOp,0}}{%
      \FormulaRefAuto{CanonicalCarrierElementIrreducible}[thm]{5}}
    \proofstepwidestar[2,5]{\JoinIrr{A,\JoinOp,0}{j}}{%
      \FormulaRefAuto{JoinIrreducibleSet}[def]{6}}
    \proofstepwidestar[1,2,4,5]{%
      \exists G\colon
        \PrincipalFilter{A,\InducedLeq}{j}
        \inj A\setminus\PrincipalFilter{A,\InducedLeq}{j}}{%
      \FormulaRefAuto{CanonicalHalfOccurrenceFilterInjection}[thm]{1,2,6,4}}
    \proofstepwidestar[]{%
      \PrincipalFilter{A,\InducedLeq}{j}\subseteq A}{%
      \FormulaRefAuto{PrincipalFilterSubsetCarrier}[thm]}
    \proofstepwidestar[1,2,4,5]{%
      \AtMostHalf{\PrincipalFilter{A,\InducedLeq}{j}}{A}}{%
      \FormulaRefAuto{AtMostHalf}[def]{9,8}}
    \proofstepwidestar[1,2,3,4]{\FranklSemi{A,\JoinOp,0}}{%
      \FormulaRefAuto{FranklSemilatticeProperty}[def]{%
        \rEI{\rAI{7,10}}}}
  \closeproofpartwide
\end{tabproofsplitwide}

\FormulaThmDeltaK[Mengenform impliziert Halbverbandsform]{%
  \begin{aligned}[t]
    &\forall M\,(\FranklFam M\to\Frankl M)\dsep
      \Finite A\dsep{}\\
    &\ZeroJoinSemi{A,\JoinOp,0}\dsep
      \exists x\in A\,(x\neq0)\vdash{}\\
    &\FranklSemi{A,\JoinOp,0}
  \end{aligned}
}{FranklSetImpliesSemilattice}{%
  \DeltaRow{Mengen}{A\dsep M\dsep x\dsep 0}
  \DeltaRow{Funktionen}{\JoinOp}
}
\begin{tabproofwide}
  \proofstepwidestar[1]{\forall M\,(\FranklFam M\to\Frankl M)}{\rA}
  \proofstepwidestar[2]{\Finite A}{\rA}
  \proofstepwidestar[3]{\ZeroJoinSemi{A,\JoinOp,0}}{\rA}
  \proofstepwidestar[4]{\exists x\in A\,(x\neq0)}{\rA}
  \proofstepwidestar[2,3,4]{\FranklFam{\mathcal U_{A,\JoinOp,0}}}{%
    \FormulaRefAuto{CanonicalUnionFamilyIsFranklFamily}[thm]{2,3,4}}
  \proofstepwidestar[1,2,3,4]{\Frankl{\mathcal U_{A,\JoinOp,0}}}{%
    \rRE{\rUE{1},5}}
  \proofstepwidestar[1,2,3,4]{%
    \exists j\,\HalfOcc{j}{\mathcal U_{A,\JoinOp,0}}}{%
    \FormulaRefAuto{Frankl-Eigenschaft}[def]{6}}
  \proofstepwidestar[8]{\HalfOcc{j}{\mathcal U_{A,\JoinOp,0}}}{\rA}
  \proofstepwidestar[2,3,4,8]{\FranklSemi{A,\JoinOp,0}}{%
    \FormulaRefAuto{CanonicalHalfOccurrenceTransport}[thm]{2,3,4,8}}
  \proofstepwidestar[1,2,3,4]{\FranklSemi{A,\JoinOp,0}}{\rEE{7,8,9}}
\end{tabproofwide}

\section{Von der Vereinigungsfamilie zum Halbverband}

\FormulaDefDeltaK[Schnittabgeschlossene Familie]{%
  \InterClosedFam L
  \coloneqq
  \forall X\in L\,\forall Y\in L\,(X\cap Y\in L)
}{IntersectionClosedFamily}{%
  \DeltaRow{Mengen}{L\dsep X\dsep Y}
  \DeltaRow{\textbf{Neues Symbol}}{}
  \DeltaPrem{Schnittabgeschlossene Familien}{\InterClosedFam L}
}

\par\medskip
Sei nun \(M\) eine Frankl-Familie. Wir verwenden die folgenden drei
kanonischen Objekte:
\[
  \begin{aligned}
    U_M&\coloneqq\bigcup M,&
    \mathcal L_M
      &\coloneqq\{U_M\setminus X\mid X\in M^\circ\}.
  \end{aligned}
\]
\par\noindent
Für \(X,Y\in\mathcal L_M\) setzen wir außerdem
\[
  X\FamilyJoinOp{M}{}Y
  \coloneqq
  \bigcap\{Z\in\mathcal L_M\mid X\subseteq Z\land Y\subseteq Z\}.
\]
Die Menge unter dem Durchschnitt ist nichtleer, weil sie \(U_M\) enthält.
Auf \(\mathcal L_M\) verwenden wir durchweg die Inklusionsordnung
\(\subseteq\). Weiter unten zeigen wir punktweise, dass sie mit der von
\(\FamilyJoinOp{M}\) induzierten Ordnung übereinstimmt.

\FormulaDefDeltaK[Komplementduale Konstruktion]{%
  \begin{aligned}[t]
    U_M&\coloneqq\bigcup M,&
    \mathcal L_M
      &\coloneqq\{U_M\setminus X\mid X\in M^\circ\},\\
    X\FamilyJoinOp{M}{}Y
      &\coloneqq
        \bigcap\{Z\in\mathcal L_M\mid X\cup Y\subseteq Z\}.
  \end{aligned}
}{FranklComplementConstruction}{%
  \DeltaRow{Mengen}{M\dsep X\dsep Y\dsep Z}
  \DeltaRow{Relationsoperatoren}{\subseteq}
  \DeltaRow{Funktionen}{\FamilyJoinOp{M}}
}

\paragraph{Beweisskizze.}
Aus \(M\) entsteht zunächst \(M^\circ\) und durch Komplementbildung die
schnittabgeschlossene Familie \(\mathcal L_M\). Der Durchschnitt aller
gemeinsamen Obermengen von \(X\) und \(Y\) ist ihr Paarsupremum; dadurch wird
\(\mathcal L_M\) zu einem endlichen Halbverband mit Nullelement. Liefert die
Halbverbandsform ein irreduzibles \(P\), so wählt ein privater Punkt von \(P\)
genau die Mitglieder des Hauptfilters aus. Die Komplementbijektion
transportiert die zugehörige Injektion schließlich zurück zu einem
halbhäufigen Element der ursprünglichen Familie \(M\):
\[
  M\longrightarrow M^\circ\longrightarrow\mathcal L_M
  \longrightarrow P\longrightarrow\text{privater Punkt in }M.
\]

\FormulaThmDeltaK[Beschränkter Existenzquantor über einer Adjunktion]{%
  \exists Y\in C\cup\{X\}\,P(Y)
  \leftrightarrow
  \bigl(\exists Y\in C\,P(Y)\bigr)\lor P(X)
}{BoundedExistentialAdjunction}{%
  \DeltaRow{Mengen}{C\dsep X\dsep Y}
  \DeltaRow{Prädikatensymbole}{P}
}
\begin{tabproofwide}
  \proofstepwidestar[1]{\exists Y\in C\cup\{X\}\,P(Y)}{\rA}
  \proofstepwidestar[2]{Y\in C\cup\{X\}\land P(Y)}{\rA}
  \proofstepwidestar[2]{Y\in C\cup\{X\}}{\rAEa{2}}
  \proofstepwidestar[2]{P(Y)}{\rAEb{2}}
  \proofstepwidestar[2]{Y\in C\lor Y=X}{%
    \FormulaRefAuto{x\in A\cup\{a\}\vdash x\in A\lor x=a}{3}}
  \proofstepwidestar[6]{Y\in C}{\rA}
  \proofstepwidestar[2,6]{\exists Y\in C\,P(Y)}{\rEI{\rAI{6,4}}}
  \proofstepwidestar[2,6]{\bigl(\exists Y\in C\,P(Y)\bigr)\lor P(X)}{%
    \rOIa{7}}
  \proofstepwidestar[9]{Y=X}{\rA}
  \proofstepwidestar[2,9]{P(X)}{\rIE{9,4}}
  \proofstepwidestar[2,9]{\bigl(\exists Y\in C\,P(Y)\bigr)\lor P(X)}{%
    \rOIb{10}}
  \proofstepwidestar[2]{\bigl(\exists Y\in C\,P(Y)\bigr)\lor P(X)}{%
    \rOE{5,6,8,9,11}}
  \proofstepwidestar[1]{\bigl(\exists Y\in C\,P(Y)\bigr)\lor P(X)}{%
    \rEE{1,2,12}}
  \proofstepwidestar[]{\exists Y\in C\cup\{X\}\,P(Y)\to
    \bigl((\exists Y\in C\,P(Y))\lor P(X)\bigr)}{\rRI{1,13}}
  \proofstepwidestar[15]{\bigl(\exists Y\in C\,P(Y)\bigr)\lor P(X)}{\rA}
  \proofstepwidestar[16]{\exists Y\in C\,P(Y)}{\rA}
  \proofstepwidestar[17]{Y\in C\land P(Y)}{\rA}
  \proofstepwidestar[17]{Y\in C\cup\{X\}}{%
    \FormulaRefAuto{z\in A\vdash z\in A\cup B}{\rAEa{17}}}
  \proofstepwidestar[17]{\exists Y\in C\cup\{X\}\,P(Y)}{%
    \rEI{\rAI{18,\rAEb{17}}}}
  \proofstepwidestar[16]{\exists Y\in C\cup\{X\}\,P(Y)}{%
    \rEE{16,17,19}}
  \proofstepwidestar[21]{P(X)}{\rA}
  \proofstepwidestar[]{X\in C\cup\{X\}}{\FormulaRefAuto{a\in A\cup\{a\}}}
  \proofstepwidestar[21]{\exists Y\in C\cup\{X\}\,P(Y)}{%
    \rEI{\rAI{22,21}}}
  \proofstepwidestar[15]{\exists Y\in C\cup\{X\}\,P(Y)}{%
    \rOE{15,16,20,21,23}}
  \proofstepwidestar[]{\bigl((\exists Y\in C\,P(Y))\lor P(X)\bigr)
    \to\exists Y\in C\cup\{X\}\,P(Y)}{\rRI{15,24}}
  \proofstepwidestar[]{\exists Y\in C\cup\{X\}\,P(Y)
    \leftrightarrow\bigl((\exists Y\in C\,P(Y))\lor P(X)\bigr)}{%
    \FormulaRefAuto{P\leftrightarrow Q\eqvdash
      (P\to Q)\land(Q\to P)}{\rAI{14,25}}}
\end{tabproofwide}

\FormulaThmDeltaK[Beschränkter Allquantor über einer Einermenge]{%
  \forall Y\in\{X\}\,P(Y)\leftrightarrow P(X)
}{BoundedUniversalSingleton}{%
  \DeltaRow{Mengen}{X\dsep Y}
  \DeltaRow{Prädikatensymbole}{P}
}
\begin{tabproof}
  \proofstep{1}{\forall Y\in\{X\}\,P(Y)}{\rA}
  \proofstep{}{X\in\{X\}}{\FormulaRefAuto{a\in\{a\}}}
  \proofstep{1}{P(X)}{\rRE{2,\rUE{1}}}
  \proofstep{}{\forall Y\in\{X\}\,P(Y)\to P(X)}{\rRI{1,3}}
  \proofstep{5}{P(X)}{\rA}
  \proofstep{6}{Y\in\{X\}}{\rA}
  \proofstep{6}{Y=X}{\FormulaRefAuto{x\in\{a\}\eqvdash x=a}{6}}
  \proofstep{5,6}{P(Y)}{\rIE{\FormulaRefAuto{a=b\vdash b=a}{7},5}}
  \proofstep{5}{\forall Y\in\{X\}\,P(Y)}{\rUI{\rRI{6,8}}}
  \proofstep{}{P(X)\to\forall Y\in\{X\}\,P(Y)}{\rRI{5,9}}
  \proofstep{}{\forall Y\in\{X\}\,P(Y)\leftrightarrow P(X)}{%
    \FormulaRefAuto{P\leftrightarrow Q\eqvdash
      (P\to Q)\land(Q\to P)}{\rAI{4,10}}}
\end{tabproof}

\FormulaThmDeltaK[Beschränkter Allquantor über einer Adjunktion]{%
  \forall Y\in C\cup\{X\}\,P(Y)
  \leftrightarrow
  \bigl(\forall Y\in C\,P(Y)\bigr)\land P(X)
}{BoundedUniversalAdjunction}{%
  \DeltaRow{Mengen}{C\dsep X\dsep Y}
  \DeltaRow{Prädikatensymbole}{P}
}
Zur Verkürzung des Beweises setzen wir
\[
  \alpha\coloneqq\forall Y\in C\cup\{X\}\,P(Y),
  \qquad
  \beta\coloneqq\forall Y\in C\,P(Y),
  \qquad
  \gamma\coloneqq P(X).
\]
\begin{tabproofwide}
  \proofstepwidestar[1]{\alpha}{\rA}
  \proofstepwidestar[2]{Y\in C}{\rA}
  \proofstepwidestar[2]{Y\in C\cup\{X\}}{%
    \FormulaRefAuto{z\in A\vdash z\in A\cup B}{2}}
  \proofstepwidestar[1,2]{P(Y)}{\rRE{3,\rUE{1}}}
  \proofstepwidestar[1]{\beta}{\rUI{\rRI{2,4}}}
  \proofstepwidestar[]{X\in C\cup\{X\}}{\FormulaRefAuto{a\in A\cup\{a\}}}
  \proofstepwidestar[1]{\gamma}{\rRE{6,\rUE{1}}}
  \proofstepwidestar[1]{\beta\land\gamma}{\rAI{5,7}}
  \proofstepwidestar[]{\alpha\to\beta\land\gamma}{\rRI{1,8}}
  \proofstepwidestar[10]{\beta\land\gamma}{\rA}
  \proofstepwidestar[11]{Y\in C\cup\{X\}}{\rA}
  \proofstepwidestar[11]{Y\in C\lor Y=X}{%
    \FormulaRefAuto{x\in A\cup\{a\}\vdash x\in A\lor x=a}{11}}
  \proofstepwidestar[13]{Y\in C}{\rA}
  \proofstepwidestar[10,13]{P(Y)}{\rRE{13,\rUE{\rAEa{10}}}}
  \proofstepwidestar[15]{Y=X}{\rA}
  \proofstepwidestar[10,15]{P(Y)}{%
    \rIE{\FormulaRefAuto{a=b\vdash b=a}{15},\rAEb{10}}}
  \proofstepwidestar[10,11]{P(Y)}{\rOE{12,13,14,15,16}}
  \proofstepwidestar[10]{\alpha}{\rUI{\rRI{11,17}}}
  \proofstepwidestar[]{\beta\land\gamma\to\alpha}{\rRI{10,18}}
  \proofstepwidestar[]{\alpha\leftrightarrow(\beta\land\gamma)}{%
    \FormulaRefAuto{P\leftrightarrow Q\eqvdash
      (P\to Q)\land(Q\to P)}{\rAI{9,19}}}
\end{tabproofwide}

\FormulaThmDeltaK[Vereinigung einer adjungierten Mengenfamilie]{%
  \bigcup(C\cup\{X\})=(\bigcup C)\cup X
}{UnionFamilyAdjunction}{%
  \DeltaRow{Mengen}{C\dsep X\dsep Y\dsep a}
}
\begin{tabproofwide}
  \proofstepwide{a\in\bigcup(C\cup\{X\})}{\leftrightarrow}{%
    \exists Y\in C\cup\{X\}\,(a\in Y)}{%
    \FormulaRefAuto{x \in \bigcup A \;\eqvdash\;\exists B\in A\,(x \in B)}}
  \proofstepwide{} {\leftrightarrow}{%
    \bigl(\exists Y\in C\,(a\in Y)\bigr)\lor a\in X}{%
    \FormulaRefAuto{BoundedExistentialAdjunction}[thm]}
  \proofstepwide{} {\leftrightarrow}{a\in\bigcup C\lor a\in X}{%
    \rLRS{\FormulaRefAuto{x \in \bigcup A \;\eqvdash\;\exists B\in A\,(x \in B)},2}}
  \proofstepwide{} {\leftrightarrow}{a\in(\bigcup C)\cup X}{%
    \FormulaRefAuto{z \in A \cup B \eqvdash z \in A \lor z \in B}{3}}
  \proofstepwidestar{\bigcup(C\cup\{X\})=(\bigcup C)\cup X}{%
    \FormulaRefAuto{\forall x\,(x \in A \leftrightarrow x \in B) \eqvdash A = B}{\rUI{4}}}
\end{tabproofwide}

\FormulaThmDeltaK[Durchschnitt einer einelementigen Mengenfamilie]{%
  \bigcap\{X\}=X
}{IntersectionSingletonFamily}{%
  \DeltaRow{Mengen}{X\dsep a}
}
\begin{tabproofwide}
  \proofstepwide{a\in\bigcap\{X\}}{\leftrightarrow}{%
    \forall Y\in\{X\}\,(a\in Y)}{%
    \FormulaRefAuto{M\neq\varnothing \vdash x\in\bigcap M \leftrightarrow
      \forall X\in M\,(x\in X)}{\FormulaRefAuto{a\in\{a\}}}}
  \proofstepwide{} {\leftrightarrow}{a\in X}{%
    \FormulaRefAuto{BoundedUniversalSingleton}[thm]}
  \proofstepwidestar{\bigcap\{X\}=X}{%
    \FormulaRefAuto{\forall x\,(x \in A \leftrightarrow x \in B) \eqvdash A = B}{\rUI{2}}}
\end{tabproofwide}

\FormulaThmDeltaK[Durchschnitt einer adjungierten nichtleeren Mengenfamilie]{%
  C\neq\varnothing\vdash
  \bigcap(C\cup\{X\})=(\bigcap C)\cap X
}{IntersectionFamilyAdjunction}{%
  \DeltaRow{Mengen}{C\dsep X\dsep Y\dsep a}
}
Im folgenden Beweis sei
\[
  \mathcal D\coloneqq C\cup\{X\}.
\]
\begin{tabproofwide}
  \proofstepwidestar[1]{C\neq\varnothing}{\rA}
  \proofstepwide[1]{a\in\bigcap\mathcal D}{\leftrightarrow}{%
    \forall Y\in\mathcal D\,(a\in Y)}{%
    \ProofRefStack{%
      \FormulaRefAuto{M\neq\varnothing \vdash
        x\in\bigcap M \leftrightarrow
        \forall X\in M\,(x\in X)}\\
      \FormulaRefAuto{a\in A\vdash A\neq\varnothing}[thm]\\
      \FormulaRefAuto{a\in A\cup\{a\}}[thm]}}
  \proofstepwide[1]{} {\leftrightarrow}{%
    \bigl(\forall Y\in C\,(a\in Y)\bigr)\land a\in X}{%
    \FormulaRefAuto{BoundedUniversalAdjunction}[thm]}
  \proofstepwide[1]{} {\leftrightarrow}{a\in\bigcap C\land a\in X}{%
    \ProofRefStack{%
      \FormulaRefAuto{M\neq\varnothing \vdash
        x\in\bigcap M \leftrightarrow
        \forall X\in M\,(x\in X)}\\
      \bigl(1,3\bigr)}}
  \proofstepwide[1]{} {\leftrightarrow}{a\in(\bigcap C)\cap X}{%
    \FormulaRefAuto{x\in A\cap B\eqvdash x\in A\land x\in B}{4}}
  \proofstepwidestar[1]{\bigcap\mathcal D=(\bigcap C)\cap X}{%
    \ProofRefStack{%
      \FormulaRefAuto{\forall x\,
        (x \in A \leftrightarrow x \in B) \eqvdash A = B}\\
      \bigl(\rUI{5}\bigr)}}
\end{tabproofwide}

\FormulaThmDeltaK[Der Familiendurchschnitt liegt in jedem Familienglied]{%
  M\neq\varnothing\dsep X\in M\vdash\bigcap M\subseteq X
}{FamilyIntersectionSubsetMember}{%
  \DeltaRow{Mengen}{M\dsep X\dsep a}
}
\begin{tabproof}
  \proofstep{1}{M\neq\varnothing}{\rA}
  \proofstep{2}{X\in M}{\rA}
  \proofstep{3}{a\in\bigcap M}{\rA}
  \proofstep{1,2,3}{a\in X}{%
    \FormulaRefAuto{M\neq\varnothing\dsep a\in\bigcap M\dsep
      X\in M\vdash a\in X}[thm]{1,3,2}}
  \proofstep{1,2}{\bigcap M\subseteq X}{%
    \FormulaRefAuto{A\subseteq B\coloneq
      \forall x\,(x\in A\rightarrow x\in B)}[def]{\rUI{\rRI{3,4}}}}
\end{tabproof}

\FormulaThmDeltaK[Endlicher nichtleerer Schnitt in einer schnittabgeschlossenen Familie]{%
  \Finite C\dsep C\neq\varnothing\dsep C\subseteq L\dsep
  \InterClosedFam L\vdash \bigcap C\in L
}{FiniteIntersectionClosedFamilyIntersection}{%
  \DeltaRow{Mengen}{C\dsep L\dsep X\dsep Y}
}
Setze für die Endlichkeitsinduktion
\[
  P(C)\quad:\Longleftrightarrow\quad
  \bigl(C\subseteq L\land C\neq\varnothing\to\bigcap C\in L\bigr).
\]
\begin{tabproofsplitwide}
  \proofpartwideindR[Induktionsanfang]{P(\varnothing)}{P(\varnothing)}%
    [FiniteIntersectionClosedBase]
    \proofstepwidestar[1]{\varnothing\subseteq L\land
      \varnothing\neq\varnothing}{\rA}
    \proofstepwidestar[]{\varnothing=\varnothing}{\rII}
    \proofstepwidestar[1]{\bot}{\rCE{2,\rAEb{1}}}
    \proofstepwidestar[]{P(\varnothing)}{\rRI{1,\FormulaRefAuto{\bot\vdash P}{3}}}
  \closeproofpartwideindR

  \proofpartwideindR[Induktionsschritt]{%
    \Finite C\dsep P(C)\dsep X\notin C\dsep\InterClosedFam L
    \vdash P(C\cup\{X\})
  }{%
    \Finite C\dsep P(C)\dsep X\notin C\dsep\InterClosedFam L
    \vdash P(C\cup\{X\})
  }[FiniteIntersectionClosedStep]
    \proofstepwidestar[1]{\Finite C}{\rA}
    \proofstepwidestar[2]{P(C)}{\rA}
    \proofstepwidestar[3]{X\notin C}{\rA}
    \proofstepwidestar[4]{\InterClosedFam L}{\rA}
    \proofstepwidestar[5]{%
      C\cup\{X\}\subseteq L\land C\cup\{X\}\neq\varnothing}{\rA}
    \proofstepwidestar[5]{X\in L}{%
      \FormulaRefAuto{A\subseteq B,\,x\in A\vdash x\in B}{%
        \rAEa{5},\FormulaRefAuto{a\in A\cup\{a\}}}}
    \proofstepwidestar[7]{C=\varnothing}{\rA}
    \proofstepwidestar[5,7]{C\cup\{X\}=\{X\}}{\rIE{7,\rII}}
    \proofstepwidestar[5,7]{\bigcap(C\cup\{X\})=X}{%
      \rIE{8,\FormulaRefAuto{IntersectionSingletonFamily}[thm]}}
    \proofstepwidestar[5,7]{\bigcap(C\cup\{X\})\in L}{\rIE{9,6}}
    \proofstepwidestar[11]{C\neq\varnothing}{\rA}
    \proofstepwidestar[5]{C\subseteq L}{%
      \FormulaRefAuto{A\subseteq B\dsep B\subseteq C\vdash A\subseteq C}{%
        \FormulaRefAuto{A\subseteq A\cup B},\rAEa{5}}}
    \proofstepwidestar[2,5,11]{\bigcap C\in L}{\rRE{2,\rAI{12,11}}}
    \proofstepwidestar[4,6,13]{(\bigcap C)\cap X\in L}{%
      \FormulaRefAuto{IntersectionClosedFamily}[def]{4,13,6}}
    \proofstepwidestar[11]{\bigcap(C\cup\{X\})=(\bigcap C)\cap X}{%
      \FormulaRefAuto{IntersectionFamilyAdjunction}[thm]{11}}
    \proofstepwidestar[2,4,5,11]{\bigcap(C\cup\{X\})\in L}{\rIE{15,14}}
    \proofstepwidestar[]{C=\varnothing\lor C\neq\varnothing}{%
      \FormulaRefAuto{P\lor\neg P}}
    \proofstepwidestar[2,4,5]{\bigcap(C\cup\{X\})\in L}{%
      \rOE{17,7,10,11,16}}
    \proofstepwidestar[1,2,3,4]{P(C\cup\{X\})}{\rRI{5,18}}
  \closeproofpartwideindR

  \proofpartwide{Induktionsschluss}
    \proofstepwidestar[1]{\Finite C}{\rA}
    \proofstepwidestar[2]{C\neq\varnothing}{\rA}
    \proofstepwidestar[3]{C\subseteq L}{\rA}
    \proofstepwidestar[4]{\InterClosedFam L}{\rA}
    \proofstepwidestar[1,4]{P(C)}{%
      \FormulaRefAuto{FiniteInductionRule}[ax]{IA,IS,1}}
    \proofstepwidestar[1,2,3,4]{\bigcap C\in L}{\rRE{5,\rAI{3,2}}}
  \closeproofpartwide
\end{tabproofsplitwide}

\FormulaThmDeltaK[Gesamtvereinigung einer endlichen vereinigungsabgeschlossenen Familie]{%
  \Finite M\dsep M\neq\varnothing\dsep\UnionClosedFam(M)
  \vdash\bigcup M\in M
}{FiniteUnionClosedFamilyTotalUnion}{%
  \DeltaRow{Mengen}{M\dsep C\dsep X}
}
Setze
\[
  Q(C)\quad:\Longleftrightarrow\quad
  \bigl(C\subseteq M\land C\neq\varnothing\to\bigcup C\in M\bigr).
\]
\begin{tabproofsplitwide}
  \proofpartwideindR[Induktionsanfang]{Q(\varnothing)}{Q(\varnothing)}%
    [FiniteUnionClosedBase]
    \proofstepwidestar[1]{\varnothing\subseteq M\land
      \varnothing\neq\varnothing}{\rA}
    \proofstepwidestar[]{\varnothing=\varnothing}{\rII}
    \proofstepwidestar[1]{\bot}{\rCE{2,\rAEb{1}}}
    \proofstepwidestar[]{Q(\varnothing)}{\rRI{1,\FormulaRefAuto{\bot\vdash P}{3}}}
  \closeproofpartwideindR

  \proofpartwideindR[Induktionsschritt]{%
    \Finite C\dsep Q(C)\dsep X\notin C\dsep\UnionClosedFam(M)
    \vdash Q(C\cup\{X\})
  }{%
    \Finite C\dsep Q(C)\dsep X\notin C\dsep\UnionClosedFam(M)
    \vdash Q(C\cup\{X\})
  }[FiniteUnionClosedStep]
    \proofstepwidestar[1]{\Finite C}{\rA}
    \proofstepwidestar[2]{Q(C)}{\rA}
    \proofstepwidestar[3]{X\notin C}{\rA}
    \proofstepwidestar[4]{\UnionClosedFam(M)}{\rA}
    \proofstepwidestar[5]{%
      C\cup\{X\}\subseteq M\land C\cup\{X\}\neq\varnothing}{\rA}
    \proofstepwidestar[5]{X\in M}{%
      \FormulaRefAuto{A\subseteq B,\,x\in A\vdash x\in B}{%
        \rAEa{5},\FormulaRefAuto{a\in A\cup\{a\}}}}
    \proofstepwidestar[7]{C=\varnothing}{\rA}
    \proofstepwidestar[5,7]{\bigcup(C\cup\{X\})=X}{%
      \rIE{7,\FormulaRefAuto{UnionFamilyAdjunction}[thm]}}
    \proofstepwidestar[5,7]{\bigcup(C\cup\{X\})\in M}{\rIE{8,6}}
    \proofstepwidestar[10]{C\neq\varnothing}{\rA}
    \proofstepwidestar[5]{C\subseteq M}{%
      \FormulaRefAuto{A\subseteq B\dsep B\subseteq C\vdash A\subseteq C}{%
        \FormulaRefAuto{A\subseteq A\cup B},\rAEa{5}}}
    \proofstepwidestar[2,5,10]{\bigcup C\in M}{\rRE{2,\rAI{11,10}}}
    \proofstepwidestar[4,6,12]{(\bigcup C)\cup X\in M}{%
      \FormulaRefAuto{X\in M\dsep Y\in M\vdash X\cup Y\in M}[ax]{12,6}}
    \proofstepwidestar[]{\bigcup(C\cup\{X\})=(\bigcup C)\cup X}{%
      \FormulaRefAuto{UnionFamilyAdjunction}[thm]}
    \proofstepwidestar[2,4,5,10]{\bigcup(C\cup\{X\})\in M}{\rIE{14,13}}
    \proofstepwidestar[]{C=\varnothing\lor C\neq\varnothing}{%
      \FormulaRefAuto{P\lor\neg P}}
    \proofstepwidestar[2,4,5]{\bigcup(C\cup\{X\})\in M}{%
      \rOE{16,7,9,10,15}}
    \proofstepwidestar[1,2,3,4]{Q(C\cup\{X\})}{\rRI{5,17}}
  \closeproofpartwideindR

  \proofpartwide{Induktionsschluss}
    \proofstepwidestar[1]{\Finite M}{\rA}
    \proofstepwidestar[2]{M\neq\varnothing}{\rA}
    \proofstepwidestar[3]{\UnionClosedFam(M)}{\rA}
    \proofstepwidestar[1,3]{Q(M)}{%
      \FormulaRefAuto{FiniteInductionRule}[ax]{IA,IS,1}}
    \proofstepwidestar[]{M\subseteq M}{\FormulaRefAuto{A\subseteq A}}
    \proofstepwidestar[1,2,3]{\bigcup M\in M}{\rRE{4,\rAI{5,2}}}
  \closeproofpartwide
\end{tabproofsplitwide}

\FormulaThmDeltaK[de-Morgan-Gesetz für relative Komplemente]{%
  (U\setminus X)\cap(U\setminus Y)=U\setminus(X\cup Y)
}{RelativeComplementUnionDeMorgan}{%
  \DeltaRow{Mengen}{U\dsep X\dsep Y\dsep a}
}
\begin{tabproofwide}
  \proofstepwide{a\in(U\setminus X)\cap(U\setminus Y)}{\leftrightarrow}{%
    a\in U\land a\notin X\land a\notin Y}{%
    \FormulaRefAuto{x\in A\cap B\eqvdash x\in A\land x\in B};
    \FormulaRefAuto{x \in A \setminus B \eqvdash x \in A \land x \notin B}}
  \proofstepwide{} {\leftrightarrow}{a\in U\land a\notin X\cup Y}{%
    \FormulaRefAuto{z \in A \cup B \eqvdash z \in A \lor z \in B};
    \FormulaRefAuto{\neg(P\lor Q)\eqvdash\neg P\land\neg Q}}
  \proofstepwide{} {\leftrightarrow}{a\in U\setminus(X\cup Y)}{%
    \FormulaRefAuto{x \in A \setminus B \eqvdash x \in A \land x \notin B}{2}}
  \proofstepwidestar{(U\setminus X)\cap(U\setminus Y)=U\setminus(X\cup Y)}{%
    \FormulaRefAuto{\forall x\,(x \in A \leftrightarrow x \in B) \eqvdash A = B}{\rUI{3}}}
\end{tabproofwide}

\FormulaThmDeltaK[Differenz mit der leeren Menge]{%
  A\setminus\varnothing=A
}{RelativeComplementEmpty}{%
  \DeltaRow{Mengen}{A\dsep x}
}
\begin{tabproofwide}
  \proofstepwidestar[1]{x\in A\setminus\varnothing}{\rA}
  \proofstepwidestar[1]{x\in A}{%
    \FormulaRefAuto{x\in A\setminus B\vdash x\in A}[thm]{1}}
  \proofstepwidestar[]{x\in A\setminus\varnothing\to x\in A}{\rRI{1,2}}
  \proofstepwidestar[4]{x\in A}{\rA}
  \proofstepwidestar[]{x\notin\varnothing}{%
    \FormulaRefAuto{\varnothing \coloneq
      \iota O\bigl(\forall x\,(x \not\in O)\bigr)}[def]}
  \proofstepwidestar[4]{x\in A\setminus\varnothing}{%
    \FormulaRefAuto{x\in A\dsep x\notin B\vdash x\in A\setminus B}[thm]{4,5}}
  \proofstepwidestar[]{x\in A\to x\in A\setminus\varnothing}{\rRI{4,6}}
  \proofstepwidestar[]{x\in A\setminus\varnothing\leftrightarrow x\in A}{%
    \rLRI{3,7}}
  \proofstepwidestar[]{A\setminus\varnothing=A}{%
    \FormulaRefAuto{\forall x\,(x\in A\leftrightarrow x\in B)
      \eqvdash A=B}[thm]{\rUI{8}}}
\end{tabproofwide}

\FormulaThmDeltaK[Differenz einer Menge mit sich selbst]{%
  A\setminus A=\varnothing
}{RelativeComplementSelf}{%
  \DeltaRow{Mengen}{A\dsep x}
}
\begin{tabproofwide}
  \proofstepwidestar[1]{x\in A\setminus A}{\rA}
  \proofstepwidestar[1]{x\in A\land x\notin A}{%
    \FormulaRefAuto{x\in A\setminus B\eqvdash
      x\in A\land x\notin B}[thm]{1}}
  \proofstepwidestar[1]{\bot}{\rBI{\rAEa{2},\rAEb{2}}}
  \proofstepwidestar[1]{x\in\varnothing}{%
    \FormulaRefAuto{\bot\vdash P}[thm]{3}}
  \proofstepwidestar[]{x\in A\setminus A\to x\in\varnothing}{\rRI{1,4}}
  \proofstepwidestar[6]{x\in\varnothing}{\rA}
  \proofstepwidestar[]{x\notin\varnothing}{%
    \FormulaRefAuto{\varnothing \coloneq
      \iota O\bigl(\forall x\,(x \not\in O)\bigr)}[def]}
  \proofstepwidestar[6]{\bot}{\rBI{6,7}}
  \proofstepwidestar[6]{x\in A\setminus A}{%
    \FormulaRefAuto{\bot\vdash P}[thm]{8}}
  \proofstepwidestar[]{x\in\varnothing\to x\in A\setminus A}{\rRI{6,9}}
  \proofstepwidestar[]{x\in A\setminus A\leftrightarrow x\in\varnothing}{%
    \rLRI{5,10}}
  \proofstepwidestar[]{A\setminus A=\varnothing}{%
    \FormulaRefAuto{\forall x\,(x\in A\leftrightarrow x\in B)
      \eqvdash A=B}[thm]{\rUI{11}}}
\end{tabproofwide}

\FormulaDefDeltaK[Komplementabbildung der dualen Konstruktion]{%
  \FranklFam M\vdash
  \begin{aligned}[t]
    &\kappa_M\colon M^\circ\to\mathcal L_M,\\[-2pt]
    &\forall X\in M^\circ\,
      \bigl(\kappa_M(X)\coloneqq U_M\setminus X\bigr)
  \end{aligned}
}{FranklComplementMap}{%
  \DeltaRow{Mengen}{M\dsep X}
  \DeltaRow{Funktionen}{\kappa_M}
  \DeltaRow{\textbf{Neues Symbol}}{}
  \DeltaPrem{Komplementabbildung}{\kappa_M}
}
\begin{tabproofwide}
  \proofstepwidestar[1]{\FranklFam M}{\rA}
  \proofstepwidestar[2]{X\in M^\circ}{\rA}
  \proofstepwidestar[1,2]{U_M\setminus X\in\mathcal L_M}{%
    \rAEb{\rAI{1,
      \FormulaRefAuto{FranklComplementConstruction}[def]{2,\rII}}}}
  \proofstepwidestar[1]{%
    \forall X\in M^\circ\,U_M\setminus X\in\mathcal L_M}{%
    \rUI{\rRI{2,3}}}
\end{tabproofwide}

Die letzte Beweiszeile ist die universelle Typisierung, auf die
\FormulaRefAuto{TypedFunctionDefinitionPrinciple}[thm] angewandt wird. Der
Satz rechtfertigt die unmittelbare Einführung von \(\kappa_M\) samt Typ und
Grundgleichung.

\FormulaThmDeltaK[Die Komplementabbildung ist wohldefiniert]{%
  \FranklFam M\vdash \kappa_M\colon M^\circ\to\mathcal L_M
}{FranklComplementMapFunction}{%
  \DeltaRow{Mengen}{M\dsep X\dsep Y\dsep Z}
  \DeltaRow{Funktionen}{\kappa_M}
}
\begin{tabproofwide}
  \proofstepwidestar[1]{\FranklFam M}{\rA}
  \proofstepwidestar[1]{\kappa_M\colon M^\circ\to\mathcal L_M}{%
    \FormulaRefAuto{FranklComplementMap}[def]{1}}
\end{tabproofwide}

\FormulaThmDeltaK[Die Komplementabbildung ist injektiv]{%
  \FranklFam M\vdash
  \forall X,Y\in M^\circ\,
    (\kappa_M(X)=\kappa_M(Y)\to X=Y)
}{FranklComplementMapInjective}{%
  \DeltaRow{Mengen}{M\dsep X\dsep Y}
  \DeltaRow{Funktionen}{\kappa_M}
}
\begin{tabproofwide}
  \proofstepwidestar[1]{\FranklFam M}{\rA}
  \proofstepwidestar[2]{X,Y\in M^\circ\land\kappa_M(X)=\kappa_M(Y)}{\rA}
  \proofstepwidestar[1,2]{X\subseteq U_M}{%
    \FormulaRefAuto{FranklAugmentedMemberSubsetCarrier}[thm]{1,\rAEa{2}}}
  \proofstepwidestar[1,2]{Y\subseteq U_M}{%
    \FormulaRefAuto{FranklAugmentedMemberSubsetCarrier}[thm]{1,\rAEa{\rAEb{2}}}}
  \proofstepwidestar[1]{%
    \forall Z\in M^\circ\,\kappa_M(Z)=U_M\setminus Z}{%
    \FormulaRefAuto{FranklComplementMap}[def]{1}}
  \proofstepwidestar[1,2]{\kappa_M(X)=U_M\setminus X}{%
    \rRE{\rUE{5},\rAEa{2}}}
  \proofstepwidestar[1,2]{\kappa_M(Y)=U_M\setminus Y}{%
    \rRE{\rUE{5},\rAEa{\rAEb{2}}}}
  \proofstepwidestar[1,2]{\kappa_M(X)=U_M\setminus Y}{%
    \rIE{7,\rAEb{\rAEb{2}}}}
  \proofstepwidestar[1,2]{U_M\setminus X=U_M\setminus Y}{\rIE{6,8}}
  \proofstepwidestar[1,2]{U_M\setminus Y=U_M\setminus X}{%
    \FormulaRefAuto{a=b\vdash b=a}{9}}
  \proofstepwidestar[1,2]{X=Y}{%
    \FormulaRefAuto{RelCompEqualityCancel}[thm]{3,4,10}}
  \proofstepwidestar[1]{\forall X,Y\in M^\circ
    (\kappa_M(X)=\kappa_M(Y)\to X=Y)}{%
    \rUI{\rUI{\rRI{2,11}}}}
\end{tabproofwide}

\FormulaThmDeltaK[Die Komplementabbildung ist surjektiv]{%
  \FranklFam M\vdash
  \forall Z\in\mathcal L_M\,\exists X\in M^\circ\,(\kappa_M(X)=Z)
}{FranklComplementMapSurjective}{%
  \DeltaRow{Mengen}{M\dsep X\dsep Z}
  \DeltaRow{Funktionen}{\kappa_M}
}
\begin{tabproofwide}
  \proofstepwidestar[1]{\FranklFam M}{\rA}
  \proofstepwidestar[2]{Z\in\mathcal L_M}{\rA}
  \proofstepwidestar[2]{\exists X\in M^\circ\,(Z=U_M\setminus X)}{%
    \FormulaRefAuto{FranklComplementConstruction}[def]{2}}
  \proofstepwidestar[4]{X\in M^\circ\land Z=U_M\setminus X}{\rA}
  \proofstepwidestar[1]{%
    \forall Y\in M^\circ\,\kappa_M(Y)=U_M\setminus Y}{%
    \FormulaRefAuto{FranklComplementMap}[def]{1}}
  \proofstepwidestar[1,4]{\kappa_M(X)=U_M\setminus X}{%
    \rRE{\rUE{5},\rAEa{4}}}
  \proofstepwidestar[4]{U_M\setminus X=Z}{%
    \FormulaRefAuto{a=b\vdash b=a}{\rAEb{4}}}
  \proofstepwidestar[1,4]{\kappa_M(X)=Z}{%
    \FormulaRefAuto{a=b,\,b=c\vdash a=c}{6,7}}
  \proofstepwidestar[1,2]{\exists X\in M^\circ\,(\kappa_M(X)=Z)}{%
    \rEE{3,4,\rEI{8}}}
  \proofstepwidestar[1]{\forall Z\in\mathcal L_M\exists X\in M^\circ
    (\kappa_M(X)=Z)}{\rUI{\rRI{2,9}}}
\end{tabproofwide}

\FormulaThmDeltaK[Die Komplementabbildung ist bijektiv]{%
  \FranklFam M\vdash \kappa_M\colon M^\circ\bij\mathcal L_M
}{FranklComplementMapBijective}{%
  \DeltaRow{Mengen}{M\dsep X\dsep Y\dsep Z}
  \DeltaRow{Funktionen}{\kappa_M}
}
\begin{tabproofwide}
  \proofstepwidestar[1]{\FranklFam M}{\rA}
  \proofstepwidestar[1]{\kappa_M\colon M^\circ\to\mathcal L_M}{%
    \FormulaRefAuto{FranklComplementMapFunction}[thm]{1}}
  \proofstepwidestar[1]{\forall X,Y\in M^\circ
    (\kappa_M(X)=\kappa_M(Y)\to X=Y)}{%
    \FormulaRefAuto{FranklComplementMapInjective}[thm]{1}}
  \proofstepwidestar[1]{\forall Z\in\mathcal L_M\exists X\in M^\circ
    (\kappa_M(X)=Z)}{%
    \FormulaRefAuto{FranklComplementMapSurjective}[thm]{1}}
  \proofstepwidestar[1]{\kappa_M\colon M^\circ\bij\mathcal L_M}{%
    \FormulaRefAuto{Bijektive Funktion}[def]{2,3,4}}
\end{tabproofwide}

\FormulaThmDeltaK[Mitglieder der erweiterten Familie liegen im Träger]{%
  \FranklFam M\dsep X\in M^\circ\vdash X\subseteq U_M
}{FranklAugmentedMemberSubsetCarrier}{%
  \DeltaRow{Mengen}{M\dsep X}
}
\begin{tabproof}
  \proofstep{1}{\FranklFam M}{\rA}
  \proofstep{2}{X\in M^\circ}{\rA}
  \proofstep{1}{M^\circ=M\cup\{\varnothing\}}{%
    \FormulaRefAuto{FranklEmptyAdjunction}[def]{1}}
  \proofstep{2}{X\in M\lor X=\varnothing}{%
    \FormulaRefAuto{x\in A\cup\{a\} \vdash x\in A\lor x=a}{\rIE{3,2}}}
  \proofstep{5}{X\in M}{\rA}
  \proofstep{5}{X\subseteq U_M}{\FormulaRefAuto{FamilyMemberSubsetUnion}[thm]{5}}
  \proofstep{7}{X=\varnothing}{\rA}
  \proofstep{7}{X\subseteq U_M}{%
    \rIE{7,\FormulaRefAuto{\varnothing\subseteq A}}}
  \proofstep{1,2}{X\subseteq U_M}{\rOE{4,5,6,7,8}}
\end{tabproof}

\FormulaThmDeltaK[Mitglieder der Komplementfamilie liegen im Träger]{%
  \FranklFam M\dsep X\in\mathcal L_M\vdash X\subseteq U_M
}{FranklComplementMemberSubsetCarrier}{%
  \DeltaRow{Mengen}{M\dsep X\dsep Y}
}
\begin{tabproof}
  \proofstep{1}{\FranklFam M}{\rA}
  \proofstep{2}{X\in\mathcal L_M}{\rA}
  \proofstep{2}{\exists Y\in M^\circ\,(X=U_M\setminus Y)}{%
    \FormulaRefAuto{FranklComplementConstruction}[def]{2}}
  \proofstep{4}{Y\in M^\circ\land X=U_M\setminus Y}{\rA}
  \proofstep{4}{X\subseteq U_M}{%
    \rIE{\rAEb{4},\FormulaRefAuto{A\setminus B\subseteq A}}}
  \proofstep{2}{X\subseteq U_M}{\rEE{3,4,5}}
\end{tabproof}

\FormulaThmDeltaK[Die erweiterte Familie ist endlich]{%
  \FranklFam M\vdash\Finite{M^\circ}
}{FranklAugmentedFamilyFinite}{%
  \DeltaRow{Mengen}{M}
}
\begin{tabproof}
  \proofstep{1}{\FranklFam M}{\rA}
  \proofstep{1}{\Finite M}{\FormulaRefAuto{FranklFamFinite}[thm]{1}}
  \proofstep{1}{\Finite{M\cup\{\varnothing\}}}{%
    \FormulaRefAuto{FiniteInsert}[thm]{2}}
  \proofstep{1}{M^\circ=M\cup\{\varnothing\}}{%
    \FormulaRefAuto{FranklEmptyAdjunction}[def]{1}}
  \proofstep{1}{\Finite{M^\circ}}{\rIE{4,3}}
\end{tabproof}

\FormulaThmDeltaK[Monotonie der vermeidenden Teilfamilie]{%
  A\subseteq B\vdash
  \FamAvoiding{A}{a}\subseteq\FamAvoiding{B}{a}
}{FamAvoidingMonotone}{%
  \DeltaRow{Mengen}{A\dsep B\dsep X\dsep a}
}
\begin{tabproof}
  \proofstep{1}{A\subseteq B}{\rA}
  \proofstep{2}{X\in\FamAvoiding{A}{a}}{\rA}
  \proofstep{2}{X\in A\land a\notin X}{%
    \FormulaRefAuto{\FamAvoiding{M}{a}\coloneqq
      \{X\in M\mid a\notin X\}}[def]{2}}
  \proofstep{1,2}{X\in B}{%
    \FormulaRefAuto{A\subseteq B\dsep x\in A\vdash x\in B}[thm]{%
      1,\rAEa{3}}}
  \proofstep{1,2}{X\in\FamAvoiding{B}{a}}{%
    \FormulaRefAuto{\FamAvoiding{M}{a}\coloneqq
      \{X\in M\mid a\notin X\}}[def]{\rAI{4,\rAEb{3}}}}
  \proofstep{1}{\FamAvoiding{A}{a}\subseteq\FamAvoiding{B}{a}}{%
    \FormulaRefAuto{A\subseteq B\coloneq
      \forall x\,(x\in A\rightarrow x\in B)}[def]{\rUI{\rRI{2,5}}}}
\end{tabproof}

\FormulaThmDeltaK[Adjunktion der leeren Menge erhält die enthaltende Teilfamilie]{%
  \FranklFam M\vdash
  \FamContaining{M^\circ}{a}=\FamContaining{M}{a}
}{FranklAugmentedContainingUnchanged}{%
  \DeltaRow{Mengen}{M\dsep X\dsep a}
}
\begin{tabproofsplitwide}
  \proofpartwideindR[Teilmengenrichtung zur ursprünglichen Familie]{%
    \FranklFam M\vdash
    \FamContaining{M^\circ}{a}\subseteq\FamContaining{M}{a}
  }{%
    \FranklFam M\vdash
    \FamContaining{M^\circ}{a}\subseteq\FamContaining{M}{a}
  }[FranklAugmentedContainingSubsetOriginal]
    \proofstepwidestar[1]{\FranklFam M}{\rA}
    \proofstepwidestar[2]{X\in\FamContaining{M^\circ}{a}}{\rA}
    \proofstepwidestar[2]{X\in M^\circ\land a\in X}{%
      \FormulaRefAuto{\FamContaining{M}{a}\coloneqq
        \{X\in M\mid a\in X\}}[def]{2}}
    \proofstepwidestar[1]{M^\circ=M\cup\{\varnothing\}}{%
      \FormulaRefAuto{FranklEmptyAdjunction}[def]{1}}
    \proofstepwidestar[1,2]{X\in M\cup\{\varnothing\}}{%
      \rIE{4,\rAEa{3}}}
    \proofstepwidestar[1,2]{X\in M\lor X=\varnothing}{%
      \FormulaRefAuto{x\in A\cup\{a\}\vdash x\in A\lor x=a}[thm]{5}}
    \proofstepwidestar[7]{X\in M}{\rA}
    \proofstepwidestar[2,7]{X\in\FamContaining{M}{a}}{%
      \FormulaRefAuto{\FamContaining{M}{a}\coloneqq
        \{X\in M\mid a\in X\}}[def]{\rAI{7,\rAEb{3}}}}
    \proofstepwidestar[9]{X=\varnothing}{\rA}
    \proofstepwidestar[2,9]{a\in\varnothing}{\rIE{9,\rAEb{3}}}
    \proofstepwidestar[]{a\notin\varnothing}{%
      \FormulaRefAuto{\varnothing \coloneq
        \iota O\bigl(\forall x\,(x \not\in O)\bigr)}[def]}
    \proofstepwidestar[2,9]{\bot}{\rBI{10,11}}
    \proofstepwidestar[2,9]{X\in\FamContaining{M}{a}}{%
      \FormulaRefAuto{\bot\vdash P}[thm]{12}}
    \proofstepwidestar[1,2]{X\in\FamContaining{M}{a}}{%
      \rOE{6,7,8,9,13}}
    \proofstepwidestar[1]{%
      \FamContaining{M^\circ}{a}\subseteq\FamContaining{M}{a}}{%
      \FormulaRefAuto{A\subseteq B\coloneq
        \forall x\,(x\in A\rightarrow x\in B)}[def]{\rUI{\rRI{2,14}}}}
  \closeproofpartwideindR

  \proofpartwideindR[Teilmengenrichtung zur erweiterten Familie]{%
    \FranklFam M\vdash
    \FamContaining{M}{a}\subseteq\FamContaining{M^\circ}{a}
  }{%
    \FranklFam M\vdash
    \FamContaining{M}{a}\subseteq\FamContaining{M^\circ}{a}
  }[FranklAugmentedContainingSubsetAugmented]
    \proofstepwidestar[1]{\FranklFam M}{\rA}
    \proofstepwidestar[2]{X\in\FamContaining{M}{a}}{\rA}
    \proofstepwidestar[2]{X\in M\land a\in X}{%
      \FormulaRefAuto{\FamContaining{M}{a}\coloneqq
        \{X\in M\mid a\in X\}}[def]{2}}
    \proofstepwidestar[1]{M^\circ=M\cup\{\varnothing\}}{%
      \FormulaRefAuto{FranklEmptyAdjunction}[def]{1}}
    \proofstepwidestar[2]{X\in M\cup\{\varnothing\}}{%
      \FormulaRefAuto{z\in A\vdash z\in A\cup B}[thm]{\rAEa{3}}}
    \proofstepwidestar[1,2]{X\in M^\circ}{\rIE{4,5}}
    \proofstepwidestar[1,2]{X\in\FamContaining{M^\circ}{a}}{%
      \FormulaRefAuto{\FamContaining{M}{a}\coloneqq
        \{X\in M\mid a\in X\}}[def]{\rAI{6,\rAEb{3}}}}
    \proofstepwidestar[1]{%
      \FamContaining{M}{a}\subseteq\FamContaining{M^\circ}{a}}{%
      \FormulaRefAuto{A\subseteq B\coloneq
        \forall x\,(x\in A\rightarrow x\in B)}[def]{\rUI{\rRI{2,7}}}}
  \closeproofpartwideindR

  \proofpartwide{Schluss}
    \proofstepwidestar[1]{%
      \FamContaining{M^\circ}{a}\subseteq\FamContaining{M}{a}}{%
      \FormulaRefAuto{FranklAugmentedContainingSubsetOriginal}[thm]{1}}
    \proofstepwidestar[1]{%
      \FamContaining{M}{a}\subseteq\FamContaining{M^\circ}{a}}{%
      \FormulaRefAuto{FranklAugmentedContainingSubsetAugmented}[thm]{1}}
    \proofstepwidestar[1]{%
      \FamContaining{M^\circ}{a}=\FamContaining{M}{a}}{%
      \FormulaRefAuto{A\subseteq B\dsep B\subseteq A\vdash A=B}[thm]{1,2}}
  \closeproofpartwide
\end{tabproofsplitwide}

\FormulaThmDeltaK[Die erweiterte Familie ist vereinigungsabgeschlossen]{%
  \FranklFam M\vdash\UnionClosedFam(M^\circ)
}{FranklAugmentedFamilyUnionClosed}{%
  \DeltaRow{Mengen}{M\dsep X\dsep Y}
}
\begin{tabproofwide}
  \proofstepwidestar[1]{\FranklFam M}{\rA}
  \proofstepwidestar[1]{\UnionClosedFam(M)}{%
    \FormulaRefAuto{FranklFamUnionClosed}[thm]{1}}
  \proofstepwidestar[]{M^\circ=M\cup\{\varnothing\}}{%
    \FormulaRefAuto{FranklEmptyAdjunction}[def]}
  \proofstepwidestar[4]{X,Y\in M^\circ}{\rA}
  \proofstepwidestar[4]{X\in M\lor X=\varnothing}{%
    \FormulaRefAuto{x\in A\cup\{a\} \vdash x\in A\lor x=a}{%
      \rIE{3,\rAEa{4}}}}
  \proofstepwidestar[4]{Y\in M\lor Y=\varnothing}{%
    \FormulaRefAuto{x\in A\cup\{a\} \vdash x\in A\lor x=a}{%
      \rIE{3,\rAEb{4}}}}

  \proofstepwidestar[7]{X\in M}{\rA}
  \proofstepwidestar[8]{Y\in M}{\rA}
  \proofstepwidestar[2,7,8]{X\cup Y\in M}{%
    \FormulaRefAuto{X\in M\dsep Y\in M\vdash X\cup Y\in M}[ax]{7,8}}
  \proofstepwidestar[1,7,8]{X\cup Y\in M\cup\{\varnothing\}}{%
    \FormulaRefAuto{z\in A\vdash z\in A\cup B}[thm]{9}}
  \proofstepwidestar[1,7,8]{X\cup Y\in M^\circ}{\rIE{3,10}}
  \proofstepwidestar[12]{Y=\varnothing}{\rA}
  \proofstepwidestar[12]{X\cup Y=X}{%
    \rIE{12,\FormulaRefAuto{A\cup\varnothing=A}}}
  \proofstepwidestar[7]{X\in M\cup\{\varnothing\}}{%
    \FormulaRefAuto{z\in A\vdash z\in A\cup B}[thm]{7}}
  \proofstepwidestar[7]{X\in M^\circ}{\rIE{3,14}}
  \proofstepwidestar[7,12]{X\cup Y\in M^\circ}{\rIE{13,15}}
  \proofstepwidestar[1,4,6,7]{X\cup Y\in M^\circ}{%
    \rOE{6,8,11,12,16}}

  \proofstepwidestar[18]{X=\varnothing}{\rA}
  \proofstepwidestar[19]{Y\in M}{\rA}
  \proofstepwidestar[18]{X\cup Y=Y}{%
    \rIE{18,\FormulaRefAuto{\varnothing\cup A=A}}}
  \proofstepwidestar[19]{Y\in M\cup\{\varnothing\}}{%
    \FormulaRefAuto{z\in A\vdash z\in A\cup B}[thm]{19}}
  \proofstepwidestar[19]{Y\in M^\circ}{\rIE{3,21}}
  \proofstepwidestar[18,19]{X\cup Y\in M^\circ}{\rIE{20,22}}
  \proofstepwidestar[24]{Y=\varnothing}{\rA}
  \proofstepwidestar[18,24]{X\cup Y=\varnothing}{%
    \rIE{18,\rIE{24,\FormulaRefAuto{\varnothing\cup\varnothing=\varnothing}}}}
  \proofstepwidestar[]{\varnothing\in M\cup\{\varnothing\}}{%
    \FormulaRefAuto{a\in A\cup\{a\}}[thm]}
  \proofstepwidestar[]{\varnothing\in M^\circ}{\rIE{3,26}}
  \proofstepwidestar[18,24]{X\cup Y\in M^\circ}{\rIE{25,27}}
  \proofstepwidestar[4,6,18]{X\cup Y\in M^\circ}{%
    \rOE{6,19,23,24,28}}
  \proofstepwidestar[1,4]{X\cup Y\in M^\circ}{%
    \rOE{5,7,17,18,29}}
  \proofstepwidestar[1]{\UnionClosedFam(M^\circ)}{%
    \FormulaRefAuto{Vereinigungsabgeschlossene Familie}[def]{%
      \rUI{\rUI{\rRI{4,30}}}}}
\end{tabproofwide}

\FormulaThmDeltaK[Die komplementduale Familie ist endlich]{%
  \FranklFam M\vdash\Finite{\mathcal L_M}
}{FranklComplementFamilyFinite}{%
  \DeltaRow{Mengen}{M}
  \DeltaRow{Funktionen}{\kappa_M}
}
\begin{tabproof}
  \proofstep{1}{\FranklFam M}{\rA}
  \proofstep{1}{\Finite{M^\circ}}{%
    \FormulaRefAuto{FranklAugmentedFamilyFinite}[thm]{1}}
  \proofstep{1}{\kappa_M\colon M^\circ\bij\mathcal L_M}{%
    \FormulaRefAuto{FranklComplementMapBijective}[thm]{1}}
  \proofstep{1}{\Finite{\mathcal L_M}}{%
    \FormulaRefAuto{FiniteBijectionTransport}[thm]{2,3}}
\end{tabproof}

\FormulaThmDeltaK[Die komplementduale Familie ist schnittabgeschlossen]{%
  \FranklFam M\vdash\InterClosedFam{\mathcal L_M}
}{FranklComplementFamilyIntersectionClosed}{%
  \DeltaRow{Mengen}{M\dsep X\dsep Y\dsep A\dsep B}
}
\begin{tabproofwide}
  \proofstepwidestar[1]{\FranklFam M}{\rA}
  \proofstepwidestar[2]{X,Y\in\mathcal L_M}{\rA}
  \proofstepwidestar[1]{\UnionClosedFam(M^\circ)}{%
    \FormulaRefAuto{FranklAugmentedFamilyUnionClosed}[thm]{1}}
  \proofstepwidestar[2]{\exists A\in M^\circ\,(X=U_M\setminus A)}{%
    \FormulaRefAuto{FranklComplementConstruction}[def]{\rAEa{2}}}
  \proofstepwidestar[2]{\exists B\in M^\circ\,(Y=U_M\setminus B)}{%
    \FormulaRefAuto{FranklComplementConstruction}[def]{\rAEb{2}}}
  \proofstepwidestar[6]{A\in M^\circ\land X=U_M\setminus A}{\rA}
  \proofstepwidestar[7]{B\in M^\circ\land Y=U_M\setminus B}{\rA}
  \proofstepwidestar[1,6,7]{A\cup B\in M^\circ}{%
    \FormulaRefAuto{X\in M\dsep Y\in M\vdash X\cup Y\in M}[ax]{%
      \rAEa{6},\rAEa{7}}}
  \proofstepwidestar[6,7]{X\cap Y=U_M\setminus(A\cup B)}{%
    \rIE{\rAEb{6},\rIE{\rAEb{7},
      \FormulaRefAuto{RelativeComplementUnionDeMorgan}[thm]}}}
  \proofstepwidestar[1,2,6,7]{X\cap Y\in\mathcal L_M}{%
    \FormulaRefAuto{FranklComplementConstruction}[def]{8,9}}
  \proofstepwidestar[1,2]{X\cap Y\in\mathcal L_M}{%
    \rEE{4,6,\rEE{5,7,10}}}
  \proofstepwidestar[1]{\InterClosedFam{\mathcal L_M}}{%
    \FormulaRefAuto{IntersectionClosedFamily}[def]{\rUI{\rUI{\rRI{2,11}}}}}
\end{tabproofwide}

\FormulaThmDeltaK[Die Inklusionsordnung der Komplementfamilie ist partiell]{%
  \FranklFam M\vdash\PartOrd{\mathcal L_M,\subseteq}
}{FranklComplementOrderPartial}{%
  \DeltaRow{Mengen}{M\dsep X}
  \DeltaRow{Relationsoperatoren}{\subseteq}
}
\begin{tabproof}
  \proofstep{1}{\FranklFam M}{\rA}
  \proofstep{2}{X\in\mathcal L_M}{\rA}
  \proofstep{1,2}{X\subseteq U_M}{%
    \FormulaRefAuto{FranklComplementMemberSubsetCarrier}[thm]{1,2}}
  \proofstep{1,2}{X\in\powerset(U_M)}{%
    \FormulaRefAuto{x\in\powerset(A)\eqvdash x\subseteq A}{3}}
  \proofstep{1}{\mathcal L_M\subseteq\powerset(U_M)}{%
    \FormulaRefAuto{A \subseteq B \coloneq
      \forall x\,(x\in A \rightarrow x\in B)}{\rUI{\rRI{2,4}}}}
  \proofstep{1}{\PartOrd{\mathcal L_M,\subseteq}}{%
    \FormulaRefAuto{InclusionOrderOnFamily}[thm]{5}}
\end{tabproof}

\FormulaDefDeltaK[Gemeinsame obere Elemente der Komplementordnung]{%
  \mathcal C_M(X,Y)\coloneqq
  \{Z\in\mathcal L_M\mid X\cup Y\subseteq Z\}
}{FranklComplementCommonUpperFamily}{%
  \DeltaRow{Mengen}{M\dsep X\dsep Y\dsep Z}
  \DeltaRow{\textbf{Neues Symbol}}{}
  \DeltaPrem{Gemeinsame obere Elemente}{\mathcal C_M(X,Y)}
}

\FormulaThmDeltaK[Die Familie gemeinsamer oberer Elemente ist endlich und nichtleer]{%
  \FranklFam M\dsep X,Y\in\mathcal L_M
  \vdash\Finite{\mathcal C_M(X,Y)}\land\mathcal C_M(X,Y)\neq\varnothing
}{FranklComplementCommonUpperFiniteNonempty}{%
  \DeltaRow{Mengen}{M\dsep X\dsep Y\dsep Z}
}
Im Beweis steht
\[
  \mathcal C\coloneqq\mathcal C_M(X,Y).
\]
\begin{tabproofwide}
  \proofstepwidestar[1]{\FranklFam M}{\rA}
  \proofstepwidestar[2]{X,Y\in\mathcal L_M}{\rA}
  \proofstepwidestar[1]{\Finite{\mathcal L_M}}{%
    \FormulaRefAuto{FranklComplementFamilyFinite}[thm]{1}}
  \proofstepwidestar[]{\mathcal C\subseteq\mathcal L_M}{%
    \FormulaRefAuto{FranklComplementCommonUpperFamily}[def]}
  \proofstepwidestar[1]{U_M\in M}{%
    \ProofRefStack{%
      \FormulaRefAuto{FiniteUnionClosedFamilyTotalUnion}[thm]\\
      \FormulaRefAuto{FranklFamFinite}[thm](1)\\
      \FormulaRefAuto{FranklFamNonempty}[thm](1)\\
      \FormulaRefAuto{FranklFamUnionClosed}[thm](1)}}
  \proofstepwidestar[1]{U_M\in\mathcal L_M}{%
    \ProofRefStack{%
      \FormulaRefAuto{FranklComplementConstruction}[def]\\
      \FormulaRefAuto{a\in A\cup\{a\}}[thm]\\
      \FormulaRefAuto{RelativeComplementEmpty}[thm]}}
  \proofstepwidestar[1,2]{X\cup Y\subseteq U_M}{%
    \ProofRefStack{%
      \FormulaRefAuto{A \subseteq C,\, B \subseteq C
        \vdash A \cup B \subseteq C}\\
      \FormulaRefAuto{FranklComplementMemberSubsetCarrier}[thm]
        \bigl(1,\rAEa{2}\bigr)\\
      \FormulaRefAuto{FranklComplementMemberSubsetCarrier}[thm]
        \bigl(1,\rAEb{2}\bigr)}}
  \proofstepwidestar[1,2]{U_M\in\mathcal C}{%
    \FormulaRefAuto{FranklComplementCommonUpperFamily}[def]{6,7}}
  \proofstepwidestar[1,2]{\mathcal C\neq\varnothing}{%
    \FormulaRefAuto{a\in S\vdash S\neq\varnothing}{8}}
  \proofstepwidestar[1]{\Finite{\mathcal C}}{%
    \FormulaRefAuto{FiniteSubset}[thm]{4,3}}
  \proofstepwidestar[1,2]{\Finite{\mathcal C}\land
    \mathcal C\neq\varnothing}{\rAI{10,9}}
\end{tabproofwide}

\FormulaThmDeltaK[Die konstruierte Operation ist das Paarsupremum]{%
  \FranklFam M\dsep X,Y\in\mathcal L_M\vdash
  \PairSup{\mathcal L_M}{\subseteq}{X}{Y}{X\FamilyJoinOp{M}{}Y}
}{FranklComplementPairSupremum}{%
  \DeltaRow{Mengen}{M\dsep X\dsep Y\dsep Z\dsep W}
  \DeltaRow{Relationsoperatoren}{\subseteq}
  \DeltaRow{Funktionen}{\FamilyJoinOp{M}}
}
\begin{tabproofsplitwide}
  \proofpartwideindR[Beide Operanden liegen im Schnitt der gemeinsamen oberen Elemente]{%
    \begin{aligned}[t]
      &\FranklFam M\dsep X,Y\in\mathcal L_M\vdash{}\\
      &X\subseteq\bigcap\mathcal C_M(X,Y)\land
       Y\subseteq\bigcap\mathcal C_M(X,Y)
    \end{aligned}
  }{%
    \FranklFam M\dsep X,Y\in\mathcal L_M\vdash
    X\subseteq\bigcap\mathcal C_M(X,Y)\land
    Y\subseteq\bigcap\mathcal C_M(X,Y)
  }[FranklComplementCommonUpperIntersectionBounds]
    \proofstepwidestar[1]{\FranklFam M}{\rA}
    \proofstepwidestar[2]{X,Y\in\mathcal L_M}{\rA}
    \proofstepwidestar[1,2]{\mathcal C_M(X,Y)\neq\varnothing}{%
      \rAEb{\FormulaRefAuto{FranklComplementCommonUpperFiniteNonempty}[thm]{1,2}}}
    \proofstepwidestar[4]{a\in X}{\rA}
    \proofstepwidestar[5]{Z\in\mathcal C_M(X,Y)}{\rA}
    \proofstepwidestar[5]{X\cup Y\subseteq Z}{%
      \FormulaRefAuto{FranklComplementCommonUpperFamily}[def]{5}}
    \proofstepwidestar[]{X\subseteq X\cup Y}{%
      \FormulaRefAuto{A\subseteq A\cup B}}
    \proofstepwidestar[5]{X\subseteq Z}{%
      \FormulaRefAuto{A\subseteq B\dsep B\subseteq C\vdash A\subseteq C}{7,6}}
    \proofstepwidestar[4,5]{a\in Z}{%
      \FormulaRefAuto{A\subseteq B,\,x\in A\vdash x\in B}{8,4}}
    \proofstepwidestar[4]{\forall Z\in\mathcal C_M(X,Y)\,(a\in Z)}{%
      \rUI{\rRI{5,9}}}
    \proofstepwidestar[1,2]{%
      a\in\bigcap\mathcal C_M(X,Y)\leftrightarrow
      \forall Z\in\mathcal C_M(X,Y)\,(a\in Z)}{%
      \FormulaRefAuto{M\neq\varnothing \vdash x\in\bigcap M \leftrightarrow
        \forall X\in M\,(x\in X)}{3}}
    \proofstepwidestar[1,2,4]{a\in\bigcap\mathcal C_M(X,Y)}{%
      \FormulaRefAuto{P\leftrightarrow Q,Q\vdash P}{11,10}}
    \proofstepwidestar[1,2]{X\subseteq\bigcap\mathcal C_M(X,Y)}{%
      \FormulaRefAuto{A \subseteq B \coloneq
        \forall x\,(x\in A \rightarrow x\in B)}{\rUI{\rRI{4,12}}}}
    \proofstepwidestar[14]{a\in Y}{\rA}
    \proofstepwidestar[15]{Z\in\mathcal C_M(X,Y)}{\rA}
    \proofstepwidestar[15]{X\cup Y\subseteq Z}{%
      \FormulaRefAuto{FranklComplementCommonUpperFamily}[def]{15}}
    \proofstepwidestar[]{Y\subseteq X\cup Y}{%
      \FormulaRefAuto{B\subseteq A\cup B}}
    \proofstepwidestar[15]{Y\subseteq Z}{%
      \FormulaRefAuto{A\subseteq B\dsep B\subseteq C\vdash A\subseteq C}{17,16}}
    \proofstepwidestar[14,15]{a\in Z}{%
      \FormulaRefAuto{A\subseteq B,\,x\in A\vdash x\in B}{18,14}}
    \proofstepwidestar[14]{\forall Z\in\mathcal C_M(X,Y)\,(a\in Z)}{%
      \rUI{\rRI{15,19}}}
    \proofstepwidestar[1,2]{%
      a\in\bigcap\mathcal C_M(X,Y)\leftrightarrow
      \forall Z\in\mathcal C_M(X,Y)\,(a\in Z)}{%
      \FormulaRefAuto{M\neq\varnothing \vdash x\in\bigcap M \leftrightarrow
        \forall X\in M\,(x\in X)}{3}}
    \proofstepwidestar[1,2,14]{a\in\bigcap\mathcal C_M(X,Y)}{%
      \FormulaRefAuto{P\leftrightarrow Q,Q\vdash P}{21,20}}
    \proofstepwidestar[1,2]{Y\subseteq\bigcap\mathcal C_M(X,Y)}{%
      \FormulaRefAuto{A \subseteq B \coloneq
        \forall x\,(x\in A \rightarrow x\in B)}{\rUI{\rRI{14,22}}}}
    \proofstepwidestar[1,2]{%
      X\subseteq\bigcap\mathcal C_M(X,Y)\land
      Y\subseteq\bigcap\mathcal C_M(X,Y)}{\rAI{13,23}}
  \closeproofpartwideindR

  \proofpartwideindR[Träger und obere Schranken]{%
    \begin{aligned}[t]
      &\FranklFam M\dsep X,Y\in\mathcal L_M\vdash{}\\
      &X\FamilyJoinOp{M}{}Y\in\mathcal L_M\land
       X\subseteq X\FamilyJoinOp{M}{}Y\land
       Y\subseteq X\FamilyJoinOp{M}{}Y
    \end{aligned}
  }{%
    \FranklFam M\dsep X,Y\in\mathcal L_M\vdash
    X\FamilyJoinOp{M}{}Y\in\mathcal L_M\land
    X\subseteq X\FamilyJoinOp{M}{}Y\land
    Y\subseteq X\FamilyJoinOp{M}{}Y
  }[FranklComplementPairSupremumBounds]
    \proofstepwidestar[1]{\FranklFam M}{\rA}
    \proofstepwidestar[2]{X,Y\in\mathcal L_M}{\rA}
    \proofstepwidestar[1]{\InterClosedFam{\mathcal L_M}}{%
      \FormulaRefAuto{FranklComplementFamilyIntersectionClosed}[thm]{1}}
    \proofstepwidestar[1,2]{\Finite{\mathcal C_M(X,Y)}\land
      \mathcal C_M(X,Y)\neq\varnothing}{%
      \FormulaRefAuto{FranklComplementCommonUpperFiniteNonempty}[thm]{1,2}}
    \proofstepwidestar[1,2]{\Finite{\mathcal C_M(X,Y)}}{\rAEa{4}}
    \proofstepwidestar[1,2]{\mathcal C_M(X,Y)\neq\varnothing}{\rAEb{4}}
    \proofstepwidestar[]{\mathcal C_M(X,Y)\subseteq\mathcal L_M}{%
      \FormulaRefAuto{FranklComplementCommonUpperFamily}[def]}
    \proofstepwidestar[1,2]{\bigcap\mathcal C_M(X,Y)\in\mathcal L_M}{%
      \FormulaRefAuto{FiniteIntersectionClosedFamilyIntersection}[thm]{%
        5,6,7,3}}
    \proofstepwidestar[1,2]{X\FamilyJoinOp{M}{}Y=\bigcap\mathcal C_M(X,Y)}{%
      \FormulaRefAuto{FranklComplementConstruction}[def]{1,2}}
    \proofstepwidestar[1,2]{X\FamilyJoinOp{M}{}Y\in\mathcal L_M}{\rIE{9,8}}
    \proofstepwidestar[1,2]{%
      X\subseteq\bigcap\mathcal C_M(X,Y)\land
      Y\subseteq\bigcap\mathcal C_M(X,Y)}{%
      \FormulaRefAuto{FranklComplementCommonUpperIntersectionBounds}[thm]{1,2}}
    \proofstepwidestar[1,2]{X\subseteq\bigcap\mathcal C_M(X,Y)}{\rAEa{11}}
    \proofstepwidestar[1,2]{Y\subseteq\bigcap\mathcal C_M(X,Y)}{\rAEb{11}}
    \proofstepwidestar[1,2]{X\subseteq X\FamilyJoinOp{M}{}Y}{%
      \rIE{\FormulaRefAuto{a=b\vdash b=a}{9},12}}
    \proofstepwidestar[1,2]{Y\subseteq X\FamilyJoinOp{M}{}Y}{%
      \rIE{\FormulaRefAuto{a=b\vdash b=a}{9},13}}
    \proofstepwidestar[1,2]{%
      X\FamilyJoinOp{M}{}Y\in\mathcal L_M\land
      X\subseteq X\FamilyJoinOp{M}{}Y\land
      Y\subseteq X\FamilyJoinOp{M}{}Y}{%
      \rAI{10,\rAI{14,15}}}
  \closeproofpartwideindR
  \proofpartwideindR[Kleinste obere Schranke]{%
    \begin{aligned}[t]
      &\FranklFam M\dsep X,Y\in\mathcal L_M\dsep
       W\in\mathcal L_M\dsep X\subseteq W\dsep
       Y\subseteq W\vdash{}\\
      &X\FamilyJoinOp{M}{}Y\subseteq W
    \end{aligned}
  }{%
    \FranklFam M\dsep X,Y\in\mathcal L_M\dsep
    W\in\mathcal L_M\dsep X\subseteq W\dsep
    Y\subseteq W\vdash
    X\FamilyJoinOp{M}{}Y\subseteq W
  }[FranklComplementPairSupremumLeast]
    \proofstepwidestar[1]{\FranklFam M}{\rA}
    \proofstepwidestar[2]{X,Y\in\mathcal L_M}{\rA}
    \proofstepwidestar[3]{W\in\mathcal L_M}{\rA}
    \proofstepwidestar[4]{X\subseteq W}{\rA}
    \proofstepwidestar[5]{Y\subseteq W}{\rA}
    \proofstepwidestar[4,5]{X\cup Y\subseteq W}{%
      \FormulaRefAuto{A \subseteq C,\, B \subseteq C \vdash A \cup B \subseteq C}{4,5}}
    \proofstepwidestar[3,4,5]{W\in\mathcal C_M(X,Y)}{%
      \FormulaRefAuto{FranklComplementCommonUpperFamily}[def]{3,6}}
    \proofstepwidestar[1,2]{\mathcal C_M(X,Y)\neq\varnothing}{%
      \rAEb{\FormulaRefAuto{FranklComplementCommonUpperFiniteNonempty}[thm]{1,2}}}
    \proofstepwidestar[1,2,3,4,5]{\bigcap\mathcal C_M(X,Y)\subseteq W}{%
      \FormulaRefAuto{FamilyIntersectionSubsetMember}[thm]{%
        8,7}}
    \proofstepwidestar[1,2]{X\FamilyJoinOp{M}{}Y=\bigcap\mathcal C_M(X,Y)}{%
      \FormulaRefAuto{FranklComplementConstruction}[def]{1,2}}
    \proofstepwidestar[1,2,3,4,5]{X\FamilyJoinOp{M}{}Y\subseteq W}{%
      \rIE{\FormulaRefAuto{a=b\vdash b=a}{10},9}}
  \closeproofpartwideindR
  \proofpartwideindR[Zusammenführung der Paarsupremumseigenschaften]{%
    \FranklFam M\dsep X,Y\in\mathcal L_M\vdash
    \PairSup{\mathcal L_M}{\subseteq}{X}{Y}{X\FamilyJoinOp{M}{}Y}
  }{%
    \FranklFam M\dsep X,Y\in\mathcal L_M\vdash
    \PairSup{\mathcal L_M}{\subseteq}{X}{Y}{X\FamilyJoinOp{M}{}Y}
  }[FranklComplementPairSupremumAssembly]
    \proofstepwidestar[1]{\FranklFam M}{\rA}
    \proofstepwidestar[2]{X,Y\in\mathcal L_M}{\rA}
    \proofstepwidestar[1]{\PartOrd{\mathcal L_M,\subseteq}}{%
      \FormulaRefAuto{FranklComplementOrderPartial}[thm]{1}}
    \proofstepwidestar[1,2]{X\FamilyJoinOp{M}{}Y\in\mathcal L_M}{%
      \rAEa{\FormulaRefAuto{FranklComplementPairSupremumBounds}[thm]{1,2}}}
    \proofstepwidestar[1,2]{X\subseteq X\FamilyJoinOp{M}{}Y}{%
      \rAEa{\rAEb{\FormulaRefAuto{FranklComplementPairSupremumBounds}[thm]{1,2}}}}
    \proofstepwidestar[1,2]{Y\subseteq X\FamilyJoinOp{M}{}Y}{%
      \rAEb{\rAEb{\FormulaRefAuto{FranklComplementPairSupremumBounds}[thm]{1,2}}}}
    \proofstepwidestar[2]{X\in\mathcal L_M}{\rAEa{2}}
    \proofstepwidestar[2]{Y\in\mathcal L_M}{\rAEb{2}}
    \proofstepwidestar[2]{\{X,Y\}\subseteq\mathcal L_M}{%
      \FormulaRefAuto{a \in A,\, b \in A \vdash
        \{a,b\} \subseteq A}{7,8}}
    \proofstepwidestar[1,2]{%
      X\FamilyJoinOp{M}{}Y\in\UB_{\mathcal L_M,\subseteq}(\{X,Y\})}{%
      \ProofRefStack{%
        \FormulaRefAuto{P\leftrightarrow Q, Q\vdash P}\\
        \bigl(\FormulaRefAuto{PairUpperBoundSetCharacterization}[thm]
          (3,7,8),\bigr.\\[-2pt]
        \bigl.\rAI{4,\rAI{5,6}}\bigr)}}
    \proofstepwidestar[11]{%
      W\in\UB_{\mathcal L_M,\subseteq}(\{X,Y\})}{\rA}
    \proofstepwidestar[1,2]{%
      \begin{aligned}[t]
        &W\in\UB_{\mathcal L_M,\subseteq}(\{X,Y\})\leftrightarrow{}\\[-2pt]
        &\bigl(W\in\mathcal L_M\land X\subseteq W
          \land Y\subseteq W\bigr)
      \end{aligned}}{%
      \FormulaRefAuto{PairUpperBoundSetCharacterization}[thm]{3,7,8}}
    \proofstepwidestar[1,2,11]{%
      W\in\mathcal L_M\land X\subseteq W
        \land Y\subseteq W}{%
      \FormulaRefAuto{P\leftrightarrow Q, P\vdash Q}{12,11}}
    \proofstepwidestar[1,2,11]{%
      X\FamilyJoinOp{M}{}Y\subseteq W}{%
      \ProofRefStack{%
        \FormulaRefAuto{FranklComplementPairSupremumLeast}[thm]\\
        \bigl(1,2,\rAEa{13},\rAEa{\rAEb{13}},\bigr.\\[-2pt]
        \bigl.\rAEb{\rAEb{13}}\bigr)}}
    \proofstepwidestar[1,2]{%
      \begin{aligned}
        &\forall W\in\UB_{\mathcal L_M,\subseteq}(\{X,Y\})\,{}\\[-2pt]
        &\quad X\FamilyJoinOp{M}{}Y\subseteq W
      \end{aligned}}{%
      \rUI{\rRI{11,14}}}
    \proofstepwidestar[1,2]{%
      \PairSup{\mathcal L_M}{\subseteq}{X}{Y}{X\FamilyJoinOp{M}{}Y}}{%
      \FormulaRefAuto{SupremumCandidateEqualsSupremum}[thm]{3,9,10,15}}
  \closeproofpartwideindR
\end{tabproofsplitwide}

\FormulaThmDeltaK[Existenz der Paarsuprema in der Komplementordnung]{%
  \begin{aligned}[t]
    &\FranklFam M\dsep X,Y\in\mathcal L_M\vdash{}\\[-2pt]
    &\exists Z\in\UB_{\mathcal L_M,\subseteq}(\{X,Y\})\,
      \forall W\in\UB_{\mathcal L_M,\subseteq}(\{X,Y\})\,
        Z\subseteq W
  \end{aligned}
}{FranklComplementPairSupremumExistence}{%
  \DeltaRow{Mengen}{M\dsep W\dsep X\dsep Y\dsep Z}
  \DeltaRow{Relationsoperatoren}{\subseteq}
  \DeltaRow{Funktionen}{\FamilyJoinOp{M}}
}
\begin{tabproofwide}
  \proofstepwidestar[1]{\FranklFam M}{\rA}
  \proofstepwidestar[2]{X,Y\in\mathcal L_M}{\rA}
  \proofstepwidestar[1]{\PartOrd{\mathcal L_M,\subseteq}}{%
    \FormulaRefAuto{FranklComplementOrderPartial}[thm]{1}}
  \proofstepwidestar[1,2]{%
    \begin{aligned}[t]
      &X\FamilyJoinOp{M}{}Y\in\mathcal L_M\land{}\\[-2pt]
      &X\subseteq X\FamilyJoinOp{M}{}Y\land
        Y\subseteq X\FamilyJoinOp{M}{}Y
    \end{aligned}}{%
    \FormulaRefAuto{FranklComplementPairSupremumBounds}[thm]{1,2}}
  \proofstepwidestar[1,2]{X\FamilyJoinOp{M}{}Y\in\mathcal L_M}{\rAEa{4}}
  \proofstepwidestar[1,2]{X\subseteq X\FamilyJoinOp{M}{}Y}{%
    \rAEa{\rAEb{4}}}
  \proofstepwidestar[1,2]{Y\subseteq X\FamilyJoinOp{M}{}Y}{%
    \rAEb{\rAEb{4}}}
  \proofstepwidestar[2]{X\in\mathcal L_M}{\rAEa{2}}
  \proofstepwidestar[2]{Y\in\mathcal L_M}{\rAEb{2}}
  \proofstepwidestar[1,2]{%
    \begin{aligned}[t]
      &X\FamilyJoinOp{M}{}Y\in
        \UB_{\mathcal L_M,\subseteq}(\{X,Y\})\leftrightarrow{}\\[-2pt]
      &\bigl(X\FamilyJoinOp{M}{}Y\in\mathcal L_M\land
        X\subseteq X\FamilyJoinOp{M}{}Y\land{}\\[-2pt]
      &\qquad Y\subseteq X\FamilyJoinOp{M}{}Y\bigr)
    \end{aligned}}{%
    \FormulaRefAuto{PairUpperBoundSetCharacterization}[thm]{3,8,9}}
  \proofstepwidestar[1,2]{%
    \begin{aligned}[t]
      &X\FamilyJoinOp{M}{}Y\in\mathcal L_M\land
        X\subseteq X\FamilyJoinOp{M}{}Y\land{}\\[-2pt]
      &\quad Y\subseteq X\FamilyJoinOp{M}{}Y
    \end{aligned}}{%
    \rAI{5,\rAI{6,7}}}
  \proofstepwidestar[1,2]{%
    \begin{aligned}[t]
      &X\FamilyJoinOp{M}{}Y\in{}\\[-2pt]
      &\quad\UB_{\mathcal L_M,\subseteq}(\{X,Y\})
    \end{aligned}}{%
    \ProofRefStack{\FormulaRefAuto{P\leftrightarrow Q, Q\vdash P}\\
      (10,11)}}
  \proofstepwidestar[13]{%
    W\in\UB_{\mathcal L_M,\subseteq}(\{X,Y\})}{\rA}
  \proofstepwidestar[1,2]{%
      \begin{aligned}[t]
        &W\in\UB_{\mathcal L_M,\subseteq}(\{X,Y\})\leftrightarrow{}\\[-2pt]
        &\bigl(W\in\mathcal L_M\land X\subseteq W
          \land Y\subseteq W\bigr)
      \end{aligned}}{%
    \FormulaRefAuto{PairUpperBoundSetCharacterization}[thm]{3,8,9}}
  \proofstepwidestar[1,2,13]{%
    \begin{aligned}[t]
      &W\in\mathcal L_M\land X\subseteq W\land{}\\[-2pt]
      &\quad Y\subseteq W
    \end{aligned}}{%
    \FormulaRefAuto{P\leftrightarrow Q, P\vdash Q}{14,13}}
  \proofstepwidestar[1,2,13]{W\in\mathcal L_M}{\rAEa{15}}
  \proofstepwidestar[1,2,13]{X\subseteq W}{\rAEa{\rAEb{15}}}
  \proofstepwidestar[1,2,13]{Y\subseteq W}{\rAEb{\rAEb{15}}}
  \proofstepwidestar[1,2,13]{%
    X\FamilyJoinOp{M}{}Y\subseteq W}{%
    \ProofRefStack{%
      \FormulaRefAuto{FranklComplementPairSupremumLeast}[thm]\\
      (1,2,16,17,18)}}
  \proofstepwidestar[1,2]{%
    \begin{aligned}[t]
      &\forall W\in\UB_{\mathcal L_M,\subseteq}(\{X,Y\})\,{}\\[-2pt]
      &\quad X\FamilyJoinOp{M}{}Y\subseteq W
    \end{aligned}}{%
    \rUI{\rRI{13,19}}}
  \proofstepwidestar[1,2]{%
    \begin{aligned}[t]
      &\exists Z\in\UB_{\mathcal L_M,\subseteq}(\{X,Y\})\,{}\\[-2pt]
      &\quad\forall W\in\UB_{\mathcal L_M,\subseteq}(\{X,Y\})\,
        Z\subseteq W
    \end{aligned}}{%
    \rEI{\rAI{12,20}}}
\end{tabproofwide}

\FormulaThmDeltaK[Die konstruierte Operation ist eine Paarsupremumsoperation]{%
  \FranklFam M\vdash
  \PairJoinOp{\mathcal L_M,\FamilyJoinOp{M}}
}{FranklComplementPairJoinOperation}{%
  \DeltaRow{Mengen}{M\dsep X\dsep Y}
  \DeltaRow{Relationsoperatoren}{\subseteq}
  \DeltaRow{Funktionen}{\FamilyJoinOp{M}}
}
\begin{tabproofwide}
  \proofstepwidestar[1]{\FranklFam M}{\rA}
  \proofstepwidestar[2]{X,Y\in\mathcal L_M}{\rA}
  \proofstepwidestar[1,2]{%
    \begin{aligned}[t]
      &\exists Z\in\UB_{\mathcal L_M,\subseteq}(\{X,Y\})\,
        \forall W\in\UB_{\mathcal L_M,\subseteq}(\{X,Y\})\,
          Z\subseteq W
    \end{aligned}}{%
    \ProofRefStack{%
      \FormulaRefAuto{FranklComplementPairSupremumExistence}[thm]\\
      (1,2)}}
  \proofstepwidestar[1,2]{%
    X\FamilyJoinOp{M}{}Y\in\mathcal L_M}{%
    \rAEa{\FormulaRefAuto{FranklComplementPairSupremumBounds}[thm]{1,2}}}
  \proofstepwidestar[1,2]{%
    \PairSup{\mathcal L_M}{\subseteq}{X}{Y}{%
      X\FamilyJoinOp{M}{}Y}}{%
    \ProofRefStack{\FormulaRefAuto{FranklComplementPairSupremum}[thm]\\
      (1,2)}}
  \proofstepwidestar[1]{%
    \PairJoinOp{\mathcal L_M,\FamilyJoinOp{M}}}{%
    \FormulaRefAuto{PairJoinOperation}[def]{3,4,5}}
\end{tabproofwide}

\FormulaThmDeltaK[Die konstruierte Operation bildet einen Halbverband]{%
  \FranklFam M\vdash\Semi{\mathcal L_M,\FamilyJoinOp{M}}
}{FranklComplementSemilattice}{%
  \DeltaRow{Mengen}{M}
  \DeltaRow{Relationsoperatoren}{\subseteq}
  \DeltaRow{Funktionen}{\FamilyJoinOp{M}}
}
\begin{tabproof}
  \proofstep{1}{\FranklFam M}{\rA}
  \proofstep{1}{\PartOrd{\mathcal L_M,\subseteq}}{%
    \FormulaRefAuto{FranklComplementOrderPartial}[thm]{1}}
  \proofstep{1}{\PairJoinOp{\mathcal L_M,\FamilyJoinOp{M}}}{%
    \FormulaRefAuto{FranklComplementPairJoinOperation}[thm]{1}}
  \proofstep{1}{\Semi{\mathcal L_M,\FamilyJoinOp{M}}}{%
    \FormulaRefAuto{PairJoinCharacterization}[thm]{2,3}}
\end{tabproof}

\FormulaThmDeltaK[Die induzierte Ordnung ist die Inklusionsordnung]{%
  \FranklFam M\dsep X,Y\in\mathcal L_M\vdash
  \bigl(X\InducedLeq Y\leftrightarrow X\subseteq Y\bigr)
}{FranklComplementOrderRecovered}{%
  \DeltaRow{Mengen}{M\dsep X\dsep Y}
  \DeltaRow{Relationsoperatoren}{\subseteq\dsep\InducedLeq}
  \DeltaRow{Funktionen}{\FamilyJoinOp{M}}
}
\begin{tabproof}
  \proofstep{1}{\FranklFam M}{\rA}
  \proofstep{2}{X\in\mathcal L_M}{\rA}
  \proofstep{3}{Y\in\mathcal L_M}{\rA}
  \proofstep{1}{\PartOrd{\mathcal L_M,\subseteq}}{%
    \FormulaRefAuto{FranklComplementOrderPartial}[thm]{1}}
  \proofstep{1}{\PairJoinOp{\mathcal L_M,\FamilyJoinOp{M}}}{%
    \FormulaRefAuto{FranklComplementPairJoinOperation}[thm]{1}}
  \proofstep{1,2,3}{%
    X\subseteq Y\leftrightarrow X\InducedLeq Y}{%
    \FormulaRefAuto{PairJoinRecoversOrder}[thm]{4,5,\rAI{2,3}}}
  \proofstep{1,2,3}{%
    X\InducedLeq Y\leftrightarrow X\subseteq Y}{%
    \FormulaRefAuto{P\leftrightarrow Q\vdash Q\leftrightarrow P}{6}}
\end{tabproof}
\FormulaThmDeltaK[Komplementdual einer Frankl-Familie]{%
  \FranklFam M\vdash
  \ZeroJoinSemi{\mathcal L_M,\FamilyJoinOp{M},\varnothing}
}{FranklComplementFiniteLattice}{%
  \DeltaRow{Mengen}{M\dsep X}
  \DeltaRow{Relationsoperatoren}{\subseteq}
  \DeltaRow{Funktionen}{\FamilyJoinOp{M}}
}
\begin{tabproofsplitwide}
  \proofpartwideindR[Träger des Nullelements]{%
    \FranklFam M\vdash\varnothing\in\mathcal L_M
  }{%
    \FranklFam M\vdash\varnothing\in\mathcal L_M
  }[FranklComplementZeroCarrier]
    \proofstepwidestar[1]{\FranklFam M}{\rA}
    \proofstepwidestar[1]{U_M\in M}{%
      \FormulaRefAuto{FiniteUnionClosedFamilyTotalUnion}[thm]{%
        \FormulaRefAuto{FranklFamFinite}[thm]{1},
        \FormulaRefAuto{FranklFamNonempty}[thm]{1},
        \FormulaRefAuto{FranklFamUnionClosed}[thm]{1}}}
    \proofstepwidestar[1]{\varnothing\in\mathcal L_M}{%
      \FormulaRefAuto{FranklComplementConstruction}[def]{2,
        \FormulaRefAuto{RelativeComplementSelf}[thm]}}
  \closeproofpartwideindR
  \proofpartwideindR[Supremum von Null und einem Trägerelement]{%
    \FranklFam M\dsep X\in\mathcal L_M\vdash
    \PairSup{\mathcal L_M}{\subseteq}{\varnothing}{X}{X}
  }{%
    \FranklFam M\dsep X\in\mathcal L_M\vdash
    \PairSup{\mathcal L_M}{\subseteq}{\varnothing}{X}{X}
  }[FranklComplementZeroSelfSupremum]
    \proofstepwidestar[1]{\FranklFam M}{\rA}
    \proofstepwidestar[2]{X\in\mathcal L_M}{\rA}
    \proofstepwidestar[1]{\PartOrd{\mathcal L_M,\subseteq}}{%
      \FormulaRefAuto{FranklComplementOrderPartial}[thm]{1}}
    \proofstepwidestar[1]{\varnothing\in\mathcal L_M}{%
      \FormulaRefAuto{FranklComplementZeroCarrier}[thm]{1}}
    \proofstepwidestar[1,2]{\varnothing\subseteq X}{%
      \FormulaRefAuto{\varnothing\subseteq A}}
    \proofstepwidestar[2]{X\subseteq X}{%
      \FormulaRefAuto{A\subseteq A}}
    \proofstepwidestar[1,2]{\{\varnothing,X\}\subseteq\mathcal L_M}{%
      \FormulaRefAuto{a \in A,\, b \in A \vdash
        \{a,b\} \subseteq A}{4,2}}
    \proofstepwidestar[1,2]{%
      X\in\UB_{\mathcal L_M,\subseteq}(\{\varnothing,X\})}{%
      \ProofRefStack{%
        \FormulaRefAuto{P\leftrightarrow Q, Q\vdash P}\\
        \bigl(\FormulaRefAuto{PairUpperBoundSetCharacterization}[thm]
          (3,4,2),\bigr.\\[-2pt]
        \bigl.\rAI{2,\rAI{5,6}}\bigr)}}
    \proofstepwidestar[9]{%
      W\in\UB_{\mathcal L_M,\subseteq}(\{\varnothing,X\})}{\rA}
    \proofstepwidestar[1,2]{%
      \begin{aligned}[t]
        &W\in\UB_{\mathcal L_M,\subseteq}(\{\varnothing,X\})
          \leftrightarrow{}\\[-2pt]
        &\bigl(W\in\mathcal L_M\land\varnothing\subseteq W
          \land X\subseteq W\bigr)
      \end{aligned}}{%
      \FormulaRefAuto{PairUpperBoundSetCharacterization}[thm]{3,4,2}}
    \proofstepwidestar[1,2,9]{%
      W\in\mathcal L_M\land\varnothing\subseteq W
        \land X\subseteq W}{%
      \FormulaRefAuto{P\leftrightarrow Q, P\vdash Q}{10,9}}
    \proofstepwidestar[1,2,9]{X\subseteq W}{%
      \rAEb{\rAEb{11}}}
    \proofstepwidestar[1,2]{%
      \begin{aligned}
        &\forall W\in\UB_{\mathcal L_M,\subseteq}
          (\{\varnothing,X\})\,{}\\[-2pt]
        &\quad X\subseteq W
      \end{aligned}}{%
      \rUI{\rRI{9,12}}}
    \proofstepwidestar[1,2]{%
      \PairSup{\mathcal L_M}{\subseteq}{\varnothing}{X}{X}}{%
      \FormulaRefAuto{SupremumCandidateEqualsSupremum}[thm]{3,7,8,13}}
  \closeproofpartwideindR
  \proofpartwideindR[Neutralität des Nullelements]{%
    \FranklFam M\dsep X\in\mathcal L_M\vdash
    \varnothing\FamilyJoinOp{M}{}X=X
  }{%
    \FranklFam M\dsep X\in\mathcal L_M\vdash
    \varnothing\FamilyJoinOp{M}{}X=X
  }[FranklComplementZeroIdentity]
    \proofstepwidestar[1]{\FranklFam M}{\rA}
    \proofstepwidestar[2]{X\in\mathcal L_M}{\rA}
    \proofstepwidestar[1]{\varnothing\in\mathcal L_M}{%
      \FormulaRefAuto{FranklComplementZeroCarrier}[thm]{1}}
    \proofstepwidestar[1,2]{%
      \PairSup{\mathcal L_M}{\subseteq}{\varnothing}{X}{%
        \varnothing\FamilyJoinOp{M}{}X}}{%
      \FormulaRefAuto{FranklComplementPairSupremum}[thm]{1,3,2}}
    \proofstepwidestar[1,2]{%
      \PairSup{\mathcal L_M}{\subseteq}{\varnothing}{X}{X}}{%
      \FormulaRefAuto{FranklComplementZeroSelfSupremum}[thm]{1,2}}
    \proofstepwidestar[1,2]{\varnothing\FamilyJoinOp{M}{}X=X}{%
      \FormulaRefAuto{a=b,\,c=b\vdash a=c}{4,5}}
  \closeproofpartwideindR
  \proofpartwide{Schluss}
    \proofstepwidestar[1]{\FranklFam M}{\rA}
    \proofstepwidestar[1]{\Semi{\mathcal L_M,\FamilyJoinOp{M}}}{%
      \FormulaRefAuto{FranklComplementSemilattice}[thm]{1}}
    \proofstepwidestar[1]{\varnothing\in\mathcal L_M}{%
      \FormulaRefAuto{FranklComplementZeroCarrier}[thm]{1}}
    \proofstepwidestar[4]{X\in\mathcal L_M}{\rA}
    \proofstepwidestar[1,4]{\varnothing\FamilyJoinOp{M}{}X=X}{%
      \FormulaRefAuto{FranklComplementZeroIdentity}[thm]{1,4}}
    \proofstepwidestar[1]{%
      \ZeroJoinSemi{\mathcal L_M,\FamilyJoinOp{M},\varnothing}}{%
      \FormulaRefAuto{ZeroJoinSemilattice}[def]{2,3,5}}
  \closeproofpartwide
\end{tabproofsplitwide}

\FormulaThmDeltaK[Der Komplementdual ist nichttrivial]{%
  \FranklFam M\vdash\exists X\in\mathcal L_M\,(X\neq\varnothing)
}{FranklComplementNontrivial}{%
  \DeltaRow{Mengen}{M\dsep X}
}
\begin{tabproof}
  \proofstep{1}{\FranklFam M}{\rA}
  \proofstep{1}{U_M\neq\varnothing}{%
    \FormulaRefAuto{FranklFamCarrierNonempty}[thm]{1}}
  \proofstep{1}{U_M\in\mathcal L_M}{%
    \FormulaRefAuto{FranklComplementConstruction}[def]{%
      \FormulaRefAuto{a\in A\cup\{a\}}[thm],
      \FormulaRefAuto{RelativeComplementEmpty}[thm]}}
  \proofstep{1}{\exists X\in\mathcal L_M\,(X\neq\varnothing)}{%
    \rEI{\rAI{3,2}}}
\end{tabproof}


\section{Reduktion auf Frankl-Familien mit leerer Menge}

Die Adjunktion der leeren Menge ist eine echte Normalisierung: Sie verändert
weder den Träger noch die Mengen, die ein gegebenes Trägerelement enthalten.
Endlichkeit und Vereinigungsabschluss bleiben erhalten. Daher genügt es, die
Frankl-Vermutung für Familien zu beweisen, welche die leere Menge enthalten.
Die folgenden Sätze registrieren alle dafür benötigten Einzelschritte.

\FormulaThmDeltaK[Die erweiterte Familie enthält die leere Menge]{%
  \varnothing\in M^\circ
}{FranklEmptyAdjunctionContainsEmpty}{%
  \DeltaRow{Mengen}{M}
}
\begin{tabproof}
  \proofstep{}{M^\circ=M\cup\{\varnothing\}}{%
    \FormulaRefAuto{FranklEmptyAdjunction}[def]}
  \proofstep{}{\varnothing\in M\cup\{\varnothing\}}{%
    \FormulaRefAuto{a\in A\cup\{a\}}[thm]}
  \proofstep{}{\varnothing\in M^\circ}{\rIE{1,2}}
\end{tabproof}

\FormulaThmDeltaK[Die Adjunktion der leeren Menge erhält den Träger]{%
  \bigcup M^\circ=\bigcup M
}{FranklEmptyAdjunctionCarrierUnchanged}{%
  \DeltaRow{Mengen}{M}
}
\begin{tabproof}
  \proofstep{}{M^\circ=M\cup\{\varnothing\}}{%
    \FormulaRefAuto{FranklEmptyAdjunction}[def]}
  \proofstep{}{\bigcup(M\cup\{\varnothing\})=(\bigcup M)\cup\varnothing}{%
    \FormulaRefAuto{UnionFamilyAdjunction}[thm]}
  \proofstep{}{\bigcup M^\circ=(\bigcup M)\cup\varnothing}{\rIE{1,2}}
  \proofstep{}{(\bigcup M)\cup\varnothing=\bigcup M}{%
    \FormulaRefAuto{A\cup\varnothing=A}}
  \proofstep{}{\bigcup M^\circ=\bigcup M}{%
    \FormulaRefAuto{a = b,\, b = c \vdash a = c}{3,4}}
\end{tabproof}

\FormulaThmDeltaK[Die erweiterte Familie ist nichtleer]{%
  M^\circ\neq\varnothing
}{FranklEmptyAdjunctionNonempty}{%
  \DeltaRow{Mengen}{M}
}
\begin{tabproof}
  \proofstep{}{\varnothing\in M^\circ}{%
    \FormulaRefAuto{FranklEmptyAdjunctionContainsEmpty}[thm]}
  \proofstep{}{M^\circ\neq\varnothing}{%
    \FormulaRefAuto{a\in A\vdash A\neq\varnothing}[thm]{1}}
\end{tabproof}

\FormulaThmDeltaK[Die erweiterte Frankl-Familie hat nichtleeren Träger]{%
  \FranklFam M\vdash\bigcup M^\circ\neq\varnothing
}{FranklEmptyAdjunctionCarrierNonempty}{%
  \DeltaRow{Mengen}{M}
}
\begin{tabproof}
  \proofstep{1}{\FranklFam M}{\rA}
  \proofstep{1}{\bigcup M\neq\varnothing}{%
    \FormulaRefAuto{FranklFamCarrierNonempty}[thm]{1}}
  \proofstep{}{\bigcup M^\circ=\bigcup M}{%
    \FormulaRefAuto{FranklEmptyAdjunctionCarrierUnchanged}[thm]}
  \proofstep{}{\bigcup M=\bigcup M^\circ}{%
    \FormulaRefAuto{a=b\vdash b=a}[thm]{3}}
  \proofstep{1}{\bigcup M^\circ\neq\varnothing}{\rIE{4,2}}
\end{tabproof}

Die Endlichkeit ist bereits durch
\FormulaRefAuto{FranklAugmentedFamilyFinite}[thm] und der
Vereinigungsabschluss durch
\FormulaRefAuto{FranklAugmentedFamilyUnionClosed}[thm] bewiesen. Zusammen
mit den beiden vorstehenden Aussagen folgt die Erhaltung aller vier
Voraussetzungen.

\FormulaThmDeltaK[Adjunktion der leeren Menge erhält Frankl-Familien]{%
  \FranklFam M\vdash\FranklFam{M^\circ}
}{FranklEmptyAdjunctionPreservesFamily}{%
  \DeltaRow{Mengen}{M}
}
\begin{tabproof}
  \proofstep{1}{\FranklFam M}{\rA}
  \proofstep{}{M^\circ\neq\varnothing}{%
    \FormulaRefAuto{FranklEmptyAdjunctionNonempty}[thm]}
  \proofstep{1}{\bigcup M^\circ\neq\varnothing}{%
    \FormulaRefAuto{FranklEmptyAdjunctionCarrierNonempty}[thm]{1}}
  \proofstep{1}{\Finite{M^\circ}}{%
    \FormulaRefAuto{FranklAugmentedFamilyFinite}[thm]{1}}
  \proofstep{1}{\UnionClosedFam(M^\circ)}{%
    \FormulaRefAuto{FranklAugmentedFamilyUnionClosed}[thm]{1}}
  \proofstep{1}{\FranklFam{M^\circ}}{%
    \FormulaRefAuto{Frankl-Familie}[def]{2,3,4,5}}
\end{tabproof}

Für den Rücktransport genügt es, eine Injektion auf die kleinere
vermeidende Teilfamilie einzuschränken. Die enthaltende Teilfamilie bleibt
unverändert.

\FormulaThmDeltaK[Halbhäufige Elemente werden zur ursprünglichen Familie zurücktransportiert]{%
  \FranklFam M\dsep\HalfOcc{a}{M^\circ}\vdash\HalfOcc{a}{M}
}{FranklEmptyAdjunctionHalfOccurrenceBack}{%
  \DeltaRow{Mengen}{M\dsep a}
  \DeltaRow{Funktionen}{F}
}
\begin{tabproofwide}
  \proofstepwidestar[1]{\FranklFam M}{\rA}
  \proofstepwidestar[2]{\HalfOcc{a}{M^\circ}}{\rA}
  \proofstepwidestar[2]{%
    a\in\bigcup M^\circ\land
    \exists F\colon\FamAvoiding{M^\circ}{a}
      \inj\FamContaining{M^\circ}{a}}{%
    \FormulaRefAuto{Halbhäufiges Element}[def]{2}}
  \proofstepwidestar[]{\bigcup M^\circ=\bigcup M}{%
    \FormulaRefAuto{FranklEmptyAdjunctionCarrierUnchanged}[thm]}
  \proofstepwidestar[2]{a\in\bigcup M}{\rIE{4,\rAEa{3}}}
  \proofstepwidestar[]{M\subseteq M^\circ}{%
    \rIE{\FormulaRefAuto{FranklEmptyAdjunction}[def],
      \FormulaRefAuto{A\subseteq A\cup B}[thm]}}
  \proofstepwidestar[]{%
    \FamAvoiding{M}{a}\subseteq\FamAvoiding{M^\circ}{a}}{%
    \FormulaRefAuto{FamAvoidingMonotone}[thm]{6}}
  \proofstepwidestar[1]{%
    \FamContaining{M^\circ}{a}=\FamContaining{M}{a}}{%
    \FormulaRefAuto{FranklAugmentedContainingUnchanged}[thm]{1}}
  \proofstepwidestar[2]{%
    \exists F\colon\FamAvoiding{M^\circ}{a}
      \inj\FamContaining{M^\circ}{a}}{\rAEb{3}}
  \proofstepwidestar[10]{%
    F\colon\FamAvoiding{M^\circ}{a}
      \inj\FamContaining{M^\circ}{a}}{\rA}
  \proofstepwidestar[7,10]{%
    F\restriction_{\FamAvoiding{M}{a}}\colon
      \FamAvoiding{M}{a}\inj\FamContaining{M^\circ}{a}}{%
    \FormulaRefAuto{F\colon A\inj B\dsep C\subseteq A\vdash
      F\restriction_C\colon C\inj B}[thm]{10,7}}
  \proofstepwidestar[1,7,10]{%
    \exists F\colon\FamAvoiding{M}{a}\inj\FamContaining{M}{a}}{%
    \rEI{\rIE{8,11}}}
  \proofstepwidestar[1,2]{%
    \exists F\colon\FamAvoiding{M}{a}\inj\FamContaining{M}{a}}{%
    \rEE{9,10,12}}
  \proofstepwidestar[1,2]{\HalfOcc{a}{M}}{%
    \FormulaRefAuto{Halbhäufiges Element}[def]{5,13}}
\end{tabproofwide}

\FormulaThmDeltaK[Frankl-Eigenschaft wird zur ursprünglichen Familie zurücktransportiert]{%
  \FranklFam M\dsep\Frankl{M^\circ}\vdash\Frankl M
}{FranklEmptyAdjunctionFranklBack}{%
  \DeltaRow{Mengen}{M\dsep a}
}
\begin{tabproof}
  \proofstep{1}{\FranklFam M}{\rA}
  \proofstep{2}{\Frankl{M^\circ}}{\rA}
  \proofstep{2}{\exists a\,\HalfOcc{a}{M^\circ}}{%
    \FormulaRefAuto{Frankl-Eigenschaft}[def]{2}}
  \proofstep{4}{\HalfOcc{a}{M^\circ}}{\rA}
  \proofstep{1,4}{\HalfOcc{a}{M}}{%
    \FormulaRefAuto{FranklEmptyAdjunctionHalfOccurrenceBack}[thm]{1,4}}
  \proofstep{1,4}{\Frankl M}{%
    \FormulaRefAuto{Frankl-Eigenschaft}[def]{\rEI{5}}}
  \proofstep{1,2}{\Frankl M}{\rEE{3,4,6}}
\end{tabproof}

\FormulaThmDeltaK[Reduktion auf Frankl-Familien mit leerer Menge]{%
  \begin{aligned}[t]
    &\forall M\,(\FranklFam M\to\Frankl M)\\[-2pt]
    &\quad\leftrightarrow
      \forall M\,\bigl((\FranklFam M\land\varnothing\in M)
        \to\Frankl M\bigr)
  \end{aligned}
}{FranklEmptyFamilyGlobalEquivalence}{%
  \DeltaRow{Mengen}{M}
}
\begin{tabproofwide}
  \proofstepwidestar[1]{%
    \forall M\,(\FranklFam M\to\Frankl M)}{\rA}
  \proofstepwidestar[2]{%
    \FranklFam M\land\varnothing\in M}{\rA}
  \proofstepwidestar[2]{\FranklFam M}{\rAEa{2}}
  \proofstepwidestar[1]{\FranklFam M\to\Frankl M}{\rUE{1}}
  \proofstepwidestar[1,2]{\Frankl M}{\rRE{4,3}}
  \proofstepwidestar[1]{%
    (\FranklFam M\land\varnothing\in M)\to\Frankl M}{\rRI{2,5}}
  \proofstepwidestar[1]{%
    \forall M\,\bigl((\FranklFam M\land\varnothing\in M)
      \to\Frankl M\bigr)}{\rUI{6}}
  \proofstepwidestar[]{%
    \begin{aligned}[t]
      &\forall M\,(\FranklFam M\to\Frankl M)\to{}\\[-2pt]
      &\quad\forall M\,\bigl((\FranklFam M\land\varnothing\in M)
        \to\Frankl M\bigr)
    \end{aligned}}{\rRI{1,7}}
  \proofstepwidestar[9]{%
    \forall M\,\bigl((\FranklFam M\land\varnothing\in M)
      \to\Frankl M\bigr)}{\rA}
  \proofstepwidestar[10]{\FranklFam M}{\rA}
  \proofstepwidestar[10]{\FranklFam{M^\circ}}{%
    \FormulaRefAuto{FranklEmptyAdjunctionPreservesFamily}[thm]{10}}
  \proofstepwidestar[]{\varnothing\in M^\circ}{%
    \FormulaRefAuto{FranklEmptyAdjunctionContainsEmpty}[thm]}
  \proofstepwidestar[10]{%
    \FranklFam{M^\circ}\land\varnothing\in M^\circ}{\rAI{11,12}}
  \proofstepwidestar[9]{%
    (\FranklFam{M^\circ}\land\varnothing\in M^\circ)
      \to\Frankl{M^\circ}}{\rUE{9}}
  \proofstepwidestar[9,10]{\Frankl{M^\circ}}{\rRE{14,13}}
  \proofstepwidestar[9,10]{\Frankl M}{%
    \FormulaRefAuto{FranklEmptyAdjunctionFranklBack}[thm]{10,15}}
  \proofstepwidestar[9]{\FranklFam M\to\Frankl M}{\rRI{10,16}}
  \proofstepwidestar[9]{%
    \forall M\,(\FranklFam M\to\Frankl M)}{\rUI{17}}
  \proofstepwidestar[]{%
    \begin{aligned}[t]
      &\forall M\,\bigl((\FranklFam M\land\varnothing\in M)
        \to\Frankl M\bigr)\to{}\\[-2pt]
      &\quad\forall M\,(\FranklFam M\to\Frankl M)
    \end{aligned}}{\rRI{9,18}}
  \proofstepwidestar[]{%
    \begin{aligned}[t]
      &\forall M\,(\FranklFam M\to\Frankl M)\\[-2pt]
      &\quad\leftrightarrow
        \forall M\,\bigl((\FranklFam M\land\varnothing\in M)
          \to\Frankl M\bigr)
    \end{aligned}}{%
    \FormulaRefAuto{P\leftrightarrow Q\eqvdash
      (P\to Q)\land(Q\to P)}{\rAI{8,19}}}
\end{tabproofwide}

Insbesondere darf jeder Beweisversuch die leere Menge ohne Verlust
adjungieren. Die Voraussetzung \(\bigcup M\neq\varnothing\) bleibt dabei
unverändert; die ausgeschlossene Familie \(\{\varnothing\}\) wird also nicht
versehentlich in die Vermutung aufgenommen.

\section{Privater Punkt und Rücktransport}

\FormulaThmDeltaK[Charakterisierung eines privaten Punkts]{%
  \begin{aligned}[t]
    &\ZeroJoinSemi{L,\JoinOp,\varnothing}\dsep\InterClosedFam L\dsep
      \forall Q,S\in L(Q\InducedLeq S\leftrightarrow Q\subseteq S)\dsep{}\\
    &P_\ast\in D_L(P)\dsep
      \forall Q\in D_L(P)Q\InducedLeq P_\ast\dsep
      a\in P\setminus P_\ast\vdash{}\\
    &\forall Q\in L\,(a\in Q\leftrightarrow
      Q\in\PrincipalFilter{L,\InducedLeq}{P})
  \end{aligned}
}{JoinIrreduciblePrivatePointCharacterization}{%
  \DeltaRow{Mengen}{L\dsep P\dsep P_\ast\dsep Q\dsep S\dsep a}
  \DeltaRow{Funktionen}{\JoinOp}
}
\begin{tabproofwide}
  \proofstepwidestar[1]{\ZeroJoinSemi{L,\JoinOp,\varnothing}}{\rA}
  \proofstepwidestar[2]{\InterClosedFam L}{\rA}
  \proofstepwidestar[3]{\forall Q,S\in L
    (Q\InducedLeq S\leftrightarrow Q\subseteq S)}{\rA}
  \proofstepwidestar[4]{P_\ast\in D_L(P)}{\rA}
  \proofstepwidestar[5]{\forall Q\in D_L(P)
    Q\InducedLeq P_\ast}{\rA}
  \proofstepwidestar[6]{a\in P\setminus P_\ast}{\rA}
  \proofstepwidestar[1]{\Semi{L,\JoinOp}}{%
    \FormulaRefAuto{ZeroJoinBaseSemilatticeAxiom}[ax]{1}}
  \proofstepwidestar[4]{P_\ast\in L\land
    P_\ast\InducedLeq P\land P_\ast\neq P}{%
    \FormulaRefAuto{JoinIrreducibleProperLowerSet}[def]{4}}
  \proofstepwidestar[1,4]{P\in L}{%
    \FormulaRefAuto{JoinOrderSecondCarrier}[thm]{7,\rAEa{\rAEb{8}}}}
  \proofstepwidestar[6]{a\in P}{%
    \rAEa{\FormulaRefAuto{x \in A \setminus B \eqvdash
      x \in A \land x \notin B}{6}}}
  \proofstepwidestar[6]{a\notin P_\ast}{%
    \rAEb{\FormulaRefAuto{x \in A \setminus B \eqvdash
      x \in A \land x \notin B}{6}}}
  \proofstepwidestar[12]{Q\in L}{\rA}
  \proofstepwidestar[13]{a\in Q}{\rA}
  \proofstepwidestar[14]{P\nsubseteq Q}{\rA}
  \proofstepwidestar[2,9,12]{P\cap Q\in L}{%
    \FormulaRefAuto{IntersectionClosedFamily}[def]{2,9,12}}
  \proofstepwidestar[]{P\cap Q\subseteq P}{\FormulaRefAuto{A\cap B\subseteq A}}
  \proofstepwidestar[3,9,12]{P\cap Q\InducedLeq P}{%
    \FormulaRefAuto{P\leftrightarrow Q,Q\vdash P}{%
      \rUE{\rUE{3}},16}}
  \proofstepwidestar[14]{P\cap Q\neq P}{%
    \FormulaRefAuto{P\leftrightarrow Q,\neg P\vdash\neg Q}{%
      \FormulaRefAuto{A \subseteq B \eqvdash A \cap B = A},14}}
  \proofstepwidestar[3,9,12,14]{P\cap Q\in D_L(P)}{%
    \FormulaRefAuto{JoinIrreducibleProperLowerSet}[def]{15,17,18}}
  \proofstepwidestar[3,5,9,12,14]{%
    P\cap Q\InducedLeq P_\ast}{\rRE{19,\rUE{5}}}
  \proofstepwidestar[3,5,9,12,14]{P\cap Q\subseteq P_\ast}{%
    \FormulaRefAuto{P\leftrightarrow Q,P\vdash Q}{%
      \rUE{\rUE{3}},20}}
  \proofstepwidestar[10,13]{a\in P\cap Q}{%
    \FormulaRefAuto{x\in A\cap B\eqvdash x\in A\land x\in B}{%
      \rAI{10,13}}}
  \proofstepwidestar[3,5,9,10,12,13,14]{a\in P_\ast}{%
    \FormulaRefAuto{A\subseteq B,\,x\in A\vdash x\in B}{21,22}}
  \proofstepwidestar[3,5,6,9,10,12,13,14]{\bot}{\rBI{23,11}}
  \proofstepwidestar[1,2,3,4,5,6,12,13]{P\subseteq Q}{%
    \FormulaRefAuto{rDN}{\rCI{14,24}}}
  \proofstepwidestar[1,2,3,4,5,6,12]{a\in Q\to P\subseteq Q}{%
    \rRI{13,25}}
  \proofstepwidestar[27]{P\subseteq Q}{\rA}
  \proofstepwidestar[10,27]{a\in Q}{%
    \FormulaRefAuto{A\subseteq B,\,x\in A\vdash x\in B}{27,10}}
  \proofstepwidestar[10,12]{P\subseteq Q\to a\in Q}{\rRI{27,28}}
  \proofstepwidestar[1,2,3,4,5,6,12]{a\in Q\leftrightarrow P\subseteq Q}{%
    \FormulaRefAuto{P \leftrightarrow Q \eqvdash
      (P \rightarrow Q)\land(Q\rightarrow P)}{\rAI{26,29}}}
  \proofstepwidestar[12]{Q\in
    \PrincipalFilter{L,\InducedLeq}{P}\leftrightarrow
    P\InducedLeq Q}{%
    \FormulaRefAuto{PrincipalFilter}[def]{12}}
  \proofstepwidestar[3,9,12]{P\InducedLeq Q\leftrightarrow
    P\subseteq Q}{\rUE{\rUE{3}}}
  \proofstepwidestar[3,9,12]{Q\in
    \PrincipalFilter{L,\InducedLeq}{P}\leftrightarrow
    P\subseteq Q}{\rChain{31,32}}
  \proofstepwidestar[3,9,12]{P\subseteq Q\leftrightarrow Q\in
    \PrincipalFilter{L,\InducedLeq}{P}}{%
    \FormulaRefAuto{P\leftrightarrow Q\vdash Q\leftrightarrow P}{33}}
  \proofstepwidestar[1,2,3,4,5,6,12]{a\in Q\leftrightarrow
    Q\in\PrincipalFilter{L,\InducedLeq}{P}}{\rChain{30,34}}
  \proofstepwidestar[1,2,3,4,5,6]{\forall Q\in L\,(a\in Q\leftrightarrow
    Q\in\PrincipalFilter{L,\InducedLeq}{P})}{%
    \rUI{\rRI{12,35}}}
\end{tabproofwide}

\FormulaThmDeltaK[Privater Punkt eines \(\JoinOp\)-Irreduziblen]{%
  \begin{aligned}[t]
    &\Finite L\dsep\ZeroJoinSemi{L,\JoinOp,\varnothing}\dsep
      \InterClosedFam L\dsep{}\\
    &\forall Q,S\in L(Q\InducedLeq S\leftrightarrow Q\subseteq S)
      \dsep\JoinIrr{L,\JoinOp,\varnothing}{P}\vdash{}\\
    &\exists a\in P\,\forall Q\in L\,(a\in Q\leftrightarrow
      Q\in\PrincipalFilter{L,\InducedLeq}{P})
  \end{aligned}
}{JoinIrreduciblePrivatePoint}{%
  \DeltaRow{Mengen}{L\dsep P\dsep P_\ast\dsep Q\dsep S\dsep a}
  \DeltaRow{Funktionen}{\JoinOp}
}
\begin{tabproofwide}
  \proofstepwidestar[1]{\Finite L}{\rA}
  \proofstepwidestar[2]{\ZeroJoinSemi{L,\JoinOp,\varnothing}}{\rA}
  \proofstepwidestar[3]{\InterClosedFam L}{\rA}
  \proofstepwidestar[4]{\forall Q,S\in L
    (Q\InducedLeq S\leftrightarrow Q\subseteq S)}{\rA}
  \proofstepwidestar[5]{\JoinIrr{L,\JoinOp,\varnothing}{P}}{\rA}
  \proofstepwidestar[1,2,5]{\exists P_\ast\in D_L(P)\,
    \forall Q\in D_L(P)Q\InducedLeq P_\ast}{%
    \FormulaRefAuto{JoinIrreducibleLargestProperLower}[thm]{1,2,5}}
  \proofstepwidestar[7]{P_\ast\in D_L(P)\land
    \forall Q\in D_L(P)Q\InducedLeq P_\ast}{\rA}
  \proofstepwidestar[7]{P_\ast\in L\land
    P_\ast\InducedLeq P\land P_\ast\neq P}{%
    \FormulaRefAuto{JoinIrreducibleProperLowerSet}[def]{\rAEa{7}}}
  \proofstepwidestar[4,7]{P_\ast\subseteq P}{%
    \FormulaRefAuto{P\leftrightarrow Q,P\vdash Q}{%
      \rUE{\rUE{4}},\rAEa{\rAEb{8}}}}
  \proofstepwidestar[4,7]{P_\ast\subsetneq P}{%
    \FormulaRefAuto{A\subset B\coloneq A\subseteq B\land A\neq B}[def]{%
      \rAI{9,\rAEb{\rAEb{8}}}}}
  \proofstepwidestar[4,7]{\exists a\in P\setminus P_\ast}{%
    \FormulaRefAuto{A\subset B\vdash \exists x\in B\setminus A}{10}}
  \proofstepwidestar[12]{a\in P\setminus P_\ast}{\rA}
  \proofstepwidestar[2,3,4,7,12]{\forall Q\in L\,(a\in Q\leftrightarrow
    Q\in\PrincipalFilter{L,\InducedLeq}{P})}{%
    \FormulaRefAuto{JoinIrreduciblePrivatePointCharacterization}[thm]{%
      2,3,4,\rAEa{7},\rAEb{7},12}}
  \proofstepwidestar[12]{a\in P}{%
    \rAEa{\FormulaRefAuto{x \in A \setminus B \eqvdash
      x \in A \land x \notin B}{12}}}
  \proofstepwidestar[1,2,3,4,5]{\exists a\in P\,\forall Q\in L\,
    (a\in Q\leftrightarrow
      Q\in\PrincipalFilter{L,\InducedLeq}{P})}{%
    \rEE{6,7,\rEE{11,12,\rEI{\rAI{14,13}}}}}
\end{tabproofwide}

\FormulaThmDeltaK[Komplementfaser der vermeidenden Teilfamilie]{%
  \begin{aligned}[t]
    &\FranklFam M\dsep H\subseteq\mathcal L_M\dsep a\in U_M\dsep{}\\
    &\forall Q\in\mathcal L_M\,(a\in Q\leftrightarrow Q\in H)
      \vdash
      \kappa_M^{-1}[H]=\FamAvoiding{M^\circ}{a}
  \end{aligned}
}{FranklComplementAvoidingFibre}{%
  \DeltaRow{Mengen}{M\dsep H\dsep Q\dsep X\dsep a}
  \DeltaRow{Funktionen}{\kappa_M}
}
\begin{tabproofwide}
  \proofstepwidestar[1]{\FranklFam M}{\rA}
  \proofstepwidestar[2]{H\subseteq\mathcal L_M}{\rA}
  \proofstepwidestar[3]{a\in U_M}{\rA}
  \proofstepwidestar[4]{\forall Q\in\mathcal L_M
    (a\in Q\leftrightarrow Q\in H)}{\rA}
  \proofstepwidestar[1]{\kappa_M\colon M^\circ\bij\mathcal L_M}{%
    \FormulaRefAuto{FranklComplementMapBijective}[thm]{1}}
  \proofstepwidestar[1]{%
    \forall Y\in M^\circ\,\kappa_M(Y)=U_M\setminus Y}{%
    \FormulaRefAuto{FranklComplementMap}[def]{1}}
  \proofstepwidestar[1]{%
    X\in M^\circ\rightarrow\kappa_M(X)=U_M\setminus X}{%
    \rUE{6}}
  \proofstepwide[1,2]{X\in\kappa_M^{-1}[H]}{\leftrightarrow}{%
    X\in M^\circ\land\kappa_M(X)\in H}{%
    \FormulaRefAuto{F\colon A\to B\dsep C\subseteq B
      \vdash(x\in F^{-1}[C]\leftrightarrow(x\in A\land F(x)\in C))}{5,2}}
  \proofstepwide[1,2,4]{}{\leftrightarrow}{%
    X\in M^\circ\land a\in\kappa_M(X)}{%
    \rLRS{\FormulaRefAuto{P\leftrightarrow Q\vdash Q\leftrightarrow P}{%
      \rUE{4}},8}}
  \proofstepwide[1,2,4]{}{\leftrightarrow}{%
    X\in M^\circ\land a\in U_M\setminus X}{%
    \rIE{7}}
  \proofstepwide[1,2,3,4]{}{\leftrightarrow}{%
    X\in M^\circ\land a\notin X}{%
    \FormulaRefAuto{x\in A\setminus B\eqvdash x\in A\land x\notin B}{3}}
  \proofstepwide[1,2,3,4]{}{\leftrightarrow}{%
    X\in\FamAvoiding{M^\circ}{a}}{%
    \FormulaRefAuto{\FamAvoiding{M}{a}\coloneqq
      \{X\in M\mid a\notin X\}}}
  \proofstepwidestar[1,2,3,4]{%
    \kappa_M^{-1}[H]=\FamAvoiding{M^\circ}{a}}{%
    \FormulaRefAuto{\forall x\,(x\in A\leftrightarrow x\in B)
      \eqvdash A=B}{\rUI{12}}}
\end{tabproofwide}

\FormulaThmDeltaK[Komplementfaser der enthaltenden Teilfamilie]{%
  \begin{aligned}[t]
    &\FranklFam M\dsep H\subseteq\mathcal L_M\dsep a\in U_M\dsep{}\\
    &\forall Q\in\mathcal L_M\,(a\in Q\leftrightarrow Q\in H)
      \vdash
      \kappa_M^{-1}[\mathcal L_M\setminus H]=
      \FamContaining{M^\circ}{a}
  \end{aligned}
}{FranklComplementContainingFibre}{%
  \DeltaRow{Mengen}{M\dsep H\dsep Q\dsep X\dsep a}
  \DeltaRow{Funktionen}{\kappa_M}
}
Zur Abkürzung setzen wir
\[
  \begin{aligned}
    \mathcal A&\coloneqq M^\circ,&
    \mathcal K&\coloneqq\mathcal L_M\setminus H,\\
    \mathcal F&\coloneqq\FamContaining{M^\circ}{a},&
    \kappa&\coloneqq\kappa_M.
  \end{aligned}
\]
\begin{tabproofwide}
  \proofstepwidestar[1]{\FranklFam M}{\rA}
  \proofstepwidestar[2]{H\subseteq\mathcal L_M}{\rA}
  \proofstepwidestar[3]{a\in U_M}{\rA}
  \proofstepwidestar[4]{\forall Q\in\mathcal L_M
    (a\in Q\leftrightarrow Q\in H)}{\rA}
  \proofstepwidestar[1]{\kappa\colon\mathcal A\bij\mathcal L_M}{%
    \FormulaRefAuto{FranklComplementMapBijective}[thm]{1}}
  \proofstepwidestar[1]{%
    \forall Y\in\mathcal A\,\kappa(Y)=U_M\setminus Y}{%
    \FormulaRefAuto{FranklComplementMap}[def]{1}}
  \proofstepwidestar[1]{%
    X\in\mathcal A\rightarrow\kappa(X)=U_M\setminus X}{%
    \rUE{6}}
  \proofstepwide[1]{X\in\kappa^{-1}[\mathcal K]}
    {\leftrightarrow}{%
    X\in\mathcal A\land\kappa(X)\in\mathcal K}{%
    \ProofRefStack{%
      \FormulaRefAuto{F\colon A\to B\dsep C\subseteq B
        \vdash(x\in F^{-1}[C]\leftrightarrow
        (x\in A\land F(x)\in C))}\\
      \FormulaRefAuto{A\setminus B\subseteq A}\\
      \bigl(5\bigr)}}
  \proofstepwide[1,4]{}{\leftrightarrow}{%
    X\in\mathcal A\land a\notin\kappa(X)}{%
    \ProofRefStack{%
      \FormulaRefAuto{x\in A\setminus B
        \eqvdash x\in A\land x\notin B}\\
      \FormulaRefAuto{P\leftrightarrow Q
        \eqvdash\neg P\leftrightarrow\neg Q}[thm]\\
      \bigl(\rUE{4}\bigr)}}
  \proofstepwide[1,3,4]{}{\leftrightarrow}{%
    X\in\mathcal A\land a\in X}{%
    \ProofRefStack{%
      \rIE{7}\\
      \FormulaRefAuto{x\in A\setminus B
        \eqvdash x\in A\land x\notin B}\,(3)\\
      \FormulaRefAuto{rDN}}}
  \proofstepwide[1,3,4]{}{\leftrightarrow}{%
    X\in\mathcal F}{%
    \FormulaRefAuto{\FamContaining{M}{a}\coloneqq
      \{X\in M\mid a\in X\}}}
  \proofstepwidestar[1,2,3,4]{\kappa^{-1}[\mathcal K]=\mathcal F}{%
    \ProofRefStack{%
      \FormulaRefAuto{\forall x\,(x\in A\leftrightarrow x\in B)
        \eqvdash A=B}\\
      \bigl(\rUI{11}\bigr)}}
\end{tabproofwide}

\FormulaThmDeltaK[Transport eines privaten Punkts]{%
  \begin{aligned}[t]
    &\FranklFam M\dsep H\subseteq\mathcal L_M\dsep a\in U_M\dsep{}\\
    &\forall Q\in\mathcal L_M\,(a\in Q\leftrightarrow Q\in H)
      \dsep\AtMostHalf{H}{\mathcal L_M}
      \vdash\HalfOcc{a}{M}
  \end{aligned}
}{FranklComplementPrivatePointTransport}{%
  \DeltaRow{Mengen}{M\dsep H\dsep Q\dsep X\dsep a}
  \DeltaRow{Funktionen}{\kappa_M\dsep F}
}
\begin{tabproofwide}
  \proofstepwidestar[1]{\FranklFam M}{\rA}
  \proofstepwidestar[2]{H\subseteq\mathcal L_M}{\rA}
  \proofstepwidestar[3]{a\in U_M}{\rA}
  \proofstepwidestar[4]{\forall Q\in\mathcal L_M
    (a\in Q\leftrightarrow Q\in H)}{\rA}
  \proofstepwidestar[5]{\AtMostHalf{H}{\mathcal L_M}}{\rA}
  \proofstepwidestar[1]{\kappa_M\colon M^\circ\bij\mathcal L_M}{%
    \FormulaRefAuto{FranklComplementMapBijective}[thm]{1}}
  \proofstepwidestar[1,2,3,4]{\kappa_M^{-1}[H]=
    \FamAvoiding{M^\circ}{a}}{%
    \FormulaRefAuto{FranklComplementAvoidingFibre}[thm]{1,2,3,4}}
  \proofstepwidestar[1,2,3,4]{%
    \kappa_M^{-1}[\mathcal L_M\setminus H]=
    \FamContaining{M^\circ}{a}}{%
    \FormulaRefAuto{FranklComplementContainingFibre}[thm]{1,2,3,4}}
  \proofstepwidestar[1,2,5]{%
    \exists F\colon\kappa_M^{-1}[H]
      \inj\kappa_M^{-1}[\mathcal L_M\setminus H]}{%
    \FormulaRefAuto{BijectionAtMostHalfTransport}[thm]{6,2,5}}
  \proofstepwidestar[1,2,3,4,5]{%
    \exists F\colon\FamAvoiding{M^\circ}{a}
      \inj\FamContaining{M^\circ}{a}}{\rIE{7,\rIE{8,9}}}
  \proofstepwidestar[1]{M\subseteq M^\circ}{%
    \rIE{\FormulaRefAuto{FranklEmptyAdjunction}[def]{1},
      \FormulaRefAuto{A\subseteq A\cup B}}}
  \proofstepwidestar[1]{\FamAvoiding{M}{a}\subseteq
    \FamAvoiding{M^\circ}{a}}{%
    \FormulaRefAuto{FamAvoidingMonotone}[thm]{11}}
  \proofstepwidestar[1]{\FamContaining{M^\circ}{a}=
    \FamContaining{M}{a}}{%
    \FormulaRefAuto{FranklAugmentedContainingUnchanged}[thm]{1}}
  \proofstepwidestar[14]{F\colon\FamAvoiding{M^\circ}{a}
    \inj\FamContaining{M^\circ}{a}}{\rA}
  \proofstepwidestar[12,14]{F\restriction_{\FamAvoiding{M}{a}}
    \colon\FamAvoiding{M}{a}\inj\FamContaining{M^\circ}{a}}{%
    \FormulaRefAuto{F\colon A\inj B\dsep C\subseteq A\vdash
      F\restriction_C\colon C\inj B}{14,12}}
  \proofstepwidestar[12,13,14]{\exists F\colon\FamAvoiding{M}{a}
    \inj\FamContaining{M}{a}}{\rEI{\rIE{13,15}}}
  \proofstepwidestar[1,2,3,4,5]{\exists F\colon\FamAvoiding{M}{a}
    \inj\FamContaining{M}{a}}{\rEE{10,14,16}}
  \proofstepwidestar[1,2,3,4,5]{\HalfOcc{a}{M}}{%
    \FormulaRefAuto{Halbhäufiges Element}[def]{3,17}}
\end{tabproofwide}

\FormulaDefDeltaK[Hauptfilter in der Komplementfamilie]{%
  H_M(P)\coloneqq
  \PrincipalFilter{\mathcal L_M,\InducedLeq}{P}%
}{FranklComplementPrincipalFilterAbbreviation}{%
  \DeltaRow{Mengen}{M\dsep P}
  \DeltaRow{Relationsoperatoren}{\InducedLeq}
  \DeltaRow{\textbf{Neues Symbol}}{}
  \DeltaPrem{Hauptfilter in der Komplementfamilie}{H_M(P)}
}

\FormulaThmDeltaK[Halbverbandsform impliziert Mengenform]{%
  \begin{aligned}[t]
    &\FranklFam M\dsep{}\\
    &\forall A\,\forall \JoinOp\,\forall o\,\Bigl(
      \Finite A\land\ZeroJoinSemi{A,\JoinOp,o}\land{}\\
    &\qquad
      \exists x\in A\,(x\neq o)
      \to\FranklSemi{A,\JoinOp,o}\Bigr)
      \vdash\Frankl M
  \end{aligned}
}{FranklSemilatticeImpliesSet}{%
  \DeltaRow{Mengen}{A\dsep M\dsep\mathcal L_M\dsep
    P\dsep Q\dsep X\dsep a\dsep o}
  \DeltaRow{Funktionen}{\JoinOp\dsep\FamilyJoinOp{M}\dsep F}
}

Für den ersten Teilbeweis verwenden wir die Abkürzungen
\[
  \begin{aligned}
    \mathsf H(B,\star,e)
      &\coloneqq
        \Finite B\land\ZeroJoinSemi{B,\star,e}\land
        \exists x\in B\,(x\neq e),\\
    \mathsf P(B,\star,e)
      &\coloneqq\FranklSemi{B,\star,e},\\
    L&\coloneqq\mathcal L_M,
    &\diamond&\coloneqq\FamilyJoinOp{M}.
  \end{aligned}
\]
\begin{tabproofsplitwide}
  \proofpartwideindR[Anwendung auf die Komplementfamilie]{%
    \begin{aligned}[t]
      &\forall A\,\forall\JoinOp\,\forall o\,\Bigl(
        \Finite A\land\ZeroJoinSemi{A,\JoinOp,o}\land
        \exists x\in A\,(x\neq o)\\[-2pt]
      &\qquad\to\FranklSemi{A,\JoinOp,o}\Bigr)
        \dsep\FranklFam M\vdash
        \FranklSemi{\mathcal L_M,\FamilyJoinOp{M},\varnothing}
    \end{aligned}
  }{%
    \begin{aligned}[t]
      &\forall A\,\forall\JoinOp\,\forall o\,\Bigl(
        \Finite A\land\ZeroJoinSemi{A,\JoinOp,o}\land
        \exists x\in A\,(x\neq o)\\[-2pt]
      &\qquad\to\FranklSemi{A,\JoinOp,o}\Bigr)
        \dsep\FranklFam M\vdash
        \FranklSemi{\mathcal L_M,\FamilyJoinOp{M},\varnothing}
    \end{aligned}
  }[FranklComplementSemilatticeInstance]
    \proofstepwidestar[1]{%
      \forall B\,\forall\star\,\forall e\,
        \bigl(\mathsf H(B,\star,e)\to
          \mathsf P(B,\star,e)\bigr)}{\rA}
    \proofstepwidestar[2]{\FranklFam M}{\rA}
    \proofstepwidestar[2]{\Finite L}{%
      \FormulaRefAuto{FranklComplementFamilyFinite}[thm]{2}}
    \proofstepwidestar[2]{\ZeroJoinSemi{L,\diamond,\varnothing}}{%
      \FormulaRefAuto{FranklComplementFiniteLattice}[thm]{2}}
    \proofstepwidestar[2]{\exists X\in L\,(X\neq\varnothing)}{%
      \FormulaRefAuto{FranklComplementNontrivial}[thm]{2}}
    \proofstepwidestar[1]{%
      \forall\star\,\forall e\,
        \bigl(\mathsf H(L,\star,e)\to
          \mathsf P(L,\star,e)\bigr)}{\rUE{1}}
    \proofstepwidestar[1]{%
      \forall e\,
        \bigl(\mathsf H(L,\diamond,e)\to
          \mathsf P(L,\diamond,e)\bigr)}{\rUE{6}}
    \proofstepwidestar[1]{%
      \mathsf H(L,\diamond,\varnothing)\to
        \mathsf P(L,\diamond,\varnothing)}{\rUE{7}}
    \proofstepwidestar[2]{\mathsf H(L,\diamond,\varnothing)}{%
      \rAI{3,\rAI{4,5}}}
    \proofstepwidestar[1,2]{\mathsf P(L,\diamond,\varnothing)}{%
      \rRE{8,9}}
  \closeproofpartwideindR

  \proofpartwideindR[Punktweise Form der Komplementordnung]{%
    \begin{aligned}[t]
      \FranklFam M\vdash
      \forall Q,S\in\mathcal L_M\,\bigl(&
        Q\InducedLeq S\\[-2pt]
        &\leftrightarrow Q\subseteq S\bigr)
    \end{aligned}
  }{%
    \FranklFam M\vdash
    \forall Q,S\in\mathcal L_M\,
      (Q\InducedLeq S
      \leftrightarrow Q\subseteq S)
  }[FranklComplementOrderPointwise]
    \proofstepwidestar[1]{\FranklFam M}{\rA}
    \proofstepwidestar[2]{Q\in\mathcal L_M}{\rA}
    \proofstepwidestar[3]{S\in\mathcal L_M}{\rA}
    \proofstepwidestar[1,2,3]{%
      Q\InducedLeq S
      \leftrightarrow Q\subseteq S}{%
      \FormulaRefAuto{FranklComplementOrderRecovered}[thm]{1,2,3}}
    \proofstepwidestar[1,2]{%
      \forall S\in\mathcal L_M\,
      (Q\InducedLeq S
      \leftrightarrow Q\subseteq S)}{\rUI{\rRI{3,4}}}
    \proofstepwidestar[1]{%
      \forall Q,S\in\mathcal L_M\,
      (Q\InducedLeq S
      \leftrightarrow Q\subseteq S)}{\rUI{\rRI{2,5}}}
  \closeproofpartwideindR

  \proofpartwideindR[Irreduzibler Hauptfilter]{%
    \begin{aligned}[t]
      &\FranklSemi{\mathcal L_M,\FamilyJoinOp{M},\varnothing}\vdash{}\\[-2pt]
      &\exists P\Bigl(
        \JoinIrr{\mathcal L_M,\FamilyJoinOp{M},\varnothing}{P}
        \land\AtMostHalf{H_M(P)}{\mathcal L_M}\Bigr)
    \end{aligned}
  }{%
    \begin{aligned}[t]
      &\FranklSemi{\mathcal L_M,\FamilyJoinOp{M},\varnothing}\vdash{}\\[-2pt]
      &\exists P\Bigl(
        \JoinIrr{\mathcal L_M,\FamilyJoinOp{M},\varnothing}{P}
        \land\AtMostHalf{H_M(P)}{\mathcal L_M}\Bigr)
    \end{aligned}
  }[FranklComplementIrreducibleFilter]
    \proofstepwidestar[1]{%
      \FranklSemi{\mathcal L_M,\FamilyJoinOp{M},\varnothing}}{\rA}
    \proofstepwidestar[1]{%
      \begin{aligned}[t]
        \exists P\Bigl(&
          \JoinIrr{\mathcal L_M,\FamilyJoinOp{M},\varnothing}{P}\land{}\\[-2pt]
        &\AtMostHalf{%
          \PrincipalFilter{\mathcal L_M,
            \InducedLeq}{P}}{\mathcal L_M}\Bigr)
      \end{aligned}}{%
      \FormulaRefAuto{FranklSemilatticeProperty}[def]{1}}
    \proofstepwidestar[3]{%
      \begin{aligned}[t]
        &\JoinIrr{\mathcal L_M,\FamilyJoinOp{M},\varnothing}{P}\land{}\\[-2pt]
        &\AtMostHalf{%
          \PrincipalFilter{\mathcal L_M,
            \InducedLeq}{P}}{\mathcal L_M}
      \end{aligned}}{\rA}
    \proofstepwidestar[3]{%
      \JoinIrr{\mathcal L_M,\FamilyJoinOp{M},\varnothing}{P}}{\rAEa{3}}
    \proofstepwidestar[3]{%
      \AtMostHalf{%
        \PrincipalFilter{\mathcal L_M,
          \InducedLeq}{P}}{\mathcal L_M}}{\rAEb{3}}
    \proofstepwidestar[]{%
      H_M(P)=\PrincipalFilter{%
        \mathcal L_M,\InducedLeq}{P}}{%
      \FormulaRefAuto{FranklComplementPrincipalFilterAbbreviation}[def]}
    \proofstepwidestar[]{%
      \PrincipalFilter{\mathcal L_M,
        \InducedLeq}{P}=H_M(P)}{%
      \FormulaRefAuto{a=b\vdash b=a}{6}}
    \proofstepwidestar[3]{\AtMostHalf{H_M(P)}{\mathcal L_M}}{\rIE{7,5}}
    \proofstepwidestar[3]{%
      \exists P\Bigl(
        \JoinIrr{\mathcal L_M,\FamilyJoinOp{M},\varnothing}{P}
        \land\AtMostHalf{H_M(P)}{\mathcal L_M}\Bigr)}{%
      \rEI{\rAI{4,8}}}
    \proofstepwidestar[1]{%
      \exists P\Bigl(
        \JoinIrr{\mathcal L_M,\FamilyJoinOp{M},\varnothing}{P}
        \land\AtMostHalf{H_M(P)}{\mathcal L_M}\Bigr)}{%
      \rEE{2,3,9}}
  \closeproofpartwideindR

  \proofpartwideindR[Privater Punkt des Hauptfilters]{%
    \begin{aligned}[t]
      &\FranklFam M\dsep
        \JoinIrr{\mathcal L_M,\FamilyJoinOp{M},\varnothing}{P}\vdash{}\\[-2pt]
      &\exists a\in U_M\,\forall Q\in\mathcal L_M\,
        (a\in Q\leftrightarrow Q\in H_M(P))
    \end{aligned}
  }{%
    \begin{aligned}[t]
      &\FranklFam M\dsep
        \JoinIrr{\mathcal L_M,\FamilyJoinOp{M},\varnothing}{P}\vdash{}\\[-2pt]
      &\exists a\in U_M\,\forall Q\in\mathcal L_M\,
        (a\in Q\leftrightarrow Q\in H_M(P))
    \end{aligned}
  }[FranklComplementPrivateFilterPoint]
    \proofstepwidestar[1]{\FranklFam M}{\rA}
    \proofstepwidestar[2]{%
      \JoinIrr{\mathcal L_M,\FamilyJoinOp{M},\varnothing}{P}}{\rA}
    \proofstepwidestar[1]{\Finite{\mathcal L_M}}{%
      \FormulaRefAuto{FranklComplementFamilyFinite}[thm]{1}}
    \proofstepwidestar[1]{%
      \ZeroJoinSemi{\mathcal L_M,\FamilyJoinOp{M},\varnothing}}{%
      \FormulaRefAuto{FranklComplementFiniteLattice}[thm]{1}}
    \proofstepwidestar[1]{\InterClosedFam{\mathcal L_M}}{%
      \FormulaRefAuto{FranklComplementFamilyIntersectionClosed}[thm]{1}}
    \proofstepwidestar[1]{%
      \forall Q,S\in\mathcal L_M\,
      (Q\InducedLeq S
      \leftrightarrow Q\subseteq S)}{%
      \FormulaRefAuto{FranklComplementOrderPointwise}[thm]{1}}
    \proofstepwidestar[1,2]{%
      \begin{aligned}[t]
        \exists a\in P\,\forall Q\in\mathcal L_M\,
        \bigl(a\in Q\leftrightarrow
          Q\in\PrincipalFilter{%
            \mathcal L_M,
            \InducedLeq}{P}\bigr)
      \end{aligned}}{%
      \FormulaRefAuto{JoinIrreduciblePrivatePoint}[thm]{3,4,5,6,2}}
    \proofstepwidestar[8]{%
      \begin{aligned}[t]
        a\in P\land\forall Q\in\mathcal L_M\,
        \bigl(a\in Q\leftrightarrow
          Q\in\PrincipalFilter{%
            \mathcal L_M,
            \InducedLeq}{P}\bigr)
      \end{aligned}}{\rA}
    \proofstepwidestar[8]{a\in P}{\rAEa{8}}
    \proofstepwidestar[8]{%
      \forall Q\in\mathcal L_M\,
      \bigl(a\in Q\leftrightarrow
        Q\in\PrincipalFilter{%
          \mathcal L_M,
          \InducedLeq}{P}\bigr)}{\rAEb{8}}
    \proofstepwidestar[2]{P\in\mathcal L_M}{%
      \FormulaRefAuto{JoinIrreducible}[def]{2}}
    \proofstepwidestar[1,2]{P\subseteq U_M}{%
      \FormulaRefAuto{FranklComplementMemberSubsetCarrier}[thm]{1,11}}
    \proofstepwidestar[1,2,8]{a\in U_M}{%
      \FormulaRefAuto{A\subseteq B,\,x\in A\vdash x\in B}{12,9}}
    \proofstepwidestar[]{%
      H_M(P)=\PrincipalFilter{%
        \mathcal L_M,\InducedLeq}{P}}{%
      \FormulaRefAuto{FranklComplementPrincipalFilterAbbreviation}[def]}
    \proofstepwidestar[]{%
      \PrincipalFilter{\mathcal L_M,
        \InducedLeq}{P}=H_M(P)}{%
      \FormulaRefAuto{a=b\vdash b=a}{14}}
    \proofstepwidestar[8]{%
      \forall Q\in\mathcal L_M\,(a\in Q\leftrightarrow Q\in H_M(P))}{%
      \rIE{15,10}}
    \proofstepwidestar[1,2,8]{%
      a\in U_M\land
      \forall Q\in\mathcal L_M\,(a\in Q\leftrightarrow Q\in H_M(P))}{%
      \rAI{13,16}}
    \proofstepwidestar[1,2]{%
      \exists a\in U_M\,\forall Q\in\mathcal L_M\,
        (a\in Q\leftrightarrow Q\in H_M(P))}{%
      \rEE{7,8,\rEI{17}}}
  \closeproofpartwideindR

  \proofpartwideindR[Transport des privaten Punkts]{%
    \begin{aligned}[t]
      &\FranklFam M\dsep\AtMostHalf{H_M(P)}{\mathcal L_M}\dsep
        a\in U_M\dsep{}\\[-2pt]
      &\forall Q\in\mathcal L_M\,(a\in Q\leftrightarrow Q\in H_M(P))
        \vdash\HalfOcc{a}{M}
    \end{aligned}
  }{%
    \begin{aligned}[t]
      &\FranklFam M\dsep\AtMostHalf{H_M(P)}{\mathcal L_M}\dsep
        a\in U_M\dsep{}\\[-2pt]
      &\forall Q\in\mathcal L_M\,(a\in Q\leftrightarrow Q\in H_M(P))
        \vdash\HalfOcc{a}{M}
    \end{aligned}
  }[FranklComplementPrivatePointHalfOccurrence]
    \proofstepwidestar[1]{\FranklFam M}{\rA}
    \proofstepwidestar[2]{\AtMostHalf{H_M(P)}{\mathcal L_M}}{\rA}
    \proofstepwidestar[3]{a\in U_M}{\rA}
    \proofstepwidestar[4]{%
      \forall Q\in\mathcal L_M\,(a\in Q\leftrightarrow Q\in H_M(P))}{\rA}
    \proofstepwidestar[]{%
      H_M(P)=\PrincipalFilter{%
        \mathcal L_M,\InducedLeq}{P}}{%
      \FormulaRefAuto{FranklComplementPrincipalFilterAbbreviation}[def]}
    \proofstepwidestar[]{%
      \PrincipalFilter{\mathcal L_M,
        \InducedLeq}{P}=H_M(P)}{%
      \FormulaRefAuto{a=b\vdash b=a}{5}}
    \proofstepwidestar[]{%
      \PrincipalFilter{\mathcal L_M,
        \InducedLeq}{P}\subseteq\mathcal L_M}{%
      \FormulaRefAuto{PrincipalFilterSubsetCarrier}[thm]}
    \proofstepwidestar[]{H_M(P)\subseteq\mathcal L_M}{\rIE{6,7}}
    \proofstepwidestar[1,2,3,4]{\HalfOcc{a}{M}}{%
      \FormulaRefAuto{FranklComplementPrivatePointTransport}[thm]{1,8,3,4,2}}
  \closeproofpartwideindR

  \proofpartwide{Schluss}
    \proofstepwidestar[1]{%
      \begin{aligned}[t]
        &\forall A\,\forall\JoinOp\,\forall o\,\Bigl(\\[-2pt]
        &\quad\Finite A\land\ZeroJoinSemi{A,\JoinOp,o}\land
          \exists x\in A\,(x\neq o)\\[-2pt]
        &\quad\to\FranklSemi{A,\JoinOp,o}\Bigr)
      \end{aligned}}{\rA}
    \proofstepwidestar[2]{\FranklFam M}{\rA}
    \proofstepwidestar[1,2]{%
      \FranklSemi{\mathcal L_M,\FamilyJoinOp{M},\varnothing}}{%
      \FormulaRefAuto{FranklComplementSemilatticeInstance}[thm]{1,2}}
    \proofstepwidestar[1,2]{%
      \begin{aligned}[t]
        &\exists P\Bigl(
          \JoinIrr{\mathcal L_M,\FamilyJoinOp{M},\varnothing}{P}\land{}\\[-2pt]
        &\qquad\AtMostHalf{H_M(P)}{\mathcal L_M}\Bigr)
      \end{aligned}}{%
      \FormulaRefAuto{FranklComplementIrreducibleFilter}[thm]{3}}
    \proofstepwidestar[5]{%
      \begin{aligned}[t]
        &\JoinIrr{\mathcal L_M,\FamilyJoinOp{M},\varnothing}{P}\land{}\\[-2pt]
        &\quad\AtMostHalf{H_M(P)}{\mathcal L_M}
      \end{aligned}}{\rA}
    \proofstepwidestar[5]{%
      \JoinIrr{\mathcal L_M,\FamilyJoinOp{M},\varnothing}{P}}{\rAEa{5}}
    \proofstepwidestar[5]{\AtMostHalf{H_M(P)}{\mathcal L_M}}{\rAEb{5}}
    \proofstepwidestar[2,5]{%
      \begin{aligned}[t]
        &\exists a\in U_M\,\forall Q\in\mathcal L_M\,\\[-2pt]
        &\quad(a\in Q\leftrightarrow Q\in H_M(P))
      \end{aligned}}{%
      \FormulaRefAuto{FranklComplementPrivateFilterPoint}[thm]{2,6}}
    \proofstepwidestar[9]{%
      \begin{aligned}[t]
        &a\in U_M\land{}\\[-2pt]
        &\quad\forall Q\in\mathcal L_M\,
          (a\in Q\leftrightarrow Q\in H_M(P))
      \end{aligned}}{\rA}
    \proofstepwidestar[9]{a\in U_M}{\rAEa{9}}
    \proofstepwidestar[9]{%
      \begin{aligned}[t]
        &\forall Q\in\mathcal L_M\,\\[-2pt]
        &\quad(a\in Q\leftrightarrow Q\in H_M(P))
      \end{aligned}}{\rAEb{9}}
    \proofstepwidestar[2,5,9]{\HalfOcc{a}{M}}{%
      \ProofRefStack{%
        \FormulaRefAuto{FranklComplementPrivatePointHalfOccurrence}[thm]\\
        (2,7,10,11)}}
    \proofstepwidestar[2,5,9]{\Frankl M}{%
      \FormulaRefAuto{Frankl-Eigenschaft}[def]{\rEI{12}}}
    \proofstepwidestar[1,2]{\Frankl M}{%
      \rEE{4,5,\rEE{8,9,13}}}
  \closeproofpartwide
\end{tabproofsplitwide}

\section{Globale Äquivalenz von Mengen- und Halbverbandsform}

\FormulaDefDeltaK[Globale Mengenform der Frankl-Vermutung]{%
  \FranklSetStatement\coloneqq
  \forall M\,(\FranklFam M\to\Frankl M)%
}{FranklGlobalSetStatement}{%
  \DeltaRow{Mengen}{M}
  \DeltaRow{\textbf{Neues Symbol}}{}
  \DeltaPrem{Globale Mengenform}{\FranklSetStatement}
}

\FormulaDefDeltaK[Globale Halbverbandsform der Frankl-Vermutung]{%
  \begin{aligned}[t]
    \FranklJoinStatement\coloneqq
    \forall A\,\forall\JoinOp\,\forall o\,\Bigl(&
      \Finite A\land\ZeroJoinSemi{A,\JoinOp,o}\land{}\\[-2pt]
      &\exists x\in A\,(x\neq o)
      \to\FranklSemi{A,\JoinOp,o}\Bigr)
  \end{aligned}
}{FranklGlobalJoinStatement}{%
  \DeltaRow{Mengen}{A\dsep x\dsep o}
  \DeltaRow{Funktionen}{\JoinOp}
  \DeltaRow{\textbf{Neues Symbol}}{}
  \DeltaPrem{Globale Halbverbandsform}{\FranklJoinStatement}
}

\paragraph{Beweisskizze.}
Die Vorwärtsrichtung wendet die Mengenform auf die kanonische Profilfamilie
eines endlichen Halbverbandes an. Die Rückrichtung wendet die
Halbverbandsform auf die komplementduale Konstruktion einer Frankl-Familie
an. Der folgende Satz fasst daher nur die beiden bereits bewiesenen
Transporte zusammen.

\FormulaThmDeltaK[Globale Äquivalenz der beiden Frankl-Formen]{%
  \FranklSetStatement\leftrightarrow\FranklJoinStatement
}{FranklSetSemilatticeEquivalence}{%
  \DeltaRow{Mengen}{A\dsep M\dsep x\dsep o}
  \DeltaRow{Funktionen}{\JoinOp}
}
Die Lemmon-Tabelle verwendet die beiden registrierten Kurzzeichen. Jede
Entfaltung wird ausdrücklich auf die zugehörige Definition gestützt.
\begin{tabproofsplitwide}
  \proofpartwideindR[Von der Mengenform zur Halbverbandsform]{%
    \FranklSetStatement\vdash\FranklJoinStatement
  }{%
    \FranklSetStatement\vdash\FranklJoinStatement
  }[FranklSetToJoinGlobal]
    \proofstepwidestar[1]{\FranklSetStatement}{\rA}
    \proofstepwidestar[1]{\forall M\,(\FranklFam M\to\Frankl M)}{%
      \FormulaRefAuto{FranklGlobalSetStatement}[def]{1}}
    \proofstepwidestar[3]{\Finite A\land
      \ZeroJoinSemi{A,\JoinOp,o}\land\exists x\in A\,(x\neq o)}{\rA}
    \proofstepwidestar[3]{\Finite A}{\rAEa{3}}
    \proofstepwidestar[3]{\ZeroJoinSemi{A,\JoinOp,o}}{\rAEa{\rAEb{3}}}
    \proofstepwidestar[3]{\exists x\in A\,(x\neq o)}{\rAEb{\rAEb{3}}}
    \proofstepwidestar[1,3]{\FranklSemi{A,\JoinOp,o}}{%
      \FormulaRefAuto{FranklSetImpliesSemilattice}[thm]{2,4,5,6}}
    \proofstepwidestar[1]{%
      \begin{aligned}[t]
        &\Finite A\land\ZeroJoinSemi{A,\JoinOp,o}\land
          \exists x\in A\,(x\neq o)\\[-2pt]
        &\qquad\to\FranklSemi{A,\JoinOp,o}
      \end{aligned}}{\rRI{3,7}}
    \proofstepwidestar[1]{%
      \begin{aligned}[t]
        \forall A\,\forall\JoinOp\,\forall o\,\Bigl(&
          \Finite A\land\ZeroJoinSemi{A,\JoinOp,o}\land{}\\[-2pt]
        &\exists x\in A\,(x\neq o)
          \to\FranklSemi{A,\JoinOp,o}\Bigr)
      \end{aligned}}{\rUI{\rUI{\rUI{8}}}}
    \proofstepwidestar[1]{\FranklJoinStatement}{%
      \FormulaRefAuto{FranklGlobalJoinStatement}[def]{9}}
  \closeproofpartwideindR

  \proofpartwideindR[Von der Halbverbandsform zur Mengenform]{%
    \FranklJoinStatement\vdash\FranklSetStatement
  }{%
    \FranklJoinStatement\vdash\FranklSetStatement
  }[FranklJoinToSetGlobal]
    \proofstepwidestar[1]{\FranklJoinStatement}{\rA}
    \proofstepwidestar[1]{%
      \begin{aligned}[t]
        \forall A\,\forall\JoinOp\,\forall o\,\Bigl(&
          \Finite A\land\ZeroJoinSemi{A,\JoinOp,o}\land{}\\[-2pt]
        &\exists x\in A\,(x\neq o)
          \to\FranklSemi{A,\JoinOp,o}\Bigr)
      \end{aligned}}{%
      \FormulaRefAuto{FranklGlobalJoinStatement}[def]{1}}
    \proofstepwidestar[3]{\FranklFam M}{\rA}
    \proofstepwidestar[1,3]{\Frankl M}{%
      \FormulaRefAuto{FranklSemilatticeImpliesSet}[thm]{3,2}}
    \proofstepwidestar[1]{\FranklFam M\to\Frankl M}{\rRI{3,4}}
    \proofstepwidestar[1]{\forall M\,(\FranklFam M\to\Frankl M)}{\rUI{5}}
    \proofstepwidestar[1]{\FranklSetStatement}{%
      \FormulaRefAuto{FranklGlobalSetStatement}[def]{6}}
  \closeproofpartwideindR

  \proofpartwide{Schluss}
    \proofstepwidestar[]{\FranklSetStatement\to\FranklJoinStatement}{%
      \FormulaRefAuto{FranklSetToJoinGlobal}[thm]}
    \proofstepwidestar[]{\FranklJoinStatement\to\FranklSetStatement}{%
      \FormulaRefAuto{FranklJoinToSetGlobal}[thm]}
    \proofstepwidestar[]{\FranklSetStatement\leftrightarrow
      \FranklJoinStatement}{%
      \FormulaRefAuto{P \leftrightarrow Q \eqvdash
        (P \rightarrow Q)\land(Q\rightarrow P)}{\rAI{1,2}}}
  \closeproofpartwide
\end{tabproofsplitwide}

Der Äquivalenzsatz beweist die Übersetzung, nicht die Frankl-Vermutung selbst:
\(\FranklSetStatement\Longleftrightarrow\FranklJoinStatement\).
Ein künftiger Beweis oder ein Gegenbeispiel lässt sich damit ohne Verlust
zwischen beiden Sprachen übertragen.

\section{Status und Literaturhinweise}

Die Frankl-Vermutung ist mit Stand vom 23.~Juli~2026 offen; siehe
\href{https://doi.org/10.23952/jano.8.2026.1.02}%
  {van der Hout--Roos, \emph{Some Results and Conjectures Related to
  Frankl's Union Closed Conjecture}, 2026}.
Äquivalente Fassungen behandelt
\href{https://doi.org/10.1016/0097-3165(92)90068-6}%
  {Poonen, \emph{Union-closed families}, 1992}.

Auch die allgemeine Gitterform der Frankl-Vermutung ist offen.
\href{https://arxiv.org/abs/2503.00277}%
  {Bouchard (2025)} formuliert diese Fassung und leitet notwendige Bedingungen
für einen kleinsten Gegenbeispielverband her; die Arbeit beweist die allgemeine
Gitterform nicht. Für untere semimodulare Gitter ist sie bewiesen; siehe
\href{https://doi.org/10.1007/s003730050008}%
  {Reinhold (2000)}.

\begin{remark}
Der allgemeine quantitative Ersatz ist schwächer als die Hälfte:
\href{https://doi.org/10.37236/12232}%
  {Alweiss--Huang--Sellke, \emph{Improved lower bound for the union-closed
  sets conjecture}, 2024} beweisen die Schranke \((3-\sqrt5)/2\).
Die Übersetzung ergibt daraus ein \(\JoinOp\)-irreduzibles \(j\) mit
\[
  \lvert\PrincipalFilter{A,\InducedLeq}{j}\rvert
  \leq \frac{\sqrt5-1}{2}\,\lvert A\rvert.
\]
Dieser externe Satz wird nur eingeordnet, nicht als intern bewiesenes Theorem
registriert.
\end{remark}


\end{document}
